\documentclass[11pt,lang=cn,math=cm,thmcnt=section,usesamecnt,scheme=chinese,oneside,bibstyle=alphabetic,citestyle=alphabetic]{elegantbook}
\usepackage{mathtools}
\usepackage{unicode-math}
\setmathfont{NewCMMath-Regular.otf}
\setoperatorfont\symup
\renewcommand{\dashrightarrow}{\mathrel{\rightdasharrow}}
\AtBeginDocument{\let\mathrm\symup\let\mathbf\symbf\let\mathit\symit}
\usepackage{multicol}
\usepackage{bookmark}
\usepackage{float}
\usepackage{needspace}
\usepackage[noautomatic]{imakeidx}
\makeindex[title=中文索引,intoc,columns=2]
\makeatletter
\renewenvironment{theindex}{\chapter*{中文索引}\markboth{中文索引}{中文索引}\phantomsection\addcontentsline{toc}{chapter}{中文索引}\begin{multicols}{2}\small\setlength{\parindent}{0pt}\setlength{\parskip}{0pt}\let\item\@idxitem\renewcommand{\indexspace}{\par\vspace{2pt}}}{\end{multicols}}
\makeatother
\renewcommand{\examplename}{例题}
\renewcommand{\theexam}{}% 原书6个例子均未编号，保留原生例题环境而不引入新号
\renewcommand{\exercisename}{练习}
\renewcommand{\solutionname}{解答}
\elegantnewtheorem{conjecture}{猜想}{thmstyle}{conj}
\elegantnewtheorem{claim}{断言}{thmstyle}{claim}
\newcommand{\numberednote}[1]{\refstepcounter{ELEGANT@samecnt}\label{note:#1}\textbf{\csname theELEGANT@samecnt\endcsname}\quad}
\addbibresource{references/references.bib}
\DeclareSortingTemplate{source}{\sort{\field{sortkey}}}
\ExecuteBibliographyOptions{sorting=source,maxbibnames=99,backref=true}
\setlength{\emergencystretch}{2em}
\hypersetup{unicode=true,bookmarksnumbered=true,bookmarksopen=true,pdftitle={几何 Langlands 理论导论},pdfauthor={Lin Chen; Wenhan Dai},pdfsubject={完整中文译本，第1–16章},colorlinks=true,linkcolor=structurecolor,citecolor=structurecolor,urlcolor=structurecolor}
\setcounter{tocdepth}{1}
\setcounter{secnumdepth}{1}
\title{几何 Langlands 理论导论}
\subtitle{Introduction to Geometric Langlands Theory}
\author{Lin Chen 讲授\quad Wenhan Dai 记录}
\date{2023 年秋季课程\quad 2025 年修订}
\version{完整中文译本}
\cover{images/cover.png}
\logo{images/logo-blue.png}
\graphicspath{{images/}}
\newcommand{\Bun}{\operatorname{Bun}}
\newcommand{\LocSys}{\operatorname{LocSys}}
\newcommand{\Shv}{\operatorname{Shv}}
\newcommand{\DMod}{\operatorname{DMod}}
\newcommand{\QCoh}{\operatorname{QCoh}}
\newcommand{\Coh}{\operatorname{Coh}}
\newcommand{\IndCoh}{\operatorname{IndCoh}}
\newcommand{\Pic}{\operatorname{Pic}}
\newcommand{\Hecke}{\operatorname{Hecke}}
\newcommand{\Whit}{\operatorname{Whit}}
\newcommand{\Gr}{\operatorname{Gr}}
\newcommand{\Fl}{\mathcal{F}\!\ell}
\newcommand{\Op}{\operatorname{Op}}
\newcommand{\Aut}{\operatorname{Aut}}
\newcommand{\Av}{\operatorname{Av}}
\newcommand{\Tor}{\operatorname{Tor}}
\newcommand{\Hom}{\operatorname{Hom}}
\newcommand{\Ext}{\operatorname{Ext}}
\newcommand{\End}{\operatorname{End}}
\newcommand{\Spec}{\operatorname{Spec}}
\newcommand{\Sym}{\operatorname{Sym}}
\newcommand{\Rep}{\operatorname{Rep}}
\newcommand{\Gal}{\operatorname{Gal}}
\newcommand{\Fr}{\operatorname{Fr}}
\newcommand{\Tr}{\operatorname{Tr}}
\newcommand{\supp}{\operatorname{supp}}
\newcommand{\GL}{\operatorname{GL}}
\newcommand{\SL}{\operatorname{SL}}
\newcommand{\rank}{\operatorname{rank}}
\newcommand{\length}{\operatorname{length}}
\newcommand{\CT}{\operatorname{CT}}
\newcommand{\RGamma}{\mathrm{R}\Gamma}
\newcommand{\ndeg}{\mathring{\deg}}
\newcommand{\cF}{\mathcal{F}}
\newcommand{\cL}{\mathcal{L}}
\newcommand{\cK}{\mathcal{K}}
\newcommand{\cO}{\mathcal{O}}
\newcommand{\cD}{\mathcal{D}}
\newcommand{\cN}{\mathcal{N}}
\newcommand{\cA}{\mathcal{A}}
\newcommand{\A}{\mathbb{A}}
\newcommand{\Gm}{\mathbb{G}_{m}}
\newcommand{\Ga}{\mathbb{G}_{a}}
\newcommand{\Qell}{\overline{\mathbb{Q}}_{\ell}}
\newcommand{\checkG}{\check{G}}
\newcommand{\chref}[1]{第\ref{ch:#1}章}
\newcommand{\secref}[1]{第\ref{sec:#1}节}
\newcommand{\thref}[1]{定理\ref{thm:#1}}
\newcommand{\lemref}[1]{引理\ref{lem:#1}}
\newcommand{\conjref}[1]{猜想\ref{conj:#1}}
\newcommand{\origref}[2]{\href[page=#1]{references/original.pdf}{#2}}
\newcommand{\srcpage}[1]{\par\label{src:#1}\ignorespaces}
\newlength{\origimageunit}
\setlength{\origimageunit}{0.7176470588\paperwidth}% 原PDF正文439.2bp／页面612bp
\newcommand{\origdiagram}[2]{\begin{figure}[H]\centering\includegraphics[width=#2\origimageunit]{#1}\end{figure}}
\newcommand{\origdiagramtag}[4]{\begin{figure}[H]\begin{equation}\vcenter{\hbox{\includegraphics[width=#2\origimageunit]{#1}}}\tag{#3}\label{#4}\end{equation}\end{figure}}
% ===== 译注（2026-11 校订新增）：每条第 1 参数是批注类型 =====
\newcounter{transnote}
\setcounter{transnote}{0}
\newenvironment{transnote}[1]{%
  \par\addvspace{0.45\baselineskip}%
  \refstepcounter{transnote}\label{tn:\arabic{transnote}}%
  \begin{list}{}{%
    \small
    \setlength{\leftmargin}{2.8em}\setlength{\labelwidth}{2.5em}%
    \setlength{\labelsep}{0.5em}\setlength{\topsep}{0pt}%
    \setlength{\parsep}{0.15em}\setlength{\itemsep}{0pt}%
    \setlength{\rightmargin}{0pt}}%
  \item[\textcolor{structurecolor}{\bfseries 译注 \arabic{transnote}}]%
  \textcolor{structurecolor}{\itshape 〈#1〉}\enspace\ignorespaces}%
{\end{list}\par\addvspace{0.15\baselineskip}}
\begin{document}
\maketitle
\hypersetup{pageanchor=true}
\frontmatter
\tableofcontents
\mainmatter
\chapter{导论}\label{ch:1}
\srcpage{5}
本章概述从整体类域论通向几何 Langlands 理论的整体、局部和量子形式的路径。关键在于，Galois 数据与自守函数之间的经典关系具有几何提升：将函数替换为模叠上的层，将 Frobenius 特征值替换为 Hecke 特征同构。现代的函子化表述见文献\cite{GR24a,AG15}。

本讲义的安排如下。\chref{2}单独介绍 $\Bun_G$ 及其基本的叠论描述和阿代尔描述；同一章还表述 Hecke 特征性质，并给出 $\GL_n$ 的微小 Hecke 论证。\chref{3}讨论 Whittaker 模型，\chref{4}构造并分析 Laumon 层。第\ref{ch:5}–\ref{ch:9}章，尤其是\chref{8}中的构造，证明主要的 Fourier 变换、洁净延拓和消失结论。范畴形式在\chref{10}、\chref{11}、\chref{12}和\chref{13}中展开；Hitchin 理论和局部方面的输入见\chref{14}和\chref{15}；整体证明则安排在\chref{16}。

\section{几何类域论}\label{sec:1.1}
\index{类域论!几何类域论|hyperpage}\index{阿代尔|hyperpage}\index{非分歧|hyperpage}
设 $K$ 为整体域，$\A_K$ 为其阿代尔环，$\Gal_{\overline K/K}$ 为其绝对 Galois 群。整体类域论给出一个“几乎”同构；在本概述中，这个词有意使用得不甚精确，后续讲授将以严格陈述取代这一概括：
\[
\theta_K:K^\times\backslash\A_K^\times\longrightarrow
\Gal_{\overline K/K}^{\mathrm{ab}}\simeq\Gal_{K^{\mathrm{ab}}/K},
\]
其中 $K^{\mathrm{ab}}\subset\overline K$ 是最大阿贝尔扩张。这个同构由局部 Frobenius 元刻画：对每个有限素位 $v$，
\[\theta_K(\pi_v)=\Fr_v.\]
这里 $\pi_v$ 是 $K_v$ 的一个素元所确定的类，模 $\cO_v^\times$ 后有良好定义；$\Fr_v$ 是 $\Gal_{\overline{\kappa_v}/\kappa_v}$ 中几何 Frobenius 的一个提升，模惯性群 $I_v$ 后有良好定义。

\begin{transnote}{含糊点拨}
「几乎」在此是数学限定而非修辞。互反映射 $\theta_K$ 对数域只是单射、像稠密而不满；取 $\pi_0(K^\times\backslash\A_K^\times)$ 后核恰为单位分支，此时才得到同构（Artin 互反律的连通核即此分支）。对 $\mathbb F_q$ 上的函数域——本节随后转入的情形——精确陈述把 $\Gal^{\mathrm{ab}}$ 换成 Weil 群 $W_K^{\mathrm{ab}}$，$\theta_K$ 才是同构。下一段关于 $\mathbb Q_\ell^\times$ 完全不连通、因而无法检测连通核的说明，正是这一点的直接后果。
[来源：nLab "class field theory"，https://ncatlab.org/nlab/show/class+field+theory（Last revised 2022-10-18；2026-11-17 访问）；并见原书 p.5 第 2–3 段及 p.5 对 Weil 群的说明。]
\end{transnote}

\begin{theorem}{非分歧类域论：特征标形式}{1.1.1}
设 $\ell\ne\operatorname{char}(K)$。存在一个几乎双射
\[
\{\text{非分歧特征标 }\xi:K^\times\backslash\A_K^\times\to\Qell^\times\}
\longleftrightarrow
\{\text{非分歧特征标 }\rho:\Gal_{\overline K/K}\to\Qell^\times\},
\]
它由等式
\[\xi(\pi_v)=\rho(\Fr_v)\]
刻画。非分歧是指：对所有有限素位 $v$，$\xi$ 在 $\cO_v^\times$ 上平凡，$\rho$ 在 $I_v$ 上平凡。
\end{theorem}

对数域而言，由于 $\Qell^\times$ 完全不连通，互反映射的连通核无法被取值于 $\Qell^\times$ 的特征标检测到。对 $\mathbb F_q$ 上的函数域而言，精确表述使用 Weil 群 $W_{\overline K/K}$。

现在假设 $K$ 是 $k=\mathbb F_q$ 上一条光滑完备曲线 $X$ 的函数域。记 $\cO_K^\times:=\prod_v\cO_v^\times$。商 $K^\times\backslash\A_K^\times/\cO_K^\times$ 是离散的，Weil 定理给出
\[
K^\times\backslash\A_K^\times/\cO_K^\times\simeq\Pic(X),
\qquad\pi_v\longmapsto\cO_X(v).
\]
因此，非分歧自守特征标与 Picard 群的特征标相同：
\[
\{\text{非分歧 }\xi:K^\times\backslash\A_K^\times\to\Qell^\times\}
\simeq\{\xi:\Pic(X)\to\Qell^\times\}.
\]

\Needspace{4\baselineskip}
在 Galois 一侧，对于 $\Spec K$ 的一个几何点 $\eta$，有 $\Gal_{\overline K/K}\simeq\pi_1(\Spec K,\eta)$。非分歧表示经由以下满射分解：
\[\pi_1(\Spec K,\eta)\longrightarrow\pi_1(X,\eta)\]
\srcpage{6}
等价地，
\[
\begin{aligned}
\{\text{非分歧 }\rho:\Gal_{\overline K/K}\to\Qell^\times\}
&\simeq\{\rho:\pi_1(X,\eta)\to\Qell^\times\}\\
&\simeq\{X\text{ 上秩为一的 }\ell\text{-进局部系统 }\sigma\}.
\end{aligned}
\]
于是，类域论成为几乎双射
\[
\{\xi:\Pic(X)\to\Qell^\times\}
\longleftrightarrow\{X\text{ 上秩为一的 }\ell\text{-进局部系统 }\sigma\}.
\]

剩下的任务是将等式 $\xi(\pi_v)=\rho(\Fr_v)$ 用几何方式表达出来。设 $Y$ 是 $k=\mathbb F_q$ 上的代数簇，$\cF$ 是 $Y$ 上的 $\ell$-进层。对 $y\in Y(k)$，茎 $\cF_{\bar y}$ 带有 $\Gal_{\bar k/k}$ 的作用；定义其迹函数为
\[f_{\cF}(y):=\Tr(\Fr_y,\cF_{\bar y}).\]
这就是 Grothendieck 的层–函数对应。若 $\rho$ 对应局部系统 $\sigma$，则在 $k$-点 $v\in X(k)$ 处，
\[\rho(\Fr_v)=f_\sigma(v)=\Tr(\Fr_v,\sigma_{\bar v}).\]
对闭点也使用相同的记号，Frobenius 则取相应的几何 Frobenius。

\begin{corollary}{非分歧特征标关系的几何形式}{1.1.2}
存在几乎双射
\[
\{\xi:\Pic(X)\to\Qell^\times\}\longleftrightarrow
\{X\text{ 上秩为一的 }\ell\text{-进局部系统 }\sigma\},
\]
由 $\xi(\cO_X(v))=f_\sigma(v)$ 刻画。
\end{corollary}

现在将 Picard 群替换为相对 Picard 概形。对每个仿射测试概形 $S/k$，
\[
\Pic_{X/k}(S):=\Pic(X\times_k S)/\Pic(S),\qquad
\Pic_{X/k}(k)=\Pic(X).
\]
我们希望在 $\Pic_{X/k}$ 上得到一个几何对象 $\Aut_\sigma$，其迹函数就是与 $\sigma$ 相联系的特征标：
\origdiagram{p006-diagram.pdf}{0.8971}
对 $\cF$ 应施加的正确几何条件，是 Hecke 特征性质。

\begin{definition}{$\GL_1$ 的 Hecke 特征层}{1.1.3}
设 $\cF$ 是 $\Pic_{X/k}$ 上的层，$\sigma$ 是 $X$ 上秩为一的局部系统。若 $\cF$ 配备同构
\[\operatorname{add}^*\cF\simeq\sigma\boxtimes\cF,\]
且该同构满足连续两次修改的通常交换性条件，就称 $\cF$ 是以 $\sigma$ 为特征值的 Hecke 特征层。这里
\[
\operatorname{add}:X\times_k\Pic_{X/k}\longrightarrow\Pic_{X/k},
\qquad(v,L)\longmapsto L(v).
\]
\end{definition}

\begin{transnote}{含糊点拨}
「连续两次修改的通常交换性条件」指把 $\sigma$-等变性做成对有序点组一致的相容数据。对 $v\ne w$，先改 $v$ 再改 $w$ 与先改 $w$ 再改 $v$ 在 $\Pic_{X/k}$ 上落到同一点（$L(v)(w)=L(w)(v)$）；条件要求由此得到的两个同构 $\operatorname{add}^*_{v,w}\cF\simeq\sigma_v\otimes\sigma_w\otimes\cF$ 相等。缺了这条一致性，$\cF$ 只在每条 $\operatorname{add}$ 的纤维上「是特征的」，拼不成 $X\times_k\Pic_{X/k}$ 上带对称群相容性的层，\lemref{1.1.4} 也就无从谈起。一般 $G$ 的完整形式见 \S\ref{sec:2.6}：Hecke 函子给出的同构须与 $\Rep(\checkG)$ 中的张量积相容。
[依据：原书 p.6 定义 1.1.3 的措辞与 \S\ref{sec:2.6} 的一般定义；余圈条件本身出自 [BD91] 对 Hecke 特征层的原始构造。]
\end{transnote}

\Needspace{8\baselineskip}
\begin{lemma}{}{1.1.4}
设 $\cF$ 是以 $\sigma$ 为特征值的 Hecke 特征层，并假设 $\cF$ 在单位点 $\cO_X\in\Pic_{X/k}$ 处的限制是常值层，从而相应函数将单位元映为 $1$。在 $\Pic(X)=\Pic_{X/k}(k)$ 上令 $\xi=f_{\cF}$。则 $\xi$ 是特征标，且对每个素位 $v$，
\[f_{\cF}(\cO_X(v))=f_\sigma(v).\]
\end{lemma}
\srcpage{7}
\lemref{1.1.4}说明，Hecke 特征同构是在层的层面上对类域论中乘法性方程和 Frobenius 特征值方程的替代。

\begin{theorem}{Deligne}{1.1.5}
对 $X$ 上每个秩为一的局部系统 $\sigma$，在 $\Pic_{X/k}$ 上都存在本质唯一的 Hecke 特征层 $\Aut_\sigma$，其特征值为 $\sigma$；而且 $\Aut_\sigma$ 是局部系统。这里“本质唯一”是指，忽略与 $\Shv(\Spec k)$ 中对象的张量积后唯一；这种张量积保持特征值不变。
\end{theorem}

\thref{1.1.5}是纯几何陈述，Deligne 的证明也是几何的。因此，它就是非分歧整体几何类域论，等价地说，是 $\GL_1$ 的非分歧整体几何 Langlands 对应。

\section{非分歧整体几何 Langlands 对应}\label{sec:1.2}
\index{Langlands 对应!整体几何|hyperpage}\index{Hecke 特征层|hyperpage}\index{局部系统|hyperpage}
现在从 $\GL_1$ 转到 $\GL_n$，再转到 $k$ 上的分裂约化群 $G$。对 $X$ 上非常值约化群概形的版本也很重要，但不属于这组入门讲授的内容。

从现在起，$X$ 是完美域 $k$ 上光滑、完备且几何连通的曲线，$K=K(X)$ 是它的函数域。固定一种合理的层理论，其系数域 $e$ 代数闭。对性质良好的 $k$-概形或叠 $Y$，以 $\Shv_c(Y)^\heartsuit$ 表示阿贝尔范畴，以 $\Shv_c(Y)$ 表示导出范畴。标准选择为：
\begin{itemize}
\item Betti 情形：$k=\mathbb C$，使用复拓扑下的反常层；
\item $\ell$-进情形：$k$ 为有限域，使用反常 $\ell$-进 Weil 层；
\item de Rham 情形：$k$ 为特征 $0$ 的代数闭域，使用凝聚右 $D$ 模。
\end{itemize}

\begin{transnote}{易错预警}
上标 $\heartsuit$ 是 $t$-结构理论中的标准记号，指相应 $t$-结构的\textbf{心}（heart），不是「阿贝尔群值的层」。此处的「阿贝尔范畴」指所选层理论本身给出的对象范畴（Betti 情形即反常层范畴），而 $\Shv_c(Y)$ 是它的导出范畴。于是 $\cF\in\Shv_c(Y)^\heartsuit$ 与把 $\cF$ 当作导出范畴中的复形不可混同：书中的 Hecke 函子、Fourier 变换与上同调都在 $\Shv_c(Y)$ 层面计算，只在 $t$-正合性成立时才能把陈述降到心层面。\conjref{1.3.1} 与 \thref{1.3.2} 的差别正是这一处：前者是心范畴之间的 $t$-正合等价，后者放弃 $t$-正合性。
[来源：nLab "t-structure"，https://ncatlab.org/nlab/show/t-structure（Last revised 2023-04-20；2026-11-17 访问）；nLab "perverse sheaf"，https://ncatlab.org/nlab/show/perverse+sheaf（Last revised 2026-10-16；2026-11-17 访问）；并见原书 p.7 \S\ref{sec:1.2} 开头的约定。]
\end{transnote}

Galois 一侧的推广是形式上的：
\[
\begin{array}{c|l}
\GL_n&X\text{ 上秩为 }n\text{ 的局部系统 }\sigma,\\\hline
G&X\text{ 上的 }G\text{-局部系统 }\sigma.
\end{array}
\]
对于与 $G$ 相联系的自守层，实际特征值是系数域 $e$ 上的 Langlands 对偶群 $\checkG$ 的局部系统。

在自守一侧，$\GL_1=\Gm$ 的阿贝尔情形必须替换为其 $G$-类比，即
\[
\begin{array}{c|c|c|c}
\GL_1&K^\times\backslash\A_K^\times/\cO_K^\times&\Pic(X)&\Pic_{X/k}\\\hline
G&G(K)\backslash G(\A_K)/G(\cO_K)&?&?
\end{array}
\]
\begin{theorem}{Weil 一致化}{1.2.1}
在关于 $G$ 与 $X$ 的标准假设下，双陪集集合 $G(K)\backslash G(\A_K)/G(\cO_K)$ 自然等同于 $X$ 上主 $G$ 丛的同构类集合，也就是 $G$ 挠子的同构类集合。这些假设将在后面精确说明；这里所考虑的是曲线上的连通约化群。
\end{theorem}

与 $\Pic_{X/k}$ 不同，$G$ 丛的模对象通常不是概形。稳定子群随丛而变化，因此正确的对象是代数叠：它记住各点的自同构群。这些附加结构在单纯的双陪集集合中不可见，却对一般的 $G$ 至关重要。

\begin{proposition}{$\Bun_G$ 的定义}{1.2.2}
存在 $k$ 上的光滑 Artin 叠 $\Bun_G$，分类 $X$ 上的 $G$ 挠子族。在 $k$-点的同构类层面，
\[\Bun_G(k)\simeq G(K)\backslash G(\A_K)/G(\cO_K).\]
更精确的记号是 $\Bun_{G,X/k}$。
\end{proposition}
\srcpage{8}

Hecke 特征性质也需要推广。对 $G=\GL_1$，映射
\[\operatorname{add}:X\times_k\Pic_{X/k}\longrightarrow\Pic_{X/k},\qquad(x,L)\longmapsto L(x)\]
以典范方式在 $x$ 处修改线丛。对一般的 $G$ 挠子，在 $x$ 处并不存在唯一的典范修改，而是有许多种。因此，将映射 $\operatorname{add}$ 替换为整体 Hecke 对应
\origdiagram{p008-diagram.pdf}{0.5282}
它记录一个修改点，以及两个在该点之外彼此等同的 $G$ 丛。

假定已有 Hecke 特征层的一般定义，非分歧整体几何 Langlands 对应可以表述如下。
\begin{conjecture}{非分歧整体几何 Langlands}{1.2.3}
对 $X$ 上每个几何不可约的 $\checkG$-局部系统 $\sigma$，在 $\Bun_G$ 上都存在本质唯一的 Hecke 特征层 $\Aut_\sigma$，其特征值为 $\sigma$。
\end{conjecture}

这里，几何不可约是指：基变换到 $\bar k$ 后，相应 $\checkG$-局部系统不存在到 $\checkG$ 的真抛物子群的平坦约化；对于 $G=\GL_n$，这就是秩 $n$ 局部系统通常意义下的不可约性。如果 $\sigma$ 可约，预期得到的应是 $\Shv_c(\Bun_G)$ 中的 Hecke 特征对象，而不是心中的特征层；本质唯一性也不再成立。

对 $G=\GL_n$，存在性定理已经成立：Deligne 证明了 $n=1$ 的情形，Drinfeld 证明了 $n=2$ 的情形，Frenkel–Gaitsgory–Vilonen 证明了一般情形\cite{FGKV98}。在 $\ell$-进情形，对 $\Aut_\sigma$ 取 Frobenius 迹得到自守函数。因此，这个猜想给出函数域上从 Galois 数据到自守对象的几何构造；由 $\Aut_\sigma$ 得到的函数是尖点函数。

\section{非分歧整体几何 Langlands 等价}\label{sec:1.3}
\index{Langlands 等价!整体几何|hyperpage}\index{尖点|hyperpage}\index{奇异支撑!幂零|hyperpage}
\conjref{1.2.3}将一个自守层赋给一个不可约谱点。范畴形式则要求从整个谱模对象出发构造函子。以 $\LocSys_{\checkG}^{\mathrm{irred}}$ 表示分类 $X$ 上几何不可约 $\checkG$-局部系统族的代数几何对象。预期存在函子
\[\Coh(\LocSys_{\checkG}^{\mathrm{irred}})^\heartsuit\longrightarrow\Shv_c(\Bun_G)^\heartsuit,\qquad\delta_\sigma\longmapsto\Aut_\sigma,\]
而从理念上说，预期它与 $\Shv_c(\Bun_G)^\heartsuit$ 的尖点部分等价。

在 de Rham 情形，$\LocSys_{\checkG}^{\mathrm{irred}}$ 是光滑 Artin 叠，心层面的预期成为以下陈述。
\begin{conjecture}{尖点 de Rham 形式}{1.3.1}
存在典范的 $t$-正合等价
\[\mathbb L_G:\DMod_{\mathrm{coh}}(\Bun_G)_{\mathrm{cusp}}\simeq\Coh(\LocSys_{\checkG}^{\mathrm{irred}}),\]
使 $\mathbb L_G^{-1}$ 将不可约 $\checkG$-局部系统 $\sigma$ 处的摩天层映为以 $\sigma$ 为特征值的 Hecke 特征层 $\Aut_\sigma$。
\end{conjecture}

若将可约局部系统也包括进来，就要放弃 $t$-正合性，并在范畴层面工作。
\srcpage{9}
\begin{theorem}{朴素概括：完整 de Rham 形式}{1.3.2}
作为第一步近似，可以写出从谱一侧到自守一侧的等价
\[\mathbb L_G:\DMod_{\mathrm{coh}}(\Bun_G)\simeq\Coh(\LocSys_{\checkG}),\]
其中 $\checkG$-局部系统处的摩天层映为 Hecke 特征对象。这个式子只是一种指引性的概括：精确的非尖点 de Rham 表述，要将凝聚谱一侧替换为具有幂零奇异支撑的重整化范畴，例如 $\IndCoh_{\mathrm{Nilp}}(\LocSys_{\checkG})$。
\end{theorem}

这比经典 Langlands 对应高一个范畴层次：经典对象是函数与表示，几何对象则是层与 DG 范畴。更精确的现代表述将简化的凝聚谱一侧替换为适当的重整化范畴或具有幂零奇异支撑的范畴，例如 $\IndCoh_{\mathrm{Nilp}}(\LocSys_{\checkG})$\cite{AG15,GR24a}；但基本概括仍然是：$\Bun_G$ 上的自守范畴由 $\checkG$-局部系统上的谱范畴控制。这组讲授不讨论分歧整体理论。

\begin{transnote}{背景直觉}
「幂零奇异支撑」是 [AG15] 对谱一侧所作修正的核心。对 $\cF\in\QCoh(\LocSys_{\checkG})$，其奇异支撑 $\mathrm{SS}(\cF)\subset T^*\LocSys_{\checkG}$ 是余切丛里的闭锥；要求 $\mathrm{SS}(\cF)$ 落在幂零锥方向（即 $\LocSys^{\mathrm{nilp}}_{\checkG}$ 所指的那个锥）的对象构成全子范畴 $\IndCoh_{\mathrm{Nilp}}(\LocSys_{\checkG})$。它比 $\QCoh$ 小、比 $\Coh$ 大，并且在 Hecke 函子下稳定，因此才是与 $\DMod_{\mathrm{coh}}(\Bun_G)$ 相配的一侧。只保留此条件后，「几乎」所指的技术性缺口被填上，\S\ref{sec:1.3} 末尾的朴素整体形式才算精确。
[来源：nLab "geometric Langlands correspondence"，https://ncatlab.org/nlab/show/geometric+Langlands+correspondence（Last revised 2026-04-11；2026-11-17 访问）记 Arinkin–Gaitsgory 以 nilpotent ind-objects 替换 $\QCoh$ 并升级到稳定 $\infty$-范畴；并见 [AG15]（Selecta Math. 21 (2015), 1–199）的奇异支撑理论。]
\end{transnote}


\section{局部几何 Langlands 等价}\label{sec:1.4}
\index{Langlands 等价!局部几何|hyperpage}\index{环路群|hyperpage}\index{Whittaker 模型|hyperpage}
现在专注于 de Rham 情形。经典局部 Langlands 理论分类局部域 $K_v$ 上 $G(K_v)$ 的表示，它们构成一个 $1$-范畴。局部几何 Langlands 理论是其范畴化：它试图分类 $G(K_v)$ 的范畴表示，这些表示构成一个 $2$-范畴。

$G(K_v)$ 的几何体现是环路群 $LG$。在 $\mathbb C$ 上，它的 $\mathbb C$-点正是 $G(K_v)$ 的元素。定义 $LG\text{-}\mathrm{Mod}$ 为被幺半 DG 范畴 $\DMod(LG)$ 作用的 DG 范畴所构成的 $(\infty,2)$-范畴，其中幺半结构由卷积给出。若 $Y$ 是带有 $LG$ 作用的合理代数几何对象，则 $\DMod(Y)$ 是 $LG\text{-}\mathrm{Mod}$ 中的对象。

基本例子为 $\Gr_G$、$\Fl_G$ 和 $\Bun_G^{\infty;v}$：分别是仿射 Grassmann 簇、仿射旗 ind-簇，以及具有无限级结构的 $X$ 上 $G$ 挠子模 ind-叠，后者配备在 $v$ 的形式邻域 $D_v$ 上的平凡化。就复点而言，最后一个对象有双陪集描述
\[\prod_{w\ne v}G(\cO_w)\backslash G(\A_K)/G(K).\]

范畴 $\DMod(LG)$ 是它自身上的双模。对右作用取不变量，得到仍带有左 $LG$ 作用的 DG 范畴。于是，$\DMod(\Gr_G)$ 和 $\DMod(\Fl_G)$ 分别通过对右作用取 $L^+G$ 不变量和 $I$ 不变量得到。这里 $L^+G$ 是弧群，是 $G(\cO_v)$ 的几何对应物；$I$ 是 Iwahori 子群，是 $G(\cO_v)\times_{G(\kappa_v)}B(\kappa_v)$ 的几何对应物。

另一个重要对象是 Whittaker 模型
\[\Whit(LG)\in LG\text{-}\mathrm{Mod},\]
它通过相对于一般位置特征标 $\chi$ 取右 $LN$ 不变量而得到，其中 $N$ 是 Borel 子群 $B$ 的幂幺根基。还有一个关键对象 $\mathrm{KM}_{\mathfrak g}$，即临界水平下仿射 Lie 代数 $\widehat{\mathfrak g}$ 的表示所构成的 DG 范畴。

在 Galois 一侧，考虑穿孔圆盘 $D^\circ:=\Spec K_v$ 上 $\checkG$-局部系统的模问题 $\LocSys_{\checkG}(D^\circ)$。定义
\[\operatorname{QCohCat}(\LocSys_{\checkG}(D^\circ))\]
\srcpage{10}
为幺半范畴 $\QCoh(\LocSys_{\checkG}(D^\circ))$ 的模范畴所构成的 $(\infty,2)$-范畴，其中幺半结构为张量积。若 $Z\to\LocSys_{\checkG}(D^\circ)$ 是合理的代数几何对象，则 $\QCoh(Z)$ 就是这样的模范畴。

基本的谱一侧例子是
\[\LocSys_{\checkG}(D),\qquad\LocSys_{\checkG}^{\mathrm{nilp}}(D^\circ),\qquad\LocSys_{\checkG}(X^\circ),\]
分别分类圆盘 $D:=\Spec\cO_v$ 上的 $\checkG$-局部系统、穿孔圆盘上具有幂零奇点的 $\checkG$-局部系统，以及穿孔曲线 $X^\circ:=X\setminus\{v\}$ 上的 $\checkG$-局部系统。对约化群 $H$，有 $\LocSys_H(D)\simeq\mathrm{pt}/H$ 和 $\LocSys_H^{\mathrm{nilp}}(D^\circ)\simeq\cN/H$，其中 $\cN\subset\mathfrak h$ 是幂零锥。此外还考虑 $\LocSys_{\checkG}(D^\circ)$ 本身以及 oper 空间
\[\Op_{\checkG}(D^\circ).\]

局部猜想将上述自守一侧和谱一侧的例子配对；这就是局部几何 Langlands 的标准纲要形式\cite{FG05,FG07}。

\begin{conjecture}{局部几何 Langlands：纲要形式}{1.4.1}
存在几乎等价
\[LG\text{-}\mathrm{Mod}\simeq\operatorname{QCohCat}(\LocSys_{\checkG}(D^\circ)),\]
其中“几乎”掩盖了 $\QCoh$ 与凝聚或重整化变体之间的技术性不匹配：精确表述必须在两侧选取正确的凝聚部分。该等价具有以下对应：
\[
\begin{aligned}
\DMod(\Gr_G)&\longmapsto\QCoh(\LocSys_{\checkG}(D)),\\
\DMod(\Fl_G)&\longmapsto\QCoh(\LocSys_{\checkG}^{\mathrm{nilp}}(D^\circ)),\\
\DMod(\Bun_G^{\infty;v})&\longmapsto\QCoh(\LocSys_{\checkG}(X^\circ)),\\
\Whit(LG)&\longmapsto\QCoh(\LocSys_{\checkG}(D^\circ)),\\
\mathrm{KM}_{\mathfrak g}&\longmapsto\QCoh(\Op_{\checkG}(D^\circ)).
\end{aligned}
\]
\end{conjecture}

在列出的自守对象之间取 $\Hom$ 范畴，再与它们在谱一侧的像之间的 $\Hom$ 范畴比较，可以恢复许多熟悉的等价。在这些比较中，须将 $\DMod$ 替换为凝聚 $D$ 模，将 $\QCoh$ 替换为相应的凝聚谱范畴。
\begin{enumerate}
\item 根据 Bezrukavnikov–Finkelberg 的结果，导出几何 Satake 等价为
\[\End(\DMod(\Gr_G))\simeq\End(\QCoh(\LocSys_{\checkG}(D))).\]
\item Arkhipov–Bezrukavnikov–Ginzburg 定理为
\[
\begin{aligned}
&\Hom(\DMod(\Gr_G),\DMod(\Fl_G))\\
&\qquad\simeq\Hom(\QCoh(\LocSys_{\checkG}(D)),\QCoh(\LocSys_{\checkG}^{\mathrm{nilp}}(D^\circ))).
\end{aligned}
\]
\item 根据 Bezrukavnikov 关于仿射 Hecke 代数的两种几何实现之间等价的结果，
\[\End(\DMod(\Fl_G))\simeq\End(\QCoh(\LocSys_{\checkG}^{\mathrm{nilp}}(D^\circ))).\]
\item 有
\[
\begin{aligned}
&\Hom(\DMod(\Gr_G),\DMod(\Bun_G^{\infty;v}))\\
&\qquad\simeq\Hom(\QCoh(\LocSys_{\checkG}(D)),\QCoh(\LocSys_{\checkG}(X^\circ))).
\end{aligned}
\]
\srcpage{11}
这恢复了\thref{1.3.2}中的朴素整体形式；精确版本使用具有幂零奇异支撑的谱范畴。若将 $\DMod(\Gr_G)$ 和 $\QCoh(\LocSys_{\checkG}(D))$ 分别替换为 $\DMod(\Fl_G)$ 和 $\QCoh(\LocSys_{\checkG}^{\mathrm{nilp}}(D^\circ))$，就得到相应的驯顺分歧形式。
\item 根据 Frenkel–Gaitsgory–Vilonen 的结果\cite{FGKV98}，几何 Casselman–Shalika 公式为
\[
\begin{aligned}
&\Hom(\DMod(\Gr_G),\Whit(LG))\\
&\qquad\simeq\Hom(\QCoh(\LocSys_{\checkG}(D)),\QCoh(\LocSys_{\checkG}(D^\circ))).
\end{aligned}
\]
\item 根据 Arkhipov–Bezrukavnikov 的一个定理，
\[
\begin{aligned}
&\Hom(\DMod(\Fl_G),\Whit(LG))\\
&\qquad\simeq\Hom(\QCoh(\LocSys_{\checkG}^{\mathrm{nilp}}(D^\circ)),\QCoh(\LocSys_{\checkG}(D^\circ))).
\end{aligned}
\]
\item 以下等价称为 Feigin–Frenkel 定理\cite{Fre02}：
\[\Hom(\DMod(\Gr_G),\mathrm{KM}_{\mathfrak g})\simeq\Hom(\QCoh(\LocSys_{\checkG}(D)),\QCoh(\Op_{\checkG}(D^\circ))).\]
\item 根据 Raskin 的一个定理，
\[\Hom(\Whit(LG),\mathrm{KM}_{\mathfrak g})\simeq\Hom(\QCoh(\LocSys_{\checkG}(D^\circ)),\QCoh(\Op_{\checkG}(D^\circ))).\]
\end{enumerate}

\section{量子局部几何 Langlands 等价}\label{sec:1.5}
\index{Langlands 等价!量子局部几何|hyperpage}
扭曲 $D$ 模给出\conjref{1.4.1}中局部猜想的变形。预期的量子局部等价关于 $G$ 及其 Langlands 对偶群是对称的。
\begin{conjecture}{量子局部几何 Langlands}{1.5.1}
存在几乎等价
\[LG\text{-}\mathrm{Mod}_\kappa\simeq L\checkG\text{-}\mathrm{Mod}_{\check\kappa},\]
具有以下对应：
\[
\begin{aligned}
\DMod_\kappa(\Gr_G)&\longmapsto\DMod_{\check\kappa}(\Gr_{\checkG}),\\
\DMod_\kappa(\Fl_G)&\longmapsto\DMod_{\check\kappa}(\Fl_{\checkG}),\\
\DMod_\kappa(\Bun_G^{\infty;v})&\longmapsto\DMod_{\check\kappa}(\Bun_{\checkG}^{\infty;v}),\\
\Whit_\kappa(LG)&\longmapsto\mathrm{KM}_{\check{\mathfrak g},\check\kappa},\\
\mathrm{KM}_{\mathfrak g,\kappa}&\longmapsto\Whit_{\check\kappa}(L\checkG).
\end{aligned}
\]
\end{conjecture}

这个表格的对称性，就是 Langlands 对偶的量子版本。

\chapter{Hecke 特征性质}\label{ch:2}
% 原PDF第11页下半部
几何 Langlands 理论的自守一侧围绕 $\Bun_G$ 展开；这是 $X$ 上 $G$ 丛的模叠。有了 $\Bun_G$，就可以借助在曲线上某一点修改 $G$ 丛的对应来表述 Hecke 特征性质。

以下首先回顾定义 $\Bun_G$ 所用的预叠与叠的语言，记录其基本有限性性质，并将它的 $k$-点与通常的阿代尔双商联系起来。
\section{预叠与叠}\label{sec:2.1}
\srcpage{12}\index{预叠|hyperpage}\index{叠|hyperpage}\index{下降|hyperpage}
\begin{definition}{}{2.1.1}
$k$ 上的预叠是一个伪函子
\[y:\operatorname{Aff}_k^{\mathrm{op}}\longrightarrow\operatorname{Grpd}.\]
因此，每个仿射 $k$-概形 $S$ 给出一个群胚 $y(S)$，每个映射 $f:S_1\to S_2$ 给出拉回函子 $f^*:y(S_2)\to y(S_1)$，满足通常的相容性条件。若目标范畴是 $\operatorname{Set}$，这就是通常的预层概念。
\end{definition}

设 $\tau$ 是 $\operatorname{Aff}_k$ 上的 Grothendieck 拓扑，例如 Zariski、étale、光滑、fppf 或 fpqc 拓扑。
\begin{definition}{}{2.1.2}
若预叠 $y$ 对 $\tau$-覆盖满足下降，则称它为 $\tau$-叠。等价地，对每个 $\tau$-覆盖 $S^0\to S$ 及其 Čech 神经 $S^\bullet$，自然函子
\[y(S)\longrightarrow\lim y(S^\bullet)\]
是群胚的等价。
\end{definition}
\begin{note}\numberednote{2.1.3}\textbf{（分类叠）}
对群概形 $G/k$，$\tau$-分类叠 $B_\tau G$ 是商预叠 $\mathrm{pt}/G$ 的叠化。具体地，
\[B_\tau G(S)=\{S\text{ 上 }\tau\text{-局部平凡的 }G\text{ 挠子}\}.\]
它可以写为 $\operatorname{Stk}_k$ 中条形构造的叠论余极限：
\[B_\tau G\simeq\operatorname*{colim}(\mathrm{pt}\leftarrow G\leftleftarrows G\times G\cdots).\]
若 $G$ 在 $k$ 上光滑，则挠子在 étale 拓扑下就已局部平凡，因此
\[B_{\mathrm{et}}G=B_{\mathrm{sm}}G=B_{\mathrm{fppf}}G=B_{\mathrm{fpqc}}G.\]
\end{note}
\section{\texorpdfstring{$G$}{G} 丛的模叠}\label{sec:2.2}
\index{模叠!G丛|hyperpage}\index{Artin 叠|hyperpage}
设 $G/k$ 是光滑仿射代数群（在下文几何 Langlands 的应用中为连通约化群），$p:X\to S$ 是射影、平坦且有限表现的态射。定义
\[\Bun_{G,X/S}:=\operatorname{Maps}_S(X,B_{\mathrm{et}}G).\]
对任意 $S'\to S$，
\[\Bun_{G,X/S}(S')=\Hom(X\times_S S',B_{\mathrm{et}}G)=\{X\times_S S'\text{ 上的 étale }G\text{ 挠子}\}.\]
当 $S=\Spec k$ 时，简记为 $\Bun_G$。
\begin{theorem}{}{2.2.1}
叠 $\Bun_{G,X/S}$ 是在 $S$ 上局部有限表现的 Artin 叠，且具有仿射对角。更明确地，
\[\Delta_{\Bun_{G,X/S}}:\Bun_{G,X/S}\longrightarrow\Bun_{G,X/S}\times_S\Bun_{G,X/S}\]
是仿射态射，且 $\Bun_{G,X/S}$ 有开子叠穷竭 $\Bun_{G,X/S}=\bigcup_r U_r$，使每个 $U_r\to\Bun_{G,X/S}$ 都是开嵌入，每个 $U_r$ 都有光滑图册 $V_r\to U_r$，且 $V_r\to S$ 有限表现。
\end{theorem}
\begin{note}\numberednote{2.2.2}
若 $G/k$ 光滑且 $X\to S$ 是相对曲线，则 $\Bun_{G,X/S}\to S$ 光滑。在上述穷竭中，可以选择使 $V_r\to S$ 光滑的图册。
\end{note}
\srcpage{13}
\begin{note}\numberednote{2.2.3}\textbf{（$G=\GL_n$）}
固定 $X$ 上的丰沛线丛 $\cO(1)$。对 $r\gg0$，令 $U_r\subset\Bun_n$ 为由满足
\[\mathrm R^ip_*(\cF(r))=0\quad(i>0),\qquad p^*p_*(\cF(r))\twoheadrightarrow\cF(r)\]
的秩 $n$ 向量丛 $\cF$ 组成的开子叠。开子叠 $U_r$ 覆盖 $\Bun_n$。不过，$\Bun_G$ 通常不拟紧；对 $G=\GL_n$，它不连通，甚至一个固定的分支 $\Bun_G^\lambda$ 也未必拟紧。
\end{note}

以下阿代尔描述是 $\Bun_G$ 在函数论中的体现。令 $K=k(X)$，$K_x$ 为 $x\in X$ 处的完备局部域，并置 $\A_X=\prod_x'K_x$，$\cO_X=\prod_x\cO_x$。
\begin{theorem}{$k$-点上的阿代尔一致化}{2.2.4}
假设 $\Spec K$ 上、每个 $\Spec K_x$ 上以及每个有限扩张 $k'/k$ 的 $\Spec k'$ 上的所有 $G$ 挠子都是平凡的。则采用右挠子的约定，自然的粘合构造将 $\Bun_G$ 的 $k$-点同构类等同于双商
\[G(\cO_X)\backslash G(\A_X)/G(K).\]
采用左挠子时，同一陈述写成通常的相反次序 $G(K)\backslash G(\A_X)/G(\cO_X)$。
\end{theorem}
\begin{note}\numberednote{2.2.5}
当 $k$ 有限时，平凡性假设特别适用于单连通半单群，例如 $\SL_2$，以及 $\GL_n$。当 $k=\bar k$ 且特征为 $0$ 时，它适用于连通线性代数群 $G$；当 $k=\bar k$ 且特征为 $p$ 时，它适用于连通约化群 $G$。
\end{note}

本章其余部分介绍整体与局部 Hecke 叠，回顾定义 Hecke 函子的 Satake 层，然后专门讨论 $\GL_n$，以获得微小 Hecke 特征性质。
\section{整体 Hecke 叠}\label{sec:2.3}\index{Hecke 叠!整体|hyperpage}
从现在起假设 $X$ 是 $k$ 上的光滑曲线。整体 Hecke 叠记录 $G$ 丛在 $X$ 上一个移动点处的一次修改。
\begin{definition}{}{2.3.1}
叠 $\Hecke_{G,X}$ 定义为
\[
\begin{aligned}
\Hecke_{G,X}(S)=\{(x,P,P',\beta)\mid{}&x:S\to X,\ P,P'\in\Bun_G(S),\\
&\beta:P|_{X\times S-\Gamma_x}\xrightarrow{\sim}P'|_{X\times S-\Gamma_x}\},
\end{aligned}
\]
其中 $\Gamma_x\subset X\times S$ 是 $x$ 的图像。它有如下对应图：
\origdiagram{p013-diagram.pdf}{0.5146}
这里，$\overleftarrow h$ 记住 $P$，$\overrightarrow h$ 记住 $P'$，$\supp$ 记住点 $x$。
\end{definition}
\section{仿射 Grassmann 簇与局部模型}\label{sec:2.4}\index{仿射 Grassmann 簇|hyperpage}
$X$ 上的 Beilinson–Drinfeld 仿射 Grassmann 簇是 Hecke 叠的局部模型：
\[
\begin{aligned}
\Gr_{G,X}(S)=\{(x,P,\beta)\mid{}&x:S\to X,\ P\text{ 为 }X\times S\text{ 上的 }G\text{ 挠子},\\
&\beta:P|_{X\times S-\Gamma_x}\xrightarrow{\sim}P^0|_{X\times S-\Gamma_x}\},
\end{aligned}
\]
其中 $P^0$ 是平凡 $G$ 挠子。
\begin{theorem}{}{2.4.1}
$\Gr_{G,X}$ 是 ind-局部有限型的 ind-概形。若 $G$ 约化，则 $\Gr_{G,X}$ 在 $X$ 上 ind-射影。
\end{theorem}
\srcpage{14}

固定 $x:S\to X$。令 $D_x=\widehat{(X\times S)}_{\Gamma_x}$ 为沿图像的形式圆盘，$D_x^\circ=D_x\setminus\Gamma_x$ 为穿孔形式圆盘。定义正环路预叠和全环路预叠：
\[L^+G_X(S)=\{(x,\gamma)\mid\gamma:D_x\to G\},\qquad LG_X(S)=\{(x,\gamma)\mid\gamma:D_x^\circ\to G\}.\]
若 $S=\Spec k$ 且 $x$ 有参数 $t$，则 $D_x=\Spec k[[t]]$，$D_x^\circ=\Spec k((t))$。分类叠 $BL^+G_X$ 和 $BLG_X$ 分别分类 $D_x$ 与 $D_x^\circ$ 上的 $G$ 挠子。仿射 Grassmann 簇是商
\[\Gr_{G,X}\simeq LG_X/L^+G_X.\]

为了从这个局部模型过渡到整体修改，引入无限级丛
\[\Bun_G^\infty(S)=\{(x,P,\alpha)\mid x:S\to X,\ P\in\Bun_G(S),\ \alpha:P|_{D_x}\xrightarrow{\sim}P^0|_{D_x}\}.\]
它是 $X\times\Bun_G$ 上的 $L^+G_X$ 挠子。因此，从直观上以及精确的商意义上都有
\[\Hecke_{G,X}\simeq\Bun_G^\infty\times^{L^+G_X}\Gr_{G,X}\simeq\Gr_{G,X}\widetilde\times\Bun_G.\]
从而 $\Hecke_{G,X}$ 是 ind-局部有限型的 ind-Artin 叠。

\begin{transnote}{含糊点拨}
「无限级」与维数无关：$\alpha$ 定义在形式圆盘 $D_x=\Spec k[[t]]$（而非某阶 $D_x^{(n)}=\Spec k[t]/(t^n)$）上，因此 $\Bun_G^\infty$ 是有限级挠子 $\Bun_G^{(n)}$ 的逆向极限，只是 ind-对象，并不对应概形。扭曲积 $\Bun_G^\infty\times^{L^+G_X}\Gr_{G,X}$ 的含义是：一点由「$L^+G_X$-挠子 $\alpha$」与「$\Gr$ 的一点，即 $D_x^\circ$ 上的 $G$-丛配平凡化」组成，模去 $L^+G_X$ 的右作用后，正是「修改前的 $P$、修改后的 $P'$、以及 $D_x^\circ$ 上的同构」这三元组，也就是 $\Hecke_{G,X}$ 的定义。所以第二个同构并非巧合：$\Gr_{G,X}\widetilde\times\Bun_G$ 与扭曲积是同一叠的两种写法。
[来源：nLab "affine Grassmannian"，https://ncatlab.org/nlab/show/affine+Grassmannian（Last revised 2023-08-15；2026-11-17 访问）；并见原书 p.14 对 $D_x$ 与 $L^+G_X$ 的约定。]
\end{transnote}
\section{局部 Hecke 叠与几何 Satake}\label{sec:2.5}\index{Hecke 叠!局部|hyperpage}\index{几何 Satake|hyperpage}
对固定点 $x\in X$，记 $L^+G_x=G(\cO_x)$，$LG_x=G(K_x)$。局部 Hecke 叠为
\[\Hecke_{G,x}^{\mathrm{loc}}:=[L^+G_x\backslash LG_x/L^+G_x]\simeq[L^+G_x\backslash\Gr_{G,x}].\]
将一个修改在 $x$ 处完备化，便得到从整体 Hecke 叠到局部 Hecke 叠的相对位置映射
\[m_x:\Hecke_{G,x}\longrightarrow\Hecke_{G,x}^{\mathrm{loc}}.\]

几何 Satake 将局部 Hecke 叠上的球面层范畴等同于 Langlands 对偶群的表示范畴：
\[\operatorname{Perv}_{L^+G_x}(\Gr_{G,x})\simeq\operatorname{Perv}(\Hecke_{G,x}^{\mathrm{loc}})\simeq\Rep(\checkG).\]
对 $V\in\Rep(\checkG)$，以 $\operatorname{Sat}_V$ 表示 $\Hecke_{G,x}^{\mathrm{loc}}$ 上相应的反常层；在 $X$ 上的整体族中，也用同一记号表示相应的相对 Satake 层。

\begin{transnote}{背景直觉}
这句话的内容比「认同」更强：$\operatorname{Perv}_{L^+G}(\Gr_G)$ 上由卷积（即融合积）给出的张量结构，与 $\Rep(\checkG)$ 的张量结构等价；对应把 $\operatorname{IC}(\overline{\Gr^\lambda})$ 送到最高权表示 $V^\lambda$，其中 $\lambda$ 跑优势余权、$\overline{\Gr^\lambda}$ 是 $\Gr^\lambda$ 的闭包。$\Hecke^{\mathrm{loc}}$ 上的球面层与 $\Gr$ 上的 $L^+G$-等变反常层是同一批对象，「球面」指按 $L^+G$ 的右作用取不变。$\ell$-进情形还须用 Weil 结构处理 Frobenius 归一化，Betti 与 de Rham 情形则直接是反常层。这句话同时给出了 $\operatorname{Sat}_V$ 的定义：$V$ 的不可约分解对应相应轨道 IC 层的直和。
[来源：nLab "geometric Satake equivalence"，https://ncatlab.org/nlab/show/geometric+Satake+equivalence（Last revised 2024-07-22；2026-11-17 访问）；并见 [FGKV98]（JAMS 11 (1998), 451–484）。]
\end{transnote}
\section{Hecke 函子与特征层}\label{sec:2.6}\index{Hecke 函子|hyperpage}
对 $V\in\Rep(\checkG)$，整体 Hecke 函子为
\[
\begin{aligned}
H_V:D(\Bun_G)&\longrightarrow D(X\times\Bun_G),\\
\cF&\longmapsto(\supp,\overrightarrow h)_!(m^*\operatorname{Sat}_V\otimes\overleftarrow h^*\cF).
\end{aligned}
\]
在固定点 $x\in X$ 处，它限制为
\[
\begin{aligned}
H_{V,x}:D(\Bun_G)&\longrightarrow D(\Bun_G),\\
\cF&\longmapsto\overrightarrow h_{x,!}(m_x^*\operatorname{Sat}_V\otimes\overleftarrow h_x^*\cF).
\end{aligned}
\]
$\operatorname{Sat}_V$ 的反常归一化包括通常的移位与扭曲。

设 $E_{\checkG}$ 是 $X$ 上的 $\checkG$-局部系统。对 $V\in\Rep(\checkG)$，令 $V_E$ 为 $X$ 上相应的局部系统。
\begin{definition}{}{2.6.1}
若对每个 $V\in\Rep(\checkG)$，在 $D(X\times\Bun_G)$ 中都有函子性的同构
\[H_V(\cF)\simeq V_E\boxtimes\cF,\]
并且这些同构与 $\Rep(\checkG)$ 中的张量积相容，则称对象 $\cF\in D(\Bun_G)$ 是以 $E_{\checkG}$ 为特征值的 Hecke 特征层。特别地，在 $x\in X$ 处逐纤维有
\[H_{V,x}(\cF)\simeq(V_E)_x\otimes\cF.\]
\end{definition}
\srcpage{15}

这是在层论层面替代经典条件：一个自守函数应是所有球面 Hecke 算子的特征向量。
\section{\texorpdfstring{$\GL_n$}{GLₙ} 的微小 Hecke 算子}\label{sec:2.7}\index{微小 Hecke 算子|hyperpage}
对 $G=\GL_n$，有 $\checkG=\GL_n$。令 $\mathrm{Std}$ 为 $\checkG$ 的标准表示。外幂 $\bigwedge^i\mathrm{Std}$，$0\le i\le n$，在几何 Satake 下对应于微小 Schubert 轨道 $\Gr_G^i\subset\Gr_G$。这些是与标准 Hecke 算子有关的闭球面轨道，给出以 $0\le i\le n$ 为指标的 $n+1$ 个标准算子。若保留 $\GL_n$ 的中心次数移位，则可以用中心余权平移这些轨道；模掉这一中心方向后，$\omega_0$ 与 $\omega_n$ 被等同，因而有 $n$ 个类。

令 $\Hecke_n^i$ 为相应的整体 Hecke 对应，$\Gr_{n,X}^i$ 为 $\Gr_{\GL_n,X}$ 中相应的相对 Schubert 分层。一个 $S$-点是向量丛在 $x:S\to X$ 处的修改，其局部相对位置为 $\omega_i$；局部而言，这等价于一个支撑在 $\Gamma_x$ 上、长度为 $i$ 的商。局部模型是通常的 Grassmann 簇 $\Gr(i,n)$，且
\[\Hecke_n^i\simeq\Gr_{n,X}^i\widetilde\times\Bun_n.\]

\begin{transnote}{含糊点拨}
相对位置 $\omega_i=(1,\ldots,1,0,\ldots,0)$ 是 $L^+G$-轨道的标号：在一点取基后，修改映射在圆盘上可写成 $t^{\omega}$ 的矩阵形（Birkhoff 分解），$\omega_i$ 恰对应「余核为 $i$ 维」的那条轨道。$\Hecke_n^i$ 只取这一条，所以一个 $S$-点就是 $\mathcal V'\subset\mathcal V$ 而商支撑在 $\Gamma_x$ 上、长度 $i$；在代数闭点上商是 $k^{\oplus i}$，长度为 $i$ 的商的参数空间与 $i$ 维子空间的参数空间（即 $\Gr(i,n)$）经取对偶相同——这就是「局部模型是通常的 Grassmann 簇」的确切含义。整条轨道给出的参数空间是 $\Gr_{n,X}^i$，于是 $\Hecke_n^i\simeq\Gr_{n,X}^i\widetilde\times\Bun_n$ 逐点地看就是「底丛加上这个参数」。
[来源：nLab "affine Grassmannian"，https://ncatlab.org/nlab/show/affine+Grassmannian（Last revised 2023-08-15；2026-11-17 访问）；并见原书 p.15 对 $\omega_i$ 与 $\Gr(i,n)$ 的说明。]
\end{transnote}
相应的 Hecke 函子为
\[H_n^i:D(\Bun_n)\longrightarrow D(X\times\Bun_n),\]
\[H_n^i(\cF)=(\supp,\overrightarrow h^i)_!\overleftarrow h^{i,*}\cF[\text{移位}],\]
其中移位采用 Satake／反常归一化；对微小轨道，可以取移位 $\dim\Gr(i,n)=i(n-i)$，在 $\ell$-进归一化中配上相应的 Tate 扭转。

\begin{transnote}{易错预警}
两个因子各司其职，不能只取其一。移位 $[d_i]$ 取相对维数 $d_i=i(n-i)=\dim\Gr(i,n)$，使 $\Gr(i,n)$ 上的常值层经平移后落在反常层的 $t$-结构中；Tate 扭转 $(d_i/2)$ 再把 Frobenius 权拉回零（$\mathbb Q_\ell[d_i](d_i/2)$ 是权 $0$ 的）。两者合起来才使下面 (2.1) 的 $H_n^i(\cF)\simeq\bigwedge^iE\boxtimes\cF$ 两侧归一化一致：只移不扭，$\ell$-进情形会多出 $q_x^{d_i/2}$ 的因子；只扭不移，Hecke 函子不再 $t$-正合。
[来源：nLab "perverse sheaf"，https://ncatlab.org/nlab/show/perverse+sheaf（Last revised 2026-10-16；2026-11-17 访问）；nLab "Tate twist"，https://ncatlab.org/nlab/show/Tate+twist（Last revised 2025-08-10；2026-11-17 访问）；并见原书 p.16 式 (2.1) 的归一化。]
\end{transnote}

于是，以秩 $n$ 局部系统 $E$ 为特征值的 Hecke 特征层 $\cF$ 满足
\begin{equation}\label{eq:2.1}
H_n^i(\cF)\simeq\bigwedge^iE\boxtimes\cF,\qquad0\le i\le n.
\end{equation}
其余小节将对由 $E$ 构造的自守层证明这个公式。

现在给出 $\GL_n$ 情形的详细论证。关键是：第一 Hecke 特征性质，加上交换不同点时通常的对称性，就能推出式\eqref{eq:2.1}中所有微小特征性质。这是 Frobenius 特征值的初等对称函数关系的层论形式。
\section{微小 Hecke 对应}\label{sec:2.8}
令 $\Bun_n:=\Bun_{\GL_n}$。对 $0\le i\le n$，置 $\omega_i=(1^i,0^{n-i})$，从而 $t^{\omega_i}=\operatorname{diag}(tI_i,I_{n-i})$，并令
\[\Gr_n^i:=L^+\GL_n\cdot t^{\omega_i}\cdot L^+\GL_n/L^+\GL_n.\]
令 $\Gr_{n,X}^i$ 为 $X$ 上相应的相对 Schubert 分层。相应的整体 Hecke 叠为
\[\Hecke_n^i\simeq\Gr_{n,X}^i\widetilde\times\Bun_n,\]
并具有映射
\origdiagram{p015-diagram.pdf}{0.4303}
一个 $S$-点分类点 $x:S\to X$ 和向量丛的嵌入 $\mathcal V'\hookrightarrow\mathcal V$，使 $\mathcal V/\mathcal V'$ 是支撑于 $\Gamma_x$ 上、相对位置为 $\omega_i$ 的长度 $i$ 挠层。在一个代数闭点上，这一条件是 $\mathcal V/\mathcal V'\simeq k(x)^{\oplus i}$；等价地，$\mathcal V'\subset\mathcal V\subset\mathcal V'(x)$ 且 $\length(\mathcal V/\mathcal V')=i$。
\srcpage{16}

置 $d_i=i(n-i)$。这是 $\Hecke_{n,x}^i\to\Bun_n$ 的相对维数。归一化的 Hecke 函子为
\[H_n^i(\cF):=(\supp,\overrightarrow h)_!\overleftarrow h^*\cF[d_i](d_i/2)\in D(X\times\Bun_n).\]
以 $E$ 为特征值的第 $i$ 个 Hecke 特征性质是
\[H_n^i(\cF)\simeq\bigwedge^iE\boxtimes\cF.\]
\section{轨道记号与闭包序}\label{sec:2.9}\index{Schubert 轨道|hyperpage}\index{优势序|hyperpage}
整体 Hecke 叠通过用 $X\times\Bun_G$ 上的无限级 $L^+G$ 挠子扭曲仿射 Grassmann 簇而得到：
\[\Hecke_{G,X}\simeq\Gr_{G,X}\widetilde\times\Bun_G.\]
因此，每个 $L^+G$ 轨道 $\Gr_G^\lambda\subset\Gr_G$ 都给出相应的相对分层 $\Gr_{G,X}^\lambda\subset\Gr_{G,X}$，以及
\[\Hecke_G^\lambda\simeq\Gr_{G,X}^\lambda\widetilde\times\Bun_G.\]
对 $G=\GL_n$，优势余权为元组 $\lambda=(\lambda_1\ge\cdots\ge\lambda_n)$，且
\[\Gr_n^\lambda:=L^+\GL_n\cdot t^\lambda\cdot L^+\GL_n/L^+\GL_n,\qquad t^\lambda=\operatorname{diag}(t^{\lambda_1},\ldots,t^{\lambda_n}).\]
在一个固定的连通分支内，闭包序就是优势序：$\Gr_n^\lambda\subset\overline{\Gr_n^\mu}$ 当且仅当对所有 $1\le r<n$ 都有 $\sum_{j=1}^r\lambda_j\le\sum_{j=1}^r\mu_j$。微小余权 $\omega_i=(1^i,0^{n-i})$ 在几何 Satake 下对应于 $\bigwedge^i\mathrm{Std}$。
\section{经典自守归一化}\label{sec:2.10}
令
\[G_n^i=L^+\GL_n\cdot\operatorname{diag}(\underbrace{t,\ldots,t}_{i},1,\ldots,1)\cdot L^+\GL_n/L^+\GL_n.\]
对球面函数 $f\in C_c^\infty(L^+\GL_n\backslash L\GL_n/L^+\GL_n)$，归一化的局部 Hecke 算子为
\[T_x^i(f)(y)=q_x^{-d_i/2}\int_{G_n^i}f(yg)\,dg.\]
若局部系统 $E$ 对应 Weil 表示 $\rho:W_F\to\GL_n(\mathbb Q_\ell)$，则预期特征值为 $\operatorname{tr}(\operatorname{Frob}_x;\bigwedge^i\rho)$。等价地，
\[\det(1-\rho(\operatorname{Frob}_x)s)=\sum_{i=0}^n(-1)^i\operatorname{tr}(\operatorname{Frob}_x;\bigwedge^i\rho)s^i.\]
这是 $H_n^i(\cF)\simeq\bigwedge^iE\boxtimes\cF$ 在函数论中的体现。
\section{归约到第一 Hecke 算子}\label{sec:2.11}
假设 $\cF\in\operatorname{Perv}(\Bun_n)$ 配备同构
\[\alpha:H_n^1(\cF)\xrightarrow{\sim}E\boxtimes\cF.\]
对两个不同的点 $x\ne x'$，两次初等修改的叠有 $S_2$ 作用：考虑 $\mathcal V''\subset\mathcal V'\subset\mathcal V$，满足 $\length(\mathcal V/\mathcal V')=\length(\mathcal V'/\mathcal V'')=1$、$\supp(\mathcal V/\mathcal V')=x$、$\supp(\mathcal V'/\mathcal V'')=x'$。交换 $x$ 与 $x'$ 就将这两次初等修改的次序颠倒。特征层所要求的相容性是同构
\[H_n^1H_n^1(\cF)|_{X^2-\Delta}\xrightarrow{\sim}E\boxtimes E\boxtimes\cF|_{X^2-\Delta}\]
为 $S_2$-等变的。迭代后，在 $X^i-\Delta$ 上得到 $S_i$-等变同构
\[(H_n^1)^i(\cF)|_{X^i-\Delta}\simeq E^{\boxtimes i}\boxtimes\cF|_{X^i-\Delta}.\]
\srcpage{17}
\begin{proposition}{}{2.11.1}
在上述对称性假设下，对每个 $0\le i\le n$，都有
\[H_n^i(\cF)\simeq\bigwedge^iE\boxtimes\cF.\]
\end{proposition}
\section{几何 Satake 的输入}\label{sec:2.12}
回顾
\[\operatorname{Perv}_{L^+G}(\Gr_G)\simeq\operatorname{Perv}(L^+G\backslash LG/L^+G)\simeq\Rep(\checkG),\]
沿这一等价，对所有 $\lambda$ 都有
\[\operatorname{IC}_{\overline{\Gr_G^\lambda}}\longleftrightarrow V^\lambda.\]
对 $G=\GL_n$，这给出
\[\mathrm{Std}\longleftrightarrow(1,0,\ldots,0)=(1,0^{n-1}),\qquad\bigwedge^i\mathrm{Std}\longleftrightarrow(1^i,0^{n-i}).\]
因此，$\operatorname{IC}_{\Gr_n^1}^{\star i}$ 对应于 $\mathrm{Std}^{\otimes i}$，自然 $S_i$ 作用的交错投影算子切出 $\bigwedge^i\mathrm{Std}$。整体证明通过长度 $i$ 修改叠的 Springer 型解消来实现这一投影算子。
\section{长度 \texorpdfstring{$i$}{i} 修改的紧化叠}\label{sec:2.13}\index{Springer 解消|hyperpage}\index{Springer 层|hyperpage}
令
\[\operatorname{Mod}_n^{-i}:=\{(\mathcal V'\hookrightarrow\mathcal V)\mid\mathcal V/\mathcal V'\in\Tor^i\}.\]
它有对应
\origdiagram{p017-diagram.pdf}{0.4258}
其中 $\supp(\mathcal V/\mathcal V')$ 是相应的有效除子。令 $\Delta_i:X\to X^{(i)}$ 为 $x\mapsto i\cdot x$，并置
\[H_n^{i,+}:=X\times_{X^{(i)}}\operatorname{Mod}_n^{-i}.\]
满足 $\mathcal V/\mathcal V'\simeq k(x)^{\oplus i}$ 的层正是 $\Hecke_n^i$。因此，$H_n^{i,+}$ 通过允许支撑在 $x$ 处的非半单长度 $i$ 商，将微小对应紧化。

对 $I=\{1,\ldots,i\}$，Beilinson–Drinfeld Grassmann 簇给出
\[\Gr_{G,X}^I\widetilde\times\Bun_G=\Hecke_{G,X}^I.\]
对 $G=\GL_n$，存在从正的整体 Schubert 簇到紧化修改叠的真态射：
\[\Gr_{G,X}^{+,\le(1,\ldots,1)}\widetilde\times\Bun_n=\Hecke_{n,X}^{+,\le(1,\ldots,1)}\longrightarrow\operatorname{Mod}_n^{-i}.\]
在 $X^i-\Delta$ 上，它因子化为
\[\Gr_{G,X}^{+,\le(1,\ldots,1)}|_{X^i-\Delta}\simeq(\Gr_{G,x_1}^1\times\cdots\times\Gr_{G,x_i}^1)|_{X^i-\Delta}.\]

现在定义有序解消，也称 Springer 解消：
\[\widetilde{\operatorname{Mod}}_n^{-i}:=\{\mathcal V_0\hookrightarrow\mathcal V_1\hookrightarrow\cdots\hookrightarrow\mathcal V_i\mid\length(\mathcal V_j/\mathcal V_{j-1})=1\text{，对所有 }j\}.\]
存在真态射
\[\varphi:\widetilde{\operatorname{Mod}}_n^{-i}\longrightarrow\operatorname{Mod}_n^{-i},\qquad(\mathcal V_0\subset\cdots\subset\mathcal V_i)\longmapsto(\mathcal V_0\subset\mathcal V_i),\]
\srcpage{18}
以及有序支撑映射 $r:\widetilde{\operatorname{Mod}}_n^{-i}\to X^i$。它们组成交换图
\origdiagram{p018-square.pdf}{0.2755}
在互异点轨迹上，$\varphi$ 是有序支撑覆盖，并带有通过置换长度一商的次序给出的自然 $S_i$ 作用。
\begin{claim}{}{2.13.1}
Springer 层
\[\operatorname{Spr}_i:=\varphi_!\operatorname{IC}_{\widetilde{\operatorname{Mod}}_n^{-i}}\]
满足以下性质：
\begin{enumerate}
\item $\operatorname{Spr}_i$ 是反常层。
\item $\operatorname{Spr}_i$ 是从互异点轨迹作出的中间延拓；在该轨迹上，它是与有序支撑覆盖及 $S_i$ 的正则表示相联系的局部系统。
\item 交错直和项是微小核：
\[\Hom_{S_i}(\mathrm{sgn},\operatorname{Spr}_i)|_{\Hecke_n^i}\simeq\mathbb Q_\ell|_{\Hecke_n^i}[d_i](d_i/2),\qquad d_i=i(n-i).\]
\end{enumerate}
\end{claim}

这与通常 Springer 层的机制相同：解消为初等商排序，符号表示则提取外幂直和项。
\section{投影公式计算}\label{sec:2.14}
\enlargethispage{\baselineskip}
\begin{proof}
通向 $\Bun_n$ 的解消奇点映射为
\origdiagram{p018-triangle.pdf}{0.4258}
支撑映射组成如下支撑方块：
\origdiagram{p018-square.pdf}{0.2755}
由投影公式，
\[(\supp,\overrightarrow h)_!(\overleftarrow h^*\cF\otimes\operatorname{Spr}_i)\simeq(\mathrm{sym}\times\mathrm{id}_{\Bun_n})_!((H_n^1)^i(\cF)).\]
这里，归一化的移位与 Tate 扭转已经包含在 $\operatorname{Spr}_i$ 的 IC 归一化以及函子 $H_n^1$ 中。在 $X^i-\Delta$ 上，右边是迭代第一 Hecke 函子的对称化。利用 $S_i$-等变的特征同构以及中间延拓性质，得到
\[(\mathrm{sym}\times\mathrm{id}_{\Bun_n})_!((H_n^1)^i(\cF))\simeq\mathrm{sym}_!(E^{\boxtimes i})\boxtimes\cF.\]
取符号表示的等型直和项，再沿 $\Delta_i:X\to X^{(i)}$ 限制，便得到
\[\Delta_i^*\Hom_{S_i}(\mathrm{sgn},\mathrm{sym}_!E^{\boxtimes i})\simeq\bigwedge^iE.\]
\srcpage{19}
将这个恒等式与核的符号 Springer 描述结合，就得到
\[H_n^i(\cF)\simeq\bigwedge^iE\boxtimes\cF.\]
这就证明了命题\ref{pro:2.11.1}。
\end{proof}

\chapter{Whittaker}\label{ch:3}
% 原PDF第19页
本章把视角从 $\Bun_n$ 上的特征层转向 Whittaker 模型。函数论陈述是：给定球面 Hecke 特征值后，归一化的 Whittaker 向量唯一。其几何体现是 Satake 范畴的 Whittaker 模型；对 $\GL_n$，它恢复了由局部系统构造的尖点自守层\cite{FGKV98,Gai18}。
\section{回顾：Hecke 特征层与尖点性}\label{sec:3.1}\index{尖点性|hyperpage}\index{常数项函子|hyperpage}
令 $\Bun_n:=\Bun_{\GL_n}$。对 $0\le d\le n$，第 $d$ 个微小 Hecke 对应为
\origdiagram{p019-hecke.pdf}{0.3871}
一个 $S$-点是元组
\[(x,M,M',M'\subset M\subset M'(x)),\]
其中 $M$ 与 $M'$ 是秩 $n$ 向量丛，$M/M'$ 在 $x$ 的图像上有限局部自由，秩为 $d$。归一化的 Hecke 函子为
\[H_n^d(\cK):=(\supp,\overrightarrow h)_!\overleftarrow h^*\cK\left(\frac{d(n-d)}2\right)[d(n-d)]\in D(X\times\Bun_n).\]
对 $X$ 上的局部系统 $E$，若
\[H_n^d(\cK)\simeq\bigwedge^dE\boxtimes\cK,\]
则称 $\cK$ 满足以 $E$ 为特征值的第 $d$ 个 Hecke 特征性质。
如\chref{2}所述，若 $\cK$ 是反常层，并满足第一特征性质以及迭代修改通常的对称相容性，则它满足所有微小特征性质。

这里使用的主要存在性定理是 Frenkel–Gaitsgory–Vilonen 定理\cite{FGKV98}：若 $E$ 几何不可约，则 $\Bun_n$ 上存在尖点 Hecke 特征层 $\Aut_E$，且它在每个连通分支 $\Bun_n^d$ 上的限制不可约。

还要回顾尖点性的几何定义。设 $G$ 约化，$P\subsetneq G$ 是真抛物子群，其幂幺根基为 $U$，Levi 商为 $M=P/U$。则
\origdiagram{p019-constant.pdf}{0.3575}
定义常数项函子
\[\CT_!:=q_!p^*.\]
若对每个真抛物子群 $P$ 都有 $\CT_!(\cK)=0$，则称层 $\cK\in D(\Bun_G)$ 为尖点层。
\section{经典 Hecke 特征函数}\label{sec:3.2}
\srcpage{20}
以下函数论段落假设 $k$ 有限。令 $F=k(X)$，$\A$ 为 $F$ 的阿代尔环，并置 $G=\GL_n$。设
\[\sigma:\Gal(\overline F/F)\longrightarrow\GL_n(\mathbb Q_\ell)\]
处处非分歧，等价地，它是 $X$ 上秩 $n$ 的局部系统。球面自守函数是定义在
\[\GL_n(F)\backslash\GL_n(\A)/\GL_n(\cO)\]
上的函数。对闭点 $x\in|X|$，记 $K_x$ 为局部域，$\cO_x$ 为其整数环，$q_x=\#k(x)$，并选取素元 $t_x$。置
\[M_n^d(\cO_x):=\GL_n(\cO_x)t_x^{(1^d,0^{n-d})}\GL_n(\cO_x)\subset\GL_n(K_x)\subset\GL_n(\A).\]
未归一化的局部 Hecke 算子为
\[T_x^df(g)=\int_{M_n^d(\cO_x)}f(gh)\,dh.\]
若对每个 $x\in|X|$ 都有
\[T_x^df=q_x^{d(n-d)/2}\Tr(\bigwedge^d\sigma(\Fr_x))f,\]
则称 $f$ 满足关于 $\sigma$ 的第 $d$ 个 Hecke 特征性质。
等价地，令 $\gamma_x$ 为由
\[\Tr(\bigwedge^d\gamma_x)=q_x^{d(n-1)/2}\Tr(\bigwedge^d\sigma(\Fr_x)),\qquad0\le d\le n\]
刻画的归一化 Satake 参数。采用 Tate 扭转的约定时，记为 $\gamma_x=\sigma((n-1)/2)(\Fr_x)$。于是，同一特征值为
\[q_x^{-d(d-1)/2}\Tr(\bigwedge^d\gamma_x).\]

\begin{transnote}{含糊点拨}
$\gamma_x$ 与 $\sigma(\Fr_x)$ 只差一个标量：$(n-1)/2$ 是半和正根 $\rho$ 在基本余权 $\varpi_1$ 上取的值，也就是 Borel 子群模特征 $\delta_B$ 的平方根在 $t^{\varpi_1}$ 处的值 $q_x^{(n-1)/2}$。这正是 \S\ref{sec:3.3} 用 $\rho$-移位（「上点」$\dot G$）所做的事：把 Whittaker 特征标从环面扭一下，使 Satake 参数与 Langlands 参数对齐。至于为什么只需这些迹：$\Tr(\bigwedge^d\gamma_x)$，$d=0,\ldots,n$，给出 $\gamma_x$ 的 $n$ 个特征值的全部初等对称函数，因而这 $n+1$ 个等式唯一确定半单元素 $\gamma_x$。
[来源：nLab "Tate twist"，https://ncatlab.org/nlab/show/Tate+twist（Last revised 2025-08-10；2026-11-17 访问）；并见 [Laf18] 对归一化 Satake 参数的约定（原书 p.20）。]
\end{transnote}

对幂幺根基为 $U$ 的抛物子群 $P\subsetneq\GL_n$，若球面自守函数满足
\[\int_{U(F)\backslash U(\A)}f(ug)\,du=0,\qquad g\in\GL_n(\A),\]
则称它为尖点函数。
采用通常的有限阶行列式归一化，Lafforgue 定理将具有上述非分歧特征值的尖点自守表示赋给这样的不可约 $\sigma$\cite{Laf18}。球面向量在相差一个标量的意义下唯一；Whittaker 归一化确定这个标量。
\section{Whittaker 特征标与 \texorpdfstring{$\rho$}{ρ} 移位}\label{sec:3.3}\index{Whittaker 特征标|hyperpage}\index{rho移位@$\rho$ 移位|hyperpage}
令 $N\subset\GL_n$ 为上三角幂幺子群。固定非平凡加性特征标 $\psi:k\to\Qell^\times$。在 $x$ 处，标准非退化 Whittaker 特征标通过简单根坐标得到：
\[
\begin{aligned}
\dot N(K_x)&\longrightarrow(\dot N/[\dot N,\dot N])(K_x)
\xrightarrow{\sim}\bigoplus_{\alpha\in\Delta}K_x\otimes\Omega_x\\
&\xrightarrow{\mathrm{Res}}\bigoplus_{\alpha\in\Delta}k(x)
\xrightarrow{\Tr_{k(x)/k}\circ\Sigma}k\xrightarrow{\psi}\Qell^\times.
\end{aligned}
\]
这里 $\Omega_x$ 是 $x$ 处完备化的典范线。所得特征标记为
\[\Psi_x:\dot N(K_x)\longrightarrow\Qell^\times,\]
其中上点表示 $\rho$ 移位。对 $G=\GL_n$，局部可以写为
\[\dot G_n(K_x)=\{(a_{ij})\mid a_{ij}\in K_x\otimes\Omega_x^{\otimes(j-i)}\}.\]
\srcpage{21}
因此，第一条上副对角线上的元素属于 $K_x\Omega_x$，可以对其简单根坐标取留数。$\dot G(K_x)$ 上的 Whittaker 函数是满足
\[W(ug)=\Psi_x(u)W(g),\qquad u\in\dot N(K_x)\]
的函数 $W$。通常还要求右 $\dot G(\cO_x)$ 不变性。
\section{局部 Whittaker 特征函数}\label{sec:3.4}\index{Casselman–Shalika 公式|hyperpage}
设 $\gamma\in\GL_n(\mathbb Q_\ell)$，例如上述归一化 Satake 参数 $\gamma_x$。在 $\dot G(K_x)/\dot G(\cO_x)$ 上存在唯一归一化的 Hecke 特征 Whittaker 函数 $W_{\gamma,x}$，满足：
\begin{enumerate}
\item $W_{\gamma,x}(1)=1$，其中 $1$ 表示仿射 Grassmann 簇的基点；
\item 对 $h\in\dot G(\cO_x)$，$W_{\gamma,x}(gh)=W_{\gamma,x}(g)$；
\item 对 $u\in\dot N(K_x)$，$W_{\gamma,x}(ug)=\Psi_x(u)W_{\gamma,x}(g)$；
\item 对每个 $0\le d\le n$，$T_x^dW_{\gamma,x}=q_x^{-d(d-1)/2}\Tr(\bigwedge^d\gamma)W_{\gamma,x}$。
\end{enumerate}
选取用来书写 $t_x^\lambda$ 的局部坐标后，Whittaker 双商以余权格为指标：
\[\dot N(K_x)\backslash\dot G(K_x)/\dot G(\cO_x)\simeq\Lambda,\qquad\dot N(K_x)t_x^\lambda\dot G(\cO_x)\longleftrightarrow\lambda.\]
Casselman–Shalika 公式明确给出这些值；其几何形式是 Satake 的 Whittaker 实现\cite{FGKV98}。对余权 $\lambda=(\lambda_1,\ldots,\lambda_n)$，
\[W_{\gamma,x}(t_x^\lambda)=\begin{cases}q_x^{-\sum_i(i-1)\lambda_i}\Tr(\gamma,V^\lambda),&\lambda\text{ 优势},\\0,&\text{否则},\end{cases}\]
其中，对优势 $\lambda$，$V^\lambda$ 表示最高权为 $\lambda$ 的不可约 $\checkG$ 表示。在范畴层面，这是几何 Satake 的 Whittaker 形式：
\[\operatorname{Perv}^{\mathrm{Wh}}(\dot{\Gr}_{G,x})\simeq\Rep(\checkG),\]
在导出层面则为 $D^{\mathrm{Wh}}(\dot{\Gr}_{G,x})\simeq D^b(\Rep(\checkG))$。
\section{整体 Whittaker 函数与自守函数}\label{sec:3.5}
局部特征标相乘得到整体特征标
\[\Psi=\prod_x\Psi_x:\dot N(\A)\longrightarrow\Qell^\times,\]
由留数定理，它在 $\dot N(F)$ 上平凡。上述局部定理给出定义在 $(\dot N(\A),\Psi)\backslash\dot G(\A)/\dot G(\cO)$ 上的限制乘积
\[W_\sigma:=\prod_x'W_{\gamma_x,x}.\]
$\dot G(\A)$ 上存在唯一归一化的非分歧 Hecke 特征 Whittaker 函数，由 $W(1)=1$ 及以下条件刻画：
\begin{enumerate}
\item 对 $h\in\dot G(\cO)$，$W(gh)=W(g)$；
\item 对 $u\in\dot N(\A)$，$W(ug)=\Psi(u)W(g)$；
\item $T_x^dW=q_x^{d(n-d)/2}\Tr(\bigwedge^d\sigma(\Fr_x))W$。
\end{enumerate}

对 $\GL_n$，Whittaker 展开首先给出 mirabolic 商上的尖点函数。令
\[Q_1=\begin{pmatrix}\GL_{n-1}&V_{n-1}\\0&1\end{pmatrix}\]
为 mirabolic 子群。典范映射为
\srcpage{22}
\[
\begin{aligned}
\varphi:C^\infty((\dot N(\A),\Psi)\backslash\dot G(\A))&\xrightarrow{\sim}C^\infty(\dot Q_1(F)\backslash\dot G(\A))_{\mathrm{cusp}},\\
W&\longmapsto\left(g\longmapsto\sum_{\gamma\in\dot N(F)\backslash\dot Q_1(F)}W(\gamma g)\right).
\end{aligned}
\]
等价地，该和就是相应离散商上的积分。对归一化的 Hecke 特征函数 $W_\sigma$，Piatetski-Shapiro–Shalika 自守性定理说明，$\varphi(W_\sigma)$ 实际上左 $\dot G(F)$ 不变；这就是尖点 Hecke 特征函数 $f_\sigma$。
\section{几何预览：Whittaker 范畴}\label{sec:3.6}\index{Whittaker 范畴|hyperpage}\index{手征积分|hyperpage}
上述局部陈述具有以下几何形式\cite{Gai18}：
\[\operatorname{Wh}(\dot{\Gr}_{G,x})\simeq\Rep(\checkG).\]
在曲线上取手征积分后，预期有
\[\int_X^{\mathrm{ch}}\operatorname{Wh}(\dot{\Gr}_{G,x})\simeq\int_X^{\mathrm{ch}}\Rep(\checkG).\]
左边是对阿代尔 Whittaker 商的层论替代，右边则映到 $\checkG$-局部系统叠上的拟凝聚层。图景为
\origdiagram{p022-diagram.pdf}{0.6648}
因此，Whittaker 模型是在 $\Bun_G$ 上的 Hecke 特征层与 $\LocSys_{\checkG}$ 上的层之间建立联系的局部到整体桥梁。

为进行整体比较，以 $K=k(X)$ 表示整体函数域，$\A$ 表示其阿代尔环。在函数论陈述中 $k$ 有限；几何上，将 $\psi$ 替换为 Artin–Schreier 层。令 $G=\GL_n$，$N\subset B\subset G$ 为标准的极大幂幺子群和 Borel 子群，并固定非平凡加性特征标 $\psi:k\to\Qell^\times$。为简化记号，此处省去上点：只要出现 Whittaker 特征标 $\Psi$，阿代尔群就是上文所述的 $\rho$ 扭曲形式。令
\[Q_1=\begin{pmatrix}\GL_{n-1}&V\\0&1\end{pmatrix}\]
为 mirabolic 子群。经典的目标是 Whittaker 展开映射，这一目标在文献\cite{FGKV98,Gai18}的 Whittaker 形式体系中被几何化：
\[\Av_!:C^\infty((N(\A),\Psi)\backslash G(\A))\xrightarrow{\sim}C^\infty(Q_1(K)\backslash G(\A))_{\mathrm{cusp}}.\]
它将整体 Whittaker 函数 $W_\sigma=\prod_x'W_{\gamma_x,x}$ 映为 mirabolic 商上的尖点函数 $f_\sigma$。这里 $\gamma_x$ 是上文引入的归一化 Satake 参数，局部函数 $W_{\gamma_x,x}\in C^\infty((N(K_x),\Psi_x)\backslash G(K_x)/G(\cO_x))$ 是
\srcpage{23}
Casselman–Shalika 定理给出的归一化局部 Whittaker 函数。剩下的自守性陈述是：所得 $f_\sigma$ 不仅 $Q_1(K)$ 不变，而且实际 $G(K)$ 不变。

等价地，由于 $N\subset Q_1$，有
\[(\Av_!W)(g)=\sum_{\gamma\in N(K)\backslash Q_1(K)}W(\gamma g).\]
在函数论情形，尖点性使该求和收敛。几何上，以一般位置约化模空间之间相应映射的 $!$ 推前替代它。
\section{一般位置约化与 Whittaker 特征标}\label{sec:3.7}\index{一般位置约化|hyperpage}
令 $\Bun_G^{B\text{-gen}}$ 和 $\Bun_G^{N\text{-gen}}$ 分别表示配备一般位置定义的 $B$ 约化和 $N$ 约化的 $G$ 丛预叠。有交换图
\origdiagram{p023-square.pdf}{0.3142}
这里，$\mathrm{str}$ 记住底层 $G$ 丛，$L$ 忘去一般位置的 $N$ 结构。一个 $B$ 丛的一般位置 $N$ 约化，等价于它诱导的 $T$ 丛的一般位置平凡化，因此
\[\Bun_B^{N\text{-gen}}=\Bun_B\times_{\Bun_T}\Bun_T^{\mathrm{rat\ triv}}\simeq\Bun_B\times_{\Bun_T}\operatorname{Div}_\Lambda,\]
其中 $\operatorname{Div}_\Lambda$ 是余权幺半群的带色除子空间。这是在几何上替代允许有理约化并记录其缺陷除子的做法。

\begin{transnote}{含糊点拨}
$\Bun_T^{\mathrm{rat\ triv}}$ 指带有有理平凡化的 $T$-丛：除 $T$-丛 $\mathcal L$ 外还带一个在一般点处取得的同构 $\mathcal L_\eta\simeq\mathcal L^0_\eta$。这样的对象在同构意义下由「在哪些点、以什么余权偏离平凡」完全决定，于是与 $\Lambda$-着色的除子空间 $\operatorname{Div}_\Lambda$（有限个点各标一个 $\lambda\in\Lambda$）等同；$\Lambda$ 这里取余权幺半群。所谓缺陷除子就是这组数据：它记录一般位置约化在每一点「不平凡到什么程度」，而 $\Bun_T^{\mathrm{rat\ triv}}$ 与 $\operatorname{Div}_\Lambda$ 的等同说明这部分信息可以完全交给除子承担——这也是下面 Whittaker 层能按 $\operatorname{Div}_\Lambda$ 逐点定义的原因。
[依据：原书 p.23–24 对 $\Bun_B^{N\text{-gen}}$ 与 $\operatorname{Div}_\Lambda$ 的描述；并见 [Gai18] 对一般位置约化与 Whittaker 范畴的构造。]
\end{transnote}

对 $G=\GL_n$，这里以 $\Bun_N$ 表示具有固定 $T$ 丛的 $\rho$ 扭曲 Whittaker 叠：
\[\Bun_N=\{0=\overline F_0\subset\overline F_1\subset\cdots\subset\overline F_n,\quad\overline F_i/\overline F_{i-1}\simeq\Omega_X^{n-i}\}.\]
简单根扩张类给出映射
\[\ell:\Bun_N\longrightarrow\prod_{i=1}^{n-1}\Ga\xrightarrow{\Sigma}\Ga.\]
事实上，对 $1\le i\le n-1$，第 $i$ 个相邻扩张
\[0\longrightarrow\Omega_X^{n-i}\longrightarrow\overline F_{i+1}/\overline F_{i-1}\longrightarrow\Omega_X^{n-i-1}\longrightarrow0\]
定义 $\Ext^1(\Omega_X^{n-i-1},\Omega_X^{n-i})\simeq H^1(X,\Omega_X)\simeq k$ 中的一个元素，而 $\ell$ 是这些类的和。沿 $\ell$ 拉回 Artin–Schreier 层 $\cL_\psi$，得到整体 Whittaker 层。以 $D(\Bun_G^{N\text{-gen}})^{\mathrm{Wh}}$ 表示带有这种扭曲 $N(\A)$ 等变性的范畴。

因此，对 $G=\GL_n$，mirabolic Whittaker 展开的几何形式为
\[\Av_!:D(\Bun_G^{N\text{-gen}})^{\mathrm{Wh}}\xrightarrow{\sim}D(\Bun_G^{Q_1\text{-gen}})_{\mathrm{cusp}}.\]
\srcpage{24}
它应当与局部 Whittaker–Satake 等价以及谱一侧一起理解：
\origdiagram{p024-diagram.pdf}{0.6489}
这里 $\delta_E$ 是局部系统 $E$ 处的摩天层，左边的手征积分是局部 Whittaker 范畴的整体化。
\section{退化特征标与抛物子群}\label{sec:3.8}\index{Whittaker 特征标!退化|hyperpage}
$N(\A)$ 上在 $N(K)$ 上平凡且经由简单根商分解的特征标，由简单根坐标控制：
\[N(K)^\perp:=\{\chi:N(\A)\to\Qell^\times\mid\chi|_{N(K)}=1\}_{\mathrm{sr}}\simeq K^{\oplus(n-1)}.\]
这里“sr”表示简单根商。环面 $T(K)$ 通过简单根作用，所有坐标非零的开子集等同于 $T_{\mathrm{ad}}(K)$。因此，$T(K)$ 轨道由哪些简单根坐标非零来标记。等价地，
\[\{N(K)^\perp\text{ 中的 }T(K)\text{ 轨道}\}\longleftrightarrow\{J\subset\operatorname{Dyn}(G)\}\longleftrightarrow\{P\mid B\subset P\subset G\},\]
故有 $2^{n-1}$ 个轨道。对 Levi 子群为 $M$ 的抛物子群 $P$，令 $J(P)$ 为 $M$ 的简单根集合，以 $\Psi_P:N(\A)\to\Qell^\times$ 表示一个代表特征标，其简单根坐标恰好在 $J(P)$ 上非零。一般位置特征标为 $\Psi=\Psi_G$，零特征标则是与 $P=B$ 相联系的轨道。

扩展的 Piatetski-Shapiro–Shalika 定理使用所有这些 Fourier 系数；这是文献\cite{FGKV98}中 Whittaker 比较背后的函数论输入。现在令
\[Q=\begin{pmatrix}\GL_{n-1}&V\\0&\Gm\end{pmatrix}\]
为型 $(n-1,1)$ 的极大抛物子群。存在 Fourier 系数映射
\[C^\infty(Q(K)\backslash G(\A))\longrightarrow\prod_{P\supset B}C^\infty((Z_M(K)\ltimes N(\A),\Psi_P)\backslash G(\A)),\]
其中 $Z_M$ 是 $P$ 的 Levi 子群的中心，并稳定 $\Psi_P$；该记号表示 $Z_M(K)$ 对特征标的作用平凡。限制到尖点部分后，所有退化分量都为零，并与 $P=G$ 的因子等同，即非退化 Whittaker 系数。从 $Q$ 转到 $Q_1$ 后，这就恢复了上述 mirabolic 定理。
\section{沿 mirabolic 子群迭代 Fourier 变换}\label{sec:3.9}\index{Fourier 变换!迭代|hyperpage}
在迭代证明中使用 mirabolic 子群 $Q_1=L_1\ltimes V$，其中 $L_1\simeq\GL_{n-1}$，而不使用完整的极大抛物子群。$V(K)\backslash V(\A)$ 上的 Fourier 变换使用
\[V(K)^\perp=\operatorname{Ch}(V(\A)/V(K)).\]

\begin{transnote}{易错预警}
$\operatorname{Ch}(-)$ 是特征标群（Pontryagin 对偶）的记号，全书仅此处出现：$V(\A)$ 是自对偶的，$V(K)\backslash V(\A)$ 紧且完全不连通，其对偶由 $V(\A)$ 上在 $V(K)$ 上平凡的连续特征标组成，仍以 $V(\A)/V(K)$ 为参数。因此 $V(K)^\perp$ 不是「正交补空间」，而是消没子，等于 $\operatorname{Ch}(V(\A)/V(K))$。正因为要数的是对偶这一边的轨道，才轮到 $L_1(K)$ 在 $V(K)^\perp$ 上作用——在 $V(\A)/V(K)$ 上看不出这里的分裂。
[来源：nLab "Pontrjagin dual"，https://ncatlab.org/nlab/show/Pontryagin+duality（2026-11-17 访问）；并见原书 p.24 对 $V(K)^\perp$ 的使用。]
\end{transnote}
$L_1(K)$ 在 $V(K)^\perp$ 上的作用有两个轨道：
\srcpage{25}
\begin{enumerate}
\item 经过非零特征标 $\eta_{n-1}$ 的开轨道，其稳定子是较小的 mirabolic 子群 $Q'_1(K)\subset\GL_{n-1}(K)$；
\item 闭轨道 $\{0\}$，其稳定子为 $L_1(K)$。
\end{enumerate}
因此，Fourier 变换给出分解
\[
\begin{aligned}
C^\infty(Q_1(K)\backslash G(\A))\simeq{}&C^\infty((Q'_1(K)\ltimes V(\A),\eta_{n-1})\backslash G(\A))\\
&\times C^\infty(L_1(K)\ltimes V(\A)\backslash G(\A)).
\end{aligned}
\]
若 $n\ge3$，则 $Q'_1=L'_1\ltimes V'$ 是 $\GL_{n-1}$ 中的 mirabolic 子群。第一个因子又是秩 $n-1$ 的 Whittaker 型空间，因为 $Q'_1(\A)$ 作用于 $C^\infty((V(\A),\eta_{n-1})\backslash G(\A))$。第二个因子包含在 $C^\infty(Q'_1(K)\ltimes V(\A)\backslash G(\A))$ 中，并被尖点性消去。迭代这种降低秩的 Fourier 变换后，只剩下一般位置 Whittaker 系数。

\begin{transnote}{含糊点拨}
$Q'_1\subset\GL_{n-1}$ 是 $\GL_{n-1}$ 的 mirabolic 子群，即稳定某个非零线性泛函的子群，由这里的 $\eta_{n-1}$ 给出；它此前未出现，正是下面所说的 $Q'_1=L'_1\ltimes V'$。被跳过的一步是 $\GL_{n-1}(K)$ 在 $V(K)^\perp$ 上的轨道分类：非零元素构成一条开轨道，稳定子就是稳定该泛函的 mirabolic 子群；$0$ 单独是闭轨道，稳定子为整个 $\GL_{n-1}(K)=L_1(K)$。Fourier 变换按这两条轨道把函数分成两块——开轨道给出秩降一的 Whittaker 型空间，闭轨道给出在 $L_1(K)$ 上几乎不变的部分，被尖点性消去。迭代这一步就是逐次降秩的过程。
[依据：原书 p.25 对两条轨道与 $Q'_1$ 的说明；并见 [FGKV98]（JAMS 11 (1998), 451–484）中 mirabolic 子群与 Whittaker 展开的用法。]
\end{transnote}

对于极大抛物子群 $Q$，其范畴类比表现为到粘合的扩展 Whittaker 范畴的全忠实函子：
\[D(\Bun_G^{Q\text{-gen}})\hookrightarrow\sum_{P\supset B}D(\Bun_G^{Z_M\text{-gen},\mathrm{Wh}_P}).\]
由 $P$ 标记的部分，将具有特征标 $\Psi_P$ 及 $Z_M(K)$ 约化的退化系数几何化。对于 $G=\GL_n$ 以及型为 $(n-1,1)$ 的 $P$，这是经典迭代 Fourier 变换证明的几何体现。
\begingroup\setlength{\belowdisplayshortskip}{2pt}
\section{使用凝聚层的紧化变体}\label{sec:3.10}\index{凝聚层!紧化|hyperpage}
采用使 $Q_1$ 稳定所选向量的右作用约定，$G=\GL_n$ 的 $Q_1$-一般位置约化可以写成一个非零有理向量：
\[\Bun_G^{Q_1\text{-gen}}=\{(M,\Omega_X^{n-1}\dashrightarrow M),\quad\text{不恒为 }0\}.\]
其正则版本为
\[\Bun'_n=\{(M,s:\Omega_X^{n-1}\to M)\mid s\text{ 在一般点非零},\ M\in\Bun_n\},\]
并有遗忘映射 $\Bun'_n\to\Bun_n$。由于 $M$ 局部自由且 $X$ 为光滑曲线，在一般点非零的映射 $s$ 作为凝聚层的映射是单射，并给出
\[0\longrightarrow\Omega_X^{n-1}\longrightarrow M\longrightarrow M'\longrightarrow0.\]
这里 $M'$ 是秩 $n-1$ 的凝聚层；由于 $s$ 的像未必饱和，它未必局部自由。扩张类属于 $\Ext^1(M',\Omega_X^{n-1})$，Serre 对偶给出
\[\Ext^1(M',\Omega_X^{n-1})^*\simeq\Hom(\Omega_X^{n-2},M').\]
这恰好是下一次 Fourier 变换所在的对偶向量空间。

引入紧化叠
\[\Coh'_m=\{(\cF,\Omega_X^{m-1}\xrightarrow{s}\cF)\mid\cF\in\Coh,\ \rank\cF=m\},\qquad{}^{\ne0}\Coh'_m\subset\Coh'_m,\]
其中开子叠由 $s$ 在一般点非零定义。更精确地，对 $m=n,n-1,\ldots,2$，令 $\mathcal E_m\to\Coh_{m-1}$ 为相对向量叠，它在 $M'$ 上的纤维为 $\Ext^1(M',\Omega_X^{m-1})$。Serre 对偶将对偶相对向量叠等同于 $\mathcal E_m^\vee\simeq\Coh'_{m-1}$，其在 $M'$ 上的纤维为 $\Hom(\Omega_X^{m-2},M')$。映射 $\mathcal E_m\to\Coh'_m$ 记住带有典范截面的中间项；置 ${}^{\ne0}\mathcal E_m:=\mathcal E_m\times_{\Coh'_m}{}^{\ne0}\Coh'_m$。相继的 Fourier 变换可概括为
\[{}^{\ne0}\Coh'_m\xleftarrow{p_m}{}^{\ne0}\mathcal E_m\xrightarrow{\mathsf F}\mathcal E_m^\vee\simeq\Coh'_{m-1},\qquad m=n,n-1,\ldots,2.\]
\srcpage{26}
\enlargethispage{\baselineskip}
第一次变换按照文献\cite{KL85}的归一化，定义 Fourier–Deligne 函子
\[\mathsf F_n:D({}^{\ne0}\Coh'_n)\longrightarrow D(\Coh'_{n-1}).\]
重复这个构造，每次将秩降低一步。因此，从 ${}^{\ne0}\Coh'_n$ 上的 $E$-Hecke 特征层出发，便得到 ${}^{\ne0}\Coh'_1$ 上的层。秩一的终点必定是 Whittaker 层；这是 Laumon 构造中所用 Whittaker 归一化的紧化形式。
\endgroup

\chapter{Laumon 层}\label{ch:4}\index{Laumon 层|hyperpage}
% 原PDF第26页中部
设 $X$ 是 $k$ 上光滑射影连通曲线，$E$ 是 $X$ 上秩 $n$ 的局部系统。局部描述是在 étale 局部意义下的，因此下文 $\A^1$ 上的例子只是局部模型。目标是在 $X$ 上挠层的叠 $\Tor$ 上构造反常层 $\cL_E$。这是 Laumon 构造 $\GL_n$ 自守层时用来替代 Whittaker 函数的层论输入\cite{FGKV98}；这一构造的抛物变体与分歧变体见文献\cite{Hei04}。
\section{来自 Whittaker 与 Satake 的动机}\label{sec:4.1}
Whittaker 讲授（见\chref{3}）中理想的图景为
\origdiagram{p026-diagram.pdf}{0.5829}
这里 $\delta_E$ 是局部系统 $E$ 处的摩天层。应当把希望得到的 Whittaker 对象理解为
\[W_E\simeq\delta_E*W_0,\]
其中 $W_0$ 是真空 Whittaker 层。将 $\checkG$ 的表示 $V$ 作用于整个 $\checkG$-局部系统 $E$，得到 $X$ 上的局部系统 $V_E$，其在 $x$ 处的纤维为 $V_{E,x}$。对 $G=\GL_n$，这是对秩 $n$ 局部系统 $E$ 的通常操作。Laumon 层是在 $G=\GL_n$ 的正 Hecke 轨迹上对这一原则的具体实现。

\begin{transnote}{含糊点拨}
卷积 $*$ 是 Whittaker 范畴 $\operatorname{Wh}(\dot{\Gr}_{G,x})$ 上沿 Hecke 对应所作的 $!$-卷积：两个对象卷起来，就是先用一次修改把其中一个搬运过去、再叠上另一个，等价于把「作用」这一操作并入范畴。$W_0$ 是真空 Whittaker 层，即对应平凡 $\checkG$-表示（真空向量）的那个对象；$\delta_E$ 是 $E$ 处的摩天层，作用是把 $\checkG$ 的表示 $V$ 的实现「钉」在局部系统 $E$ 上。于是 $\delta_E*W_0$ 恰好是「把 $V_E$ 作用于真空」所得的 Whittaker 对象，也就是几何 Satake 的 $\operatorname{Sat}_{V_E}$ 在 Whittaker 一侧的对应物——这就是把 $W_E$ 写成这种卷积形式的意思。
[依据：原书 p.26–27 \S\ref{sec:4.1} 的图景；并见 [Gai18] 对 Whittaker 范畴上卷积的构造与 [FGKV98]（JAMS 11 (1998), 451–484）的 Whittaker 与 Satake 对应。]
\end{transnote}$\dot{\Gr}_{G,x}$ 中的上点是用于 Whittaker 层的同一个 $\rho$ 扭曲；下文的球面 Satake 输入写在普通仿射 Grassmann 簇上。

局部输入是几何 Satake：
\[\operatorname{Perv}_{G(\cO_x)}(\Gr_{G,x})\simeq\Rep(\checkG),\qquad\operatorname{IC}_{\overline{\Gr_{G,x}^\lambda}}\longleftrightarrow V^\lambda,\quad\lambda\in\Lambda^+.\]
采用 $\Gr_{G,x}\simeq G(K_x)/G(\cO_x)$ 及 $\Gr_{G,x}^\lambda=G(\cO_x)t_x^\lambda G(\cO_x)/G(\cO_x)$。

还存在几何 Casselman–Shalika 形式
\[\operatorname{Wh}(\dot{\Gr}_{G,x})\simeq\Rep(\checkG),\qquad W_\lambda\longleftrightarrow V^\lambda,\]
它与 Satake 作用相容。采用\chref{3}的 $\rho$ 扭曲记号，Whittaker 轨道为
\[S^\lambda:=\dot N(K_x)t_x^\lambda\dot G(\cO_x)/\dot G(\cO_x),\qquad\lambda\in\Lambda.\]
\srcpage{27}
$S^\lambda$ 上存在 $(\dot N(K_x),\Psi_x)$-等变层，当且仅当 $\lambda$ 优势；事实上，
\[\operatorname{Stab}_{\dot N(K_x)}(t_x^\lambda)\subset\ker\Psi_x\quad\Longleftrightarrow\quad\lambda\in\Lambda^+.\]
对这样的 $\lambda$，$S^\lambda$ 上相应的扭曲局部系统可洁净延拓到 $\dot{\Gr}_{G,x}$。一维模型是 Artin–Schreier 层：对 $j:\A^1\hookrightarrow\mathbb P^1$ 及 $i_\infty:\{\infty\}\hookrightarrow\mathbb P^1$，
\[i_\infty^*j_*\cL_\psi=0,\qquad\text{从而 }j_!\cL_\psi\simeq j_*\cL_\psi.\]
关键是 $j_*\cL_\psi$ 在边界点 $\infty$ 处的茎消失；对于零延拓，$i_\infty^!j_!$ 的消失则是恒然成立的。这是非退化 Whittaker 局部系统洁净延拓的一维模型。

\begin{transnote}{含糊点拨}
Artin–Schreier 层 $\cL_\psi$ 是 $\A^1$ 上由覆盖 $x\mapsto x^p-x$（$k$ 特征 $p$）所给的 $\mathbb Z/p$-局部系统在 $\psi$ 下的扭：它的单值性正是 $\psi$，因此在 $\infty$ 处有正则奇点。关键的一步是 $j_!$ 与 $j_*$ 只在边界处不同：对开浸入 $j$，锥 $\operatorname{Cone}(j_!\to j_*)$ 支撑在边界点 $\infty$ 上，故 $i_\infty^*j_*\cL_\psi=0$ 恰说明 $j_!\cL_\psi\to j_*\cL_\psi$ 在 $\infty$ 处的茎也为零，两者遂同构。这就是「洁净」的全部内容：边界茎消失，$!$ 延拓与 $*$ 延拓重合。零延拓那一侧的 $i_\infty^!j_!$ 恒为零，所以只需检查 $*$ 一侧。
[来源：nLab "perverse sheaf"，https://ncatlab.org/nlab/show/perverse+sheaf（Last revised 2026-10-16；2026-11-17 访问）；并见 [KL85] 与 [Gai18] 对洁净延拓的处理。]
\end{transnote}
\section{从 Satake 核到挠层}\label{sec:4.2}\index{挠层|hyperpage}
对 $V\in\Rep(\checkG)$，Satake 在局部球面仿射 Grassmann 簇上给出反常层 $\operatorname{Sat}(V)$。正则表示 $\cO(\checkG)$ 是 $\Rep(\checkG)$ 中的 ind-对象；形式上对它应用 Satake，得到一个 ind-球面 Hecke 核，而不是单个有限型反常层。对于 $G=\GL_n$，这个核是分类三元组 $(M,M',\beta:M\simeq M'\text{ 在一般点处})$ 的整体 Hecke 叠的局部模型。阿代尔双商 $\GL_n(\cO)\backslash\GL_n(\A)/\GL_n(\cO)$ 只记录选取一般位置平凡化之后的相对阿代尔位置，因此，缺少这一附加选择时，不能将它等同于完整的整体 Hecke 叠。

对 $G=\GL_n$，取出正轨迹
\[\Gr_{n,x}^+=\bigcup_{\lambda_1\ge\cdots\ge\lambda_n\ge0}\Gr_{n,x}^\lambda=(\operatorname{Mat}_{n\times n}(\cO_x)\cap\GL_n(K_x))/\GL_n(\cO_x).\]
它由在单位圆盘上定义且具有挠余核的修改组成。允许任意中心移位后，从这个正部分得到完整的仿射 Grassmann 簇。特定移位
\[(\cO_x^n\xrightarrow{\alpha}\mathcal E)\longmapsto(\cO_x(-x)^n\hookrightarrow\cO_x^n\xrightarrow{\alpha}\mathcal E)\]
使余权增加 $(1,\ldots,1)$。在整体正 Hecke 叠上，有映射
\[p:\Hecke_{\GL_n}^+\longrightarrow\Tor,\qquad(M\hookrightarrow M')\longmapsto M'/M,\]
其中 $\Hecke_{\GL_n}^+$ 分类具有挠商的秩 $n$ 丛嵌入。因此，泛正 Hecke 核应当是 $\Tor$ 上一个层的拉回。这个层就是 $\cL_E$。
\section{挠层叠}\label{sec:4.3}
定义
\[\Tor(S):=\{M\in\Coh(X\times S)\mid M\text{ 在 }S\text{ 上平坦，每个几何纤维均为挠层}\}.\]
它按长度分解：
\[\Tor=\coprod_{d\ge0}\Tor^d.\]
每个 $\Tor^d$ 都是代数叠。当 $X$ 真且 $H^0(X,\cO_X)=k$ 时，可以从 Quot 概形的开部分取得一个方便的图册：
\[\operatorname{Quot}_{\cO_X^{\oplus d}}^{\circ,d}=\{\cO_X^{\oplus d}\twoheadrightarrow M\mid M\in\Tor^d,\ H^0(\cO_X^{\oplus d})\xrightarrow{\sim}H^0(M)\}.\]
改变 $H^0(M)$ 中所选基，得到 $\GL_d$ 作用，相应商叠为
\[[\operatorname{Quot}_{\cO_X^{\oplus d}}^{\circ,d}/\GL_d]\simeq\Tor^d.\]
\srcpage{28}
\begin{example}
若 $X=\A^1$，长度 $d$ 的挠层等同于一个带有一个自同态的 $d$ 维向量空间，因此
\[\Tor^d\simeq\mathfrak{gl}_d/\GL_d.\]
\end{example}
\section{按分拆的整体分层}\label{sec:4.4}\index{分拆|hyperpage}\index{挠层!整体分层|hyperpage}
令
\[\operatorname{Part}^d=\{(d_1\ge d_2\ge\cdots\ge d_s>0)\mid d_1+\cdots+d_s=d\}.\]
叠 $\Tor^d$ 按局部闭子叠 $\Tor^{(d_1,\ldots,d_s)}$ 分层。置 $d_{s+1}=0$。由 $(d_1,\ldots,d_s)$ 标记的层是映射
\[X^{(d_1-d_2)}\times X^{(d_2-d_3)}\times\cdots\times X^{(d_s)}\longrightarrow\Tor^d\]
的像，其中
\[(D_1,\ldots,D_s)\longmapsto\cO_{D_1+\cdots+D_s}\oplus\cO_{D_2+\cdots+D_s}\oplus\cdots\oplus\cO_{D_s}.\]
分层 $\Tor^{(d)}$ 是开的；它是
\[X^{(d)}\longrightarrow\Tor^d,\qquad D\longmapsto\cO_D\]
的像。对 $X=\A^1$，这意味着
\[\Tor^{(d)}\simeq\mathfrak{gl}_d^{\mathrm{reg}}/\GL_d,\qquad\Tor^{(d)}|_{X^{(d)}-\Delta}\simeq\mathfrak{gl}_d^{\mathrm{rss}}/\GL_d,\qquad X^{(d)}\simeq\mathfrak t_d/W.\]
因此，挠层叠是在曲线情形下对 $\mathfrak{gl}_d$ 的伴随商图景的类比。
\section{\texorpdfstring{$\Tor^d$}{Torᵈ} 的 Springer 解消}\label{sec:4.5}
令 $\widetilde{\Tor}^d$ 为长度 $d$ 挠层中完整旗的叠：
\[\widetilde{\Tor}^d=\{0=M_0\subset M_1\subset\cdots\subset M_d=M\mid\length(M_i/M_{i-1})=1\}.\]
有自然映射 $\pi:\widetilde{\Tor}^d\to\Tor^d$ 及 $o_d:\widetilde{\Tor}^d\to X^d$，其中 $\pi$ 忘去旗，$o_d$ 记录商 $M_i/M_{i-1}$ 的有序支撑。记 $s_d:\Tor^d\to X^{(d)}$ 和 $r_d:X^d\to X^{(d)}$ 分别为支撑映射与对称化映射。若 $q:\Tor^{(d)}\to X^{(d)}$ 为循环层上的支撑映射，则图为
\origdiagram{p028-diagram.pdf}{0.2413}
在具有无重数支撑的正则半单轨迹 $\Tor^{d,\mathrm{rss}}$ 上，$s_d$ 将 $\Tor^{d,\mathrm{rss}}$ 展示为 $X^{(d),\mathrm{rss}}$ 上的标量自同构胚叠；$\widetilde{\Tor}^{d,\mathrm{rss}}$ 则是 $X^{d,\mathrm{rss}}$ 上的有序标量自同构胚叠。将这些标量自同构刚性化后，$\pi$ 成为通常的有序 $S_d$ 覆盖 $X^{d,\mathrm{rss}}\to X^{(d),\mathrm{rss}}$。这与 Grothendieck–Springer 图平行\cite{GPR20,AHJR14}：

\begin{transnote}{背景直觉}
「标量自同构胚叠」记录的是每条支撑所带的 $\Gm$ 自同构：无重数支撑的点 $\{x_1,\ldots,x_d\}$ 上，挠层在每个 $x_i$ 处可乘一个标量而保持同构类，这组标量构成 $\Gm^d$ 的带；胚叠的 band 就是它，所以它不是 $S_d$-覆盖而比它多出一层 $\Gm$ 的规范群。刚性化（rigidification）的做法是沿这个 band 把对象再添一层「标记」以消掉多余自同构，压掉 $\Gm^d$ 之后剩下的恰好是 $S_d$-挠子，于是 $\pi$ 变成通常的有序覆盖 $X^{d,\mathrm{rss}}\to X^{(d),\mathrm{rss}}$。这套「先把 $\Gm$-自主性记进胚叠、再刚性化掉」的手法是 [GPR20] 处理 Grothendieck–Springer 图的标准步骤。
[来源：nLab "gerbe"，https://ncatlab.org/nlab/show/gerbe（Last revised 2026-02-14；2026-11-17 访问）；并见 [GPR20] 与 [AHJR14] 对 Grothendieck–Springer 图中胚叠与刚性化的用法。]
\end{transnote}
\[\mathfrak c_d\xrightarrow{s}\mathfrak{gl}_d^{\mathrm{reg}}/\GL_d\longrightarrow\mathfrak c_d,\]
\[\widetilde{\mathfrak{gl}}_d/\GL_d\longrightarrow\mathfrak{gl}_d/\GL_d\longrightarrow\mathfrak c_d,\qquad\mathfrak t_d\longrightarrow\mathfrak c_d=\mathfrak t_d/W,\]
其中 $s$ 是 Kostant 截面。局部于 $X$，这一类比是字面上的：一点处的挠层就是幂零自同态，挠层的子层旗就是 Springer 旗。
\section{Springer–Laumon 层与 Laumon 直和项}\label{sec:4.6}
\srcpage{29}
定义次数 $d$ 的 Springer–Laumon 层为
\[\operatorname{Spr}_E^d:=\pi_!o_d^*(E^{\boxtimes d})[d]\in D_c^b(\Tor^d).\]
移位 $[d]$ 是这里等变反常约定所采用的归一化。$S_d$ 作用是 Springer 单值作用\cite{AHJR14}：在无重数支撑轨迹上，刚性化后 $\pi$ 是有序支撑的 $S_d$ 覆盖，所得作用通过中间延拓扩展到 $\operatorname{Spr}_E^d$；等价地，也可通过通常的简单反射对应得到。它不是通过任意置换一条嵌套旗而在整个旗叠上给出的字面意义上的 $S_d$ 作用。
\begin{theorem}{Laumon 延拓定理\cite{KL85}}{4.6.1}
令 $j_{\mathrm{rss}}:\Tor^{d,\mathrm{rss}}\hookrightarrow\Tor^d$ 为正则半单轨迹，$s_{\mathrm{rss}}:\Tor^{d,\mathrm{rss}}\to X^{(d),\mathrm{rss}}$ 为支撑映射，$r_{\mathrm{rss}}:X^{d,\mathrm{rss}}\to X^{(d),\mathrm{rss}}$ 为有序覆盖。层 $\operatorname{Spr}_E^d$ 是反常层，且是从正则半单轨迹作出的中间延拓：
\[\operatorname{Spr}_E^d\simeq j_{\mathrm{rss},!*}(s_{\mathrm{rss}}^*r_{\mathrm{rss},*}(E^{\boxtimes d})[d]).\]
较大的循环正则轨迹 $\Tor^{(d)}$ 也确定同一个对象；它在该轨迹上的限制，是在 $X^{(d)}$ 内跨越对角线作中间延拓所得的对象。
\end{theorem}

Laumon 层是这一 $S_d$ 作用的平凡等型直和项：
\[\cL_E^d:=\Hom_{S_d}(\mathrm{triv},\operatorname{Spr}_E^d)=(\operatorname{Spr}_E^d)^{S_d}\in\operatorname{Perv}(\Tor^d).\]
等价地，若在 $X^{(d),\mathrm{rss}}$ 上令 $E^{(d)}:=(r_{\mathrm{rss},*}E^{\boxtimes d})^{S_d}$，则
\[\cL_E^d\simeq j_{\mathrm{rss},!*}(s_{\mathrm{rss}}^*E^{(d)}[d]).\]
这个局部系统在无重数轨迹上秩为 $n^d$，不应与 $\Sym^dE$ 混淆。
\begin{lemma}{}{4.6.2}
若 $E$ 几何不可约，则 $\cL_E^d\in\operatorname{Perv}(\Tor^d)$ 几何不可约。
\end{lemma}
\begin{proof}\textbf{概要。}
基变换到 $\bar k$ 后，Schur 引理给出 $\End_{\pi_1(X_{\bar k})}(E_{\bar x})=e$，其中 $e$ 是系数域。无重数配置轨迹上的局部系统 $E^{(d)}$，对应于不可约外张量积 $E_{\bar x}^{\otimes d}$ 作为花环积群 $\pi_1(X_{\bar k})^d\rtimes S_d$ 的表示。它的自同态代数为
\[\End_{\pi_1(X_{\bar k})^d\rtimes S_d}(E_{\bar x}^{\otimes d})=(e^{\otimes d})^{S_d}=e.\]
因此，它不可约。拉回到标量自同构胚叠保持不可约性，中间延拓保持底层反常层的不可约性。
\end{proof}

最后，对 $\cL_E\in\operatorname{Perv}(\Tor)$ 定义
\[\cL_E^d:=\cL_E|_{\Tor^d}.\]
以下小节局部分析这个对象：限制到支撑于单一点的挠层后，普通 Springer 理论将局部反常对象 $\cL_{E,x}^d[-d]$ 分解为幂零层的 IC 层，其系数是 $E_x$ 的最高权表示。
\section{归约到单点处的局部挠层}\label{sec:4.7}
\srcpage{30}
为进行局部计算，以 $x$ 表示 $X$ 的一个几何点。令 $i_x:\Tor_x^d\hookrightarrow\Tor^d$ 为 $s_d$ 在 $d\cdot x$ 上的纤维。因此，$\Tor_x^d$ 分类支撑于 $x$ 的长度 $d$ 挠层。置
\[\cL_{E,x}^d:=i_x^*\cL_E^d.\]
由于这一纤维在对称幂方向上余维为 $d$，局部反常对象是 $\cL_{E,x}^d[-d]$。
\begin{lemma}{不交支撑的因子化}{4.7.1}
若 $x\ne y$，直和给出局部挠层叠的同构
\[\Tor_x^a\times\Tor_y^b\xrightarrow{\sim}\Tor_{ax+by}^{a+b},\qquad\cL_{E,x}^a\boxtimes\cL_{E,y}^b\xrightarrow{\sim}\cL_{E,ax+by}^{a+b}.\]
更一般地，对具有不同支撑点的除子 $D=\sum_xd_xx$，
\[\prod_x\Tor_x^{d_x}\xrightarrow{\sim}\Tor_D^d,\qquad\boxtimes_x\cL_{E,x}^{d_x}\simeq\cL_{E,D}^d.\]
\end{lemma}

因此，$\cL_E^d$ 的结构归约为计算 $\cL_{E,x}^d[-d]\in\operatorname{Perv}(\Tor_x^d)$。
\section{按分拆的局部分层}\label{sec:4.8}\index{挠层!局部分层|hyperpage}
设 $\lambda=(\lambda_1\ge\lambda_2\ge\cdots\ge\lambda_s>0)$ 为 $d$ 的分拆，置 $\lambda_{s+1}=0$。记 $|\lambda|=d$，$\ell(\lambda)=s$。对有效除子 $D$，记 $\cO_D:=\cO_X/\cO_X(-D)$。整体分层 $\Tor^\lambda\subset\Tor^d$ 由图册
\[X^{(\lambda_1-\lambda_2)}\times X^{(\lambda_2-\lambda_3)}\times\cdots\times X^{(\lambda_s)}\longrightarrow\Tor^\lambda\]
描述，它将 $(D_1,\ldots,D_s)$ 映为
\[\cO_{D_1+\cdots+D_s}\oplus\cO_{D_2+\cdots+D_s}\oplus\cdots\oplus\cO_{D_s}.\]
分层 $\Tor^\lambda$ 局部闭地映入 $\Tor^d$。

$x$ 处的局部层是剩余部分 $\Tor_x^\lambda\subset\Tor_x^d$，对应于挠层
\[T_{\lambda,x}:=\bigoplus_{i=1}^s\cO_X/\cO_X(-\lambda_ix).\]
等价地，过渡到几何点 $x$ 处的完备局部环并选取素元 $t$ 后，它是 $k[[t]]$ 模
\[\bigoplus_{i=1}^sk[[t]]/(t^{\lambda_i}).\]
叠 $\Tor_x^d$ 的维数为 $-d$；在下述幂零锥模型中，它的各分层是幂零轨道的商叠。

\begin{transnote}{含糊点拨}
维数可以算出来：幂零锥 $\cN_d\subset\mathfrak{gl}_d$ 由 $\det=0$（一个方程）在 $d^2$ 维的 $\mathfrak{gl}_d$ 中切出，故 $\dim\cN_d=d^2-d$；商去 $\GL_d$ 的 $d^2$ 维轨道，得 $\dim\Tor_x^d=(d^2-d)-d^2=-d$。负维数不是笔误：商叠的维数按「维数减群维数」计，正是「$\GL_d$ 的作用比底空间还大」的表现。各分层 $\Tor_x^\lambda$ 则是幂零轨道 $\cO_\lambda$ 的商叠，维数同样要减去 $d^2$——这与 \S\ref{sec:4.8} 末尾的幂零锥模型 $\Tor_x^d\simeq\cN_d/\GL_d$ 一致。
[来源：nLab "algebraic stack"，https://ncatlab.org/nlab/show/algebraic+stack（Last revised 2023-04-16；2026-11-17 访问）；并见 [GPR20] 对 $\Tor^d$ 的幂零锥描述（原书 p.30）。]
\end{transnote}
\begin{definition}{}{4.8.1}
对分拆 $\lambda\vdash d$，定义
\[B_{\lambda,x}:=\operatorname{IC}_{\overline{\Tor_x^\lambda}}.\]
\end{definition}
\section{Laumon 层的局部分解}\label{sec:4.9}
令 $V=E_x$ 为 $E$ 在 $x$ 处的纤维。对满足 $s\le n$ 的分拆 $\lambda=(\lambda_1\ge\cdots\ge\lambda_s>0)$，添零将其扩展为 $\GL(V)\simeq\GL_n$ 的优势权：
\[\lambda=(\lambda_1,\ldots,\lambda_s,0,\ldots,0).\]
令 $V^\lambda$ 为 $\GL(V)$ 中最高权为 $\lambda$ 的不可约最高权表示。

等价地，若 $\chi^\lambda\in\operatorname{Irr}(S_d)$ 是与 $\lambda$ 对应的 Specht 表示，则 Schur–Weyl 对偶给出
\[V^\lambda\simeq(V^{\otimes d}\otimes\chi^\lambda)^{S_d}.\]
若 $s>n$，这个向量空间为零。
\srcpage{31}
\begin{proposition}{局部公式}{4.9.1}
对每个几何点 $x\in X$ 与整数 $d\ge0$，置 $n(\lambda):=\sum_{i\ge1}(i-1)\lambda_i$。按通常的 Casselman–Shalika 归一化，将其看作 $n(\lambda)=\langle\lambda,(0,1,\ldots,n-1)\rangle$。则
\[\cL_{E,x}^d[-d]\simeq\bigoplus_{\substack{\lambda\vdash d\\\ell(\lambda)\le n}}B_{\lambda,x}\otimes E_x^\lambda(-n(\lambda)).\]
\end{proposition}
\section{与仿射 Grassmann 簇的关系}\label{sec:4.10}
令 $\Gr_{n,x}$ 为 $\GL_n$ 在 $x$ 处的仿射 Grassmann 簇。令 $\Gr_{n,x}^{d,\mathrm{mod}}$ 为次数 $d$ 的正修改轨迹，即其点是 $\cO_x^n$ 的长度 $d$ 商的轨迹；这一记号不应与正则幂零轨迹混淆。令
\[f:\Gr_{n,x}^{d,\mathrm{mod}}\longrightarrow\Tor_x^d\]
为将修改映为其商挠层的映射。它的像是最多可由 $n$ 个元素生成的挠层所构成的开子叠；在这一像上，$f$ 光滑，相对维数为 $nd$。对 $|\lambda|=d$ 且 $\ell(\lambda)\le n$，令 $\Gr_{n,x}^\lambda$ 为型 $\lambda$ 的 Schubert 分层。图为
\origdiagram{p031-square.pdf}{0.2596}
置
\[A_\lambda:=\operatorname{IC}_{\overline{\Gr_{n,x}^\lambda}}.\]
采用 IC 包括反常上同调移位但不包括 Tate 扭转的约定，光滑拉回给出
\[f^*B_{\lambda,x}\simeq A_\lambda[-dn].\]
等价地，$f^*B_{\lambda,x}[dn]\simeq A_\lambda$。局部 Laumon 公式中的 Tate 因子 $-n(\lambda)$ 属于 Whittaker／Casselman–Shalika 归一化，而不属于 IC 层的光滑拉回比较。
\newpage
\section{用 Springer 理论证明局部公式}\label{sec:4.11}
\begin{proof}
现在使用经典 Springer 对应\cite{GPR20,AHJR14}说明局部分解的证明。在 $x$ 的 étale 局部工作，或将 $X$ 替换为 $x$ 处的形式圆盘。选取素元后，支撑于 $x$ 的长度 $d$ 挠层就是一个带有幂零自同态的 $d$ 维向量空间。因此
\[\Tor_x^d\simeq\cN_d/\GL_d,\]
其中 $\cN_d\subset\mathfrak{gl}_d$ 是幂零锥。非局部陈述 $\Tor^d\simeq\mathfrak{gl}_d/\GL_d$ 对 $X=\A^1$ 是字面成立的；对一般曲线，它是在次数 $d$ 除子附近相应的 étale 局部模型。在这一局部等同下，映射 $\widetilde{\Tor}^d\to\Tor^d$ 成为 Grothendieck–Springer 解消，限制到 $\Tor_x^d$ 上则成为幂零锥的 Springer 解消。

基变换到几何点 $\bar x$，并暂时省去 Tate 扭转。局部反常 Springer 层分解为
\[\operatorname{Spr}_{\bar x}^d[-d]\simeq\bigoplus_{\chi\in\operatorname{Irr}(S_d)}\chi\otimes\operatorname{IC}_{\overline{\cO_\chi}}.\]
\srcpage{32}
这里，$\operatorname{Irr}(S_d)$ 等同于 $d$ 的分拆，也等同于 $\mathfrak{gl}_d$ 中的幂零轨道。采用以下约定：
\[\lambda\longleftrightarrow\cO_\lambda,\qquad\text{平凡表示}\longleftrightarrow\text{正则幂零轨道 }(d).\]
这确定了不取转置的约定。

对 Laumon 层，纤维 $E_x^{\otimes d}$ 随 Springer 层一起带入，$S_d$ 通过置换 $d$ 个因子作用于这个张量因子。因此
\[\operatorname{Spr}_{E,x}^d[-d]\simeq\bigoplus_{\chi\in\operatorname{Irr}(S_d)}\chi\otimes E_x^{\otimes d}\otimes\operatorname{IC}_{\overline{\cO_\chi}}.\]
取 $S_d$ 不变直和项得到
\[(\operatorname{Spr}_{E,x}^d[-d])^{S_d}\simeq\bigoplus_{\chi\in\operatorname{Irr}(S_d)}(\chi\otimes E_x^{\otimes d})^{S_d}\otimes\operatorname{IC}_{\overline{\cO_\chi}}.\]
现在写 $\chi=\chi^\lambda$，其中 $d=\lambda_1+\cdots+\lambda_s$、$\lambda_1\ge\cdots\ge\lambda_s>0$。添零将 $\lambda$ 扩展到长度 $n$。由 Schur–Weyl 对偶，
\[(\chi^\lambda\otimes E_x^{\otimes d})^{S_d}\simeq E_x^\lambda,\]
若 $s>\dim(E_x)=n$ 则它消失。最后，在 $\Tor_x^d\simeq\cN_d/\GL_d$ 下，有
\[\operatorname{IC}_{\overline{\cO_\lambda}}=B_{\lambda,x}.\]
恢复 Tate 扭转，便得到命题\ref{pro:4.9.1}关于 $\cL_{E,x}^d[-d]$ 的公式。
\end{proof}

\chapter{Fourier 变换}\label{ch:5}\index{Fourier–Deligne 变换|hyperpage}
% 原PDF第32页下半部
\section{回顾：从 Laumon 层到初步的自守层}\label{sec:5.1}
继续取 $G=\GL_n$，并令 $E$ 为 $X$ 上秩 $n$ 的局部系统。\chref{4}构造了 $\cL_E\in\operatorname{Perv}(\Tor)$。这里的目标是用它在 $\Bun_n$ 上构造自守层 $\Aut_E$。

令 $\Bun'_n$ 为对 $(M,s)$ 的叠，其中 $M$ 是秩 $n$ 向量丛，$s:\Omega_X^{n-1}\hookrightarrow M$ 为单射。令 $Q_n$ 为正的整体 Hecke–Whittaker 叠
\[Q_n=\left\{\begin{array}{l}0=M_0\subset M_1\subset\cdots\subset M_n,\quad M_i/M_{i-1}\simeq\Omega_X^{n-i},\\M_n\hookrightarrow M\in C_n\text{，余核为挠层}\end{array}\right\}.\]
等价地，这是与 $\rho$ 扭曲 Whittaker 叠作纤维化的正 Hecke 对应的紧化整体版本；它不是局部仿射 Grassmann 簇与 $\Bun_N$ 的字面乘积。它有三个显然的映射：$\pi:Q_n\to\Bun'_n$ 记住 $\Omega_X^{n-1}=M_1\hookrightarrow M$；$p_1:Q_n\to\Tor$ 记录挠商 $M/M_n$；$p_2:Q_n\to\Bun_N\to\Ga$ 是 Whittaker 特征标，即旗中相邻扩张类的和。
\begin{transnote}{含糊点拨}
「$\rho$ 扭曲」指把 $\Bun_N$ 上的 Whittaker 层按半和正根 $\rho$ 作 Tate 扭转：沿 Hecke 修改搬运 $N$-丛时，次数按 $\langle\rho,\cdot\rangle$ 平移，不扭则 $p_2^*\cL_\psi$ 与 $\pi$、$p_1$ 的等变性不能同时成立。称其为「紧化整体版本」，是指旗 $0=M_0\subset\cdots\subset M_n$、挠商 $M/M_n$ 与 Whittaker 特征标捆在同一叠里：末项 $M_n\hookrightarrow M$ 由同一条链给出，不能拆成 $\Gr\times\Bun_N$ 的直积；被丢掉的正是有挠缺陷（$M/M_n\ne 0$）的修改，而那恰是 $j_{k,!}$ 要延拓过去的部分。[依据：原书 p.33 本段；并见 p.22 §3.1 的 $\rho$-移位。]
\end{transnote}

令 $\cL_\psi$ 为 $\Ga$ 上的 Artin–Schreier 层。在 $\length(M/M_n)=d$ 的分支上，初步对象为
\[\Aut_E^{\prime,d}:=\pi_!(p_1^*\cL_E^d\otimes p_2^*\cL_\psi)[\dim_{\mathrm{rel}}](\dim_{\mathrm{rel}}/2).\]
方括号及圆括号表示 Fourier 归一化所用的相对移位与 Tate 扭转。感叹号表示紧支撑推前；映射 $\pi$ 一般不是真态射。

所希望的性质如下。
\begin{enumerate}
\item 若 $E$ 不可约，则对 $d\gg0$，$\Aut_E^{\prime,d}$ 是不可约反常层。
\item 对 $d\gg0$，$\Aut_E^{\prime,d}$ 沿 $\Bun'_n\to\Bun_n^d$ 下降，得到 $\Aut_E^d\in\operatorname{Perv}(\Bun_n^d)$。
\srcpage{33}
\item $\Aut_E$ 满足截断的最高次 Hecke 关系。采用\chref{2}的 Hecke 对应约定，
\[t_x^*\Aut_E^d\simeq\det(E_x)\otimes\Aut_E^{d+n},\]
其中 $t_x:\Bun_n^{d+n}\xrightarrow{\sim}\Bun_n^d$，$M\mapsto M(-x)$。等价地，$(t_x)_!\Aut_E^{d+n}\simeq\det(E_x)^{-1}\otimes\Aut_E^d$。
\end{enumerate}
困难在于，从形式上没有任何理由保证紧支撑推前 $\pi_!$ 保持反常性或不可约性。本章的关键是：将构造改写为 Fourier–Deligne 变换\cite{KL85}后，这个问题就能够处理。
\section{向量丛的 Fourier 变换}\label{sec:5.2}
设 $V\to S$ 为秩 $r$ 向量丛，$V^\vee\to S$ 为其对偶，$\langle-,-\rangle:V\times_S V^\vee\to\A^1$ 为配对。Fourier–Deligne 变换\cite{KL85}为
\[\mathsf F_{V/S}(\cK):=p_!^\vee(p^*\cK\otimes\langle-,-\rangle^*\cL_\psi)[r](r/2),\qquad\cK\in D_c^b(V).\]
采用这一归一化，它对于反常 $t$-结构是 $t$-正合的，并保持单反常层的不可约性。因此，问题是将定义 $\Aut'_E$ 的层识别为一个洁净延拓的 Fourier 变换。
\begin{transnote}{含糊点拨}
归一化 $[r](r/2)$ 的作用是把对偶的两支拉平：$p^\vee_!$ 与 $p^\vee_*$ 经 Verdier 对偶只差一个移位和扭转，取 $[r](r/2)$ 后 $\mathsf F_{V/S}$ 对反常 $t$-结构 $t$-正合。$r/2$ 可以是半整数，故系数域须含 $\sqrt q$，$\Qell$ 满足这一点。「保持单反常层的不可约性」不是算出来的，而是 $t$-正合加光滑性的形式后果：$t$-正合函子把单对象送到单对象（可能带平移）。[依据：原书 p.33–34 本段；并见 [KL85] 第 2 章的 $\mathsf F$ 归一化与 t-正合性。]
\end{transnote}
\section{\texorpdfstring{$G=\GL_2$}{G=GL₂} 的情形}\label{sec:5.3}
对 $G=\GL_2$，记
\[\overline{\Bun}'_2=\{0\longrightarrow\Omega_X\longrightarrow M\longrightarrow L\longrightarrow0,\quad L\in\Coh_1\}\]
为凝聚扩张叠。用于自守层的叠 $\Bun'_2$，是其中中间项 $M$ 局部自由的开子叠。令 $\Coh_1$ 为秩一凝聚层的叠，并令
\[\Coh_1^+:=\{(L,\cO_X\xrightarrow{s}L)\},\qquad\Coh'_1\subset\Coh_1^+\]
为 $s$ 在一般点单射的开子叠。在 $\Coh'_1$ 上，$s$ 的余核是挠层；令 $q:\Coh'_1\to\Tor$ 为这个商映射。因此，下文所用的层是 $q^*\cL_E$，而不是 $\cL_E$ 到 $\Coh'_1$ 的字面限制。

存在笛卡尔方块
\origdiagram{p033-square.pdf}{0.5692}
其中，给定扩张 $0\to\Omega_X\to M\to L\to0$ 以及 $\cO_X\to L$ 后，丛 $M'$ 是拉回。在 $\Coh_1$ 中满足 $H^1(X,L)=\Ext^1(\cO_X,L)=0$ 的开子叠上，$\Coh_1^+\to\Coh_1$ 是纤维为 $\Hom(\cO_X,L)$ 的向量丛。纤维为 $\Ext^1(L,\Omega_X)$ 的扩张叠是对偶向量丛，因为 Serre 对偶给出
\[\Ext^1(L,\Omega_X)\simeq\Hom(\cO_X,L)^\vee.\]
中间项局部自由的开子叠，是 $\Bun'_2$ 的相应部分。因此，必须在 $\Coh_1$ 的这个开子叠上工作，并且必须使用 $\Coh_1^+$ 而不是 $\Coh'_1$，因为零截面是向量丛的一部分。

Whittaker 映射 $\square\to\Bun_N\to\Ga$ 恰好是配对
\[\Hom(\cO_X,L)\times\Ext^1(L,\Omega_X)\longrightarrow\Ext^1(\cO_X,\Omega_X)\simeq k.\]
\srcpage{34}
因此，在 $\Coh_1$ 的这个开子叠上，
\[\Aut'_E\simeq\mathsf F(j_!q^*\cL_E),\qquad j:\Coh'_1\hookrightarrow\Coh_1^+.\]
若 $j_!=j_{!*}=j_*$，则 $j_!q^*\cL_E$ 为单反常层，因而其 Fourier 变换是不可约反常层。这至少在 $\Bun'_2$ 上具有向量丛描述的开部分证明了所需陈述。
\section{Drinfeld 关于 \texorpdfstring{$\GL_2$}{GL₂} 的洁净延拓定理}\label{sec:5.4}\index{洁净延拓|hyperpage}
现在限制到线丛 $\Bun_1\subset\Coh_1$。由于线丛之间的映射要么为零，要么为单射，在 $\Bun_1$ 上有
\[\Bun'_1\xrightarrow{j}\Bun_1^+\xleftarrow{i}\Bun_1,\]
其中 $i$ 为零截面，$j$ 为它的开补集。
\begin{theorem}{Drinfeld，洁净延拓；见\cite{FGKV98,KL85}}{5.4.1}
假设 $E$ 几何不可约且秩为 $2$。若 $d>4(g-1)$，则层 $j_!q^*\cL_E^d$ 在 $\Bun_1^d\subset\Coh_1$ 上洁净。
\end{theorem}
\begin{proof}

在 $\Bun'_1$ 上令 $\cF=q^*\cL_E^d$。$j_!\cF$ 的洁净性，是指自然映射 $j_!\cF\to j_*\cF$ 为同构。由于锥支撑于零截面，这等价于
\[i^*j_*\cF=0.\]
由 Verdier 对偶，只需对 $\cF$ 及其对偶证明紧支撑的边界消失。对于向量丛 $\Bun_1^+\to\Bun_1$ 上 $\Gm$ 的伸缩作用，收缩原理给出紧支撑形式
\[i^!j_!\cF\simeq\pi_!\cF,\]
其中 $\pi:\Bun'_1\to\Bun_1$ 为投影。模去叠性的 $\Gm$ 因子后，$\pi$ 是 Abel–Jacobi 映射
\[\operatorname{AJ}^d:X^{(d)}\longrightarrow\Pic^d.\]
因此，边界项由 Deligne 定理消失：对几何不可约且秩 $n>1$ 的 $E$，或对秩为 $1$ 且非平凡的 $E$，
\[d>n(2g-2)\quad\Longrightarrow\quad\operatorname{AJ}_!^d(E^{(d)})=0.\]
对 $n=2$，这一条件恰好是 $d>4(g-1)$。
\begin{transnote}{含糊点拨}
这是 Deligne 关于 Abel–Jacobi 的消失定理在秩 $n$ 的推广，也正是 \thref{5.7.2} 在 $k=1$ 时的特例：$\Av_E^d$ 在 $\Bun_1$ 上就是 $\operatorname{AJ}^d_!$ 沿 $E^{(d)}$ 的积分。$E^{(d)}$ 是 $E^{\boxtimes d}$ 的 $S_d$-不变部分在对称积 $X^{(d)}$ 上的下降；界 $d>n(2g-2)$ 随秩线性放大。$n>1$ 时必须设 $E$ 几何不可约，是为了排除 $E$ 含来自真抛物子群或真子局部系统的分量——那时 $E^{(d)}$ 混入低秩成分，消失一般不成立；$n=1$ 的「非平凡」是同一条件的退化形式，此时即经典结果。书中只引用不证明。[依据：原书 p.34–35；并见 [KL85] 的 Abel–Jacobi 消失定理。]
\end{transnote}
\end{proof}
\section{\texorpdfstring{$\GL_n$}{GLₙ} 的朴素迭代}\label{sec:5.5}
对一般 $n$，希望迭代上述 Fourier 变换。令 $\Coh_k$ 为秩 $k$ 凝聚层的叠。定义
\[\Coh_k^+:=\{(M_k,s:\Omega_X^{k-1}\to M_k)\},\qquad\Coh'_k\subset\Coh_k^+\]
为 $s$ 在一般点单射的开轨迹。若 $s$ 在一般点单射，则 $\operatorname{coker}(s)$ 的秩为 $k-1$，所以 $\Coh'_k$ 映到 $\Coh_{k-1}$。递归从沿商映射 $q_1:\Coh'_1\to\Tor$ 的拉回开始：
\[F_1:=q_1^*\cL_E,\]
并朴素地置
\[F_{k+1}:=\mathsf F(j_{k,!}F_k)|_{\Coh'_{k+1}},\qquad j_k:\Coh'_k\hookrightarrow\Coh_k^+,\]
其中将 $\Coh'_{k+1}$ 看作 $\Coh_k$ 上 Serre 对偶扩张丛的开子叠。最后，$F_n$ 限制到 $\Bun'_n$ 后，就是 $\Aut'_E$ 的 Fourier 描述。

不能简单地将所有中间的 $\Coh_k$ 限制到 $\Bun_k$。若对 $k<n$ 有 $C_k\subset\Bun_k$，则该构造只能看到紧化 Whittaker 叠中间商为向量丛的部分。
\srcpage{35}
这样会漏掉某些 $M_k\to M$ 具有挠余核的缺陷层。因此，正确的开子叠 $C_k\subset\Coh_k$ 必须比 $\Bun_k$ 更大。
\section{开子叠 \texorpdfstring{$C_k$}{Cₖ} 与两个所需条件}\label{sec:5.6}\index{归一化次数|hyperpage}
这一构造把 $C_k$ 的最终定义推迟了。这里使用的两个条件如下。设 $M_k\in C_k$，以 $M_k^{\mathrm{tf}}$ 表示它的无挠商。

首先，向量丛条件为
\[\Ext^1(\Omega_X^{k-1},M_k)\simeq\Ext^1(\Omega_X^{k-1},M_k^{\mathrm{tf}})=0.\]
它意味着 $\Coh_k^+\to\Coh_k$ 是纤维为 $\Hom(\Omega_X^{k-1},M_k)$ 的向量丛。同样的 Serre 对偶计算还给出 $\Hom(M_k,\Omega_X^k)=0$，所以对偶扩张叠没有自同构方向。它的线性对偶是扩张类构成的真正向量丛
\[\mathcal E_{k+1}\longrightarrow C_k,\qquad0\longrightarrow\Omega_X^k\longrightarrow M_{k+1}\longrightarrow M_k\longrightarrow0.\]
开子叠 $C'_{k+1}\subset\mathcal E_{k+1}$，是中间项属于所选开子叠 $C_{k+1}$ 且诱导截面 $\Omega_X^k\to M_{k+1}$ 在一般点单射的轨迹。其对偶性来自 Serre 对偶：
\[\Ext^1(M_k,\Omega_X^k)\simeq\Hom(\Omega_X^{k-1},M_k)^\vee.\]

其次，归一化次数必须足够大：
\[\ndeg(M_k^{\mathrm{tf}}):=\deg(M_k^{\mathrm{tf}})-\deg(\cO_X\oplus\Omega_X\oplus\cdots\oplus\Omega_X^{k-1})>nk(2g-2).\]
从这里起，即使省去上圈，所有次数也都指归一化次数 $\ndeg$。
\section{洁净性与平均消失}\label{sec:5.7}\index{平均函子|hyperpage}\index{消失定理|hyperpage}
在需要进行下一次 Fourier 变换的阶段，预期的洁净延拓陈述为：
\begin{theorem}{洁净性}{5.7.1}
对 $1\le k<n$，延拓 $j_{k,!}F_k$ 至少在 $\Bun_k\cap C_k\subset\Coh_k$ 上洁净。
\end{theorem}
若 $F_k$ 是 $C'_k$ 上的单反常层，则这一洁净延拓仍为单层。在 $C_k$ 的定义中加入额外条件后，同样的洁净性将在整个 $C_k$ 上成立。

洁净性的输入是 Gaitsgory 关于平均函子的消失定理\cite{Gai04,Gai16b}。令 $\operatorname{Mod}_k^d$ 为 Hecke 叠
\[\operatorname{Mod}_k^d=\{M\subset M',\ \length(M'/M)=d\},\]
具有映射
\origdiagram{p035-diagram.pdf}{0.3438}
其中 $\pi(M\subset M')=M'/M$。平均函子为
\[\Av_E^d(-):=\overrightarrow h_!(\overleftarrow h^*(-)\otimes\pi^*\cL_E^d)[dk](dk/2).\]
\srcpage{36}
\begin{theorem}{消失}{5.7.2}
若 $E$ 几何不可约且秩为 $n$，$k<n$，并且 $d>nk(2g-2)$，则
\[\Av_E^d=0.\]
更一般地，当 $\rank(E)>k$ 时，消失定理适用于 $\Bun_k$ 上的平均函子，其界为 $d>(2g-2)k\cdot\rank(E)$。
\begin{transnote}{含糊点拨}
两个界要合起来读：第一个界 $d>nk(2g-2)$ 里 $n=\rank E$、$k$ 是 $\Bun_k$ 的下标；第二个界 $d>(2g-2)k\rank E$ 把 $k=\rank E$ 的情形单列出来。因为 $\Av_E^d$ 的核是 $\pi^*\cL_E^d$，而 $\cL_E^d$ 的「体积」随 $\rank E$ 与 $d$ 同时增长，两条界只是同一增长率的不同写法。在整章证明里，这个定理只负责「挠部分为零」这一环，其余交给 Fourier 变换的 $t$-正合性；证明本身不在书中展开。[依据：原书 p.36 定理 5.7.2 及其后的说明；并见 [Gai04]、[Gai16b]。]
\end{transnote}
\end{theorem}

当 $k=1$ 时，这就是 Deligne 关于 Abel–Jacobi 的消失定理。\chref{6}用\thref{5.7.2}证明洁净性定理。

\chapter{洁净性}\label{ch:6}\index{洁净性|hyperpage}
% 原PDF第36页中部
\section{开子叠 \texorpdfstring{$C_k$}{Cₖ} 上的 Fourier 塔}\label{sec:6.1}
本章展开 Fourier 变换递归。对 $1\le k\le n$，令
\[\Coh_k:=\{M_k\in\Coh(X)\mid\rank(M_k)=k\}.\]
有两个截面叠：
\[\Coh'_k:=\{\Omega_X^{k-1}\to M_k\text{ 在一般点单射}\},\qquad\overline{\Coh}'_k:=\{\Omega_X^{k-1}\to M_k\}.\]
上横线允许截面不是单射或为零。固定开子叠 $C_k\subset\Coh_k$，并记
\[C'_k:=\Coh'_k\times_{\Coh_k}C_k,\qquad\overline C'_k:=\overline{\Coh}'_k\times_{\Coh_k}C_k.\]
要求这些开子叠满足以下条件。
\begin{enumerate}
\item\label{cond:6.1.1} 对 $M_k\in C_k$，有 $\Ext^1(\Omega_X^{k-1},M_k)=0$。因此，$\overline C'_k\to C_k$ 是纤维为 $\Hom(\Omega_X^{k-1},M_k)$ 的向量丛。
\item\label{cond:6.1.2} 若 $M_k^{\mathrm{tf}}$ 是无挠商，则
\[\ndeg(M_k^{\mathrm{tf}})>nk(2g-2).\]
证明使用这一归一化次数形式。
\item\label{cond:6.1.3} 若 $M_k\in C_k$ 且 $\Omega_X^{k-1}\hookrightarrow M_k$ 在一般点单射，则 $M_k/\Omega_X^{k-1}\in C_{k-1}$。
\end{enumerate}
由 Serre 对偶，$\overline C'_k\to C_k$ 的线性对偶的纤维为
\[\Ext^1(M_k,\Omega_X^k).\]
记这一对偶向量丛为 $\mathcal E_{k+1}\to C_k$；它的点是扩张类
\[0\longrightarrow\Omega_X^k\longrightarrow M_{k+1}\longrightarrow M_k\longrightarrow0.\]
这就是下一次 Fourier 变换从层级 $k$ 移到层级 $k+1$ 的原因。

归纳定义 $F_k$。首先，
\[F_1:=q_1^*\cL_E[\text{移位}](\text{扭曲}),\]
其中 $q_1:C'_1\to\Tor$ 取 $\cO_X=\Omega_X^0\hookrightarrow M_1$ 的商。若 $j_k:C'_k\hookrightarrow\overline C'_k$ 为开嵌入，则
\[F_{k+1}:=\mathsf F_{\overline C'_k/C_k}(j_{k,!}F_k)|_{\mathcal E_{k+1}^{\ne0}}.\]
这里 $\mathcal E_{k+1}^{\ne0}$ 是中间项属于 $C_{k+1}$ 且典范截面 $\Omega_X^k\to M_{k+1}$ 在一般点单射的开部分。不会造成混淆时，借助从扩张类到中间项的遗忘映射，将这一限制看作 $C'_{k+1}$ 上的层。函子 $\mathsf F_{\overline C'_k/C_k}$ 是 $C_k$ 上两个对偶向量丛之间的 Fourier–Deligne 变换。
\begin{transnote}{含糊点拨}
这里的「相对」有两层：$\mathcal E_{k+1}$ 与 $\overline C'_k$ 都是 $C_k$ 上的向量丛（后者是其紧化），所以变换是逐纤维做的 §5.2，移位与扭转按相对维数取，而 $C_k$ 本身不动。先取 $j_{k,!}$ 再变换，是为了让结果只落在 $\mathcal E_{k+1}^{\ne 0}$ 上：中间项在一般点单射正是下一层叠 $C'_{k+1}$ 的定义条件，缺了 $!$-延拓，支撑会伸到有挠缺陷的边界上。[依据：原书 p.37–38 本段与 §6.2 的定义；并见 [KL85] 的 Fourier 变换。]
\end{transnote}
\srcpage{37}

在 $\Bun'_n\cap C'_n$ 上，结果是初步的自守层：
\[F_n|_{\Bun'_n\cap C'_n}=\Aut'_E|_{\Bun'_n\cap C'_n}.\]
更一般地，$F_n$ 由紧化 Whittaker 叠计算：
\[Q_n=\left\{\begin{array}{l}0=M_0\subset M_1\subset\cdots\subset M_n,\quad M_i/M_{i-1}\simeq\Omega_X^{n-i},\\M_n\hookrightarrow M\in C_n\text{，余核为挠层}\end{array}\right\}.\]
它取代未紧化的 Whittaker 叠。由于必须保留挠缺陷，允许目标属于 $\Coh_n$，而不仅属于 $\Bun_n$。
\section{洁净性定理}\label{sec:6.2}
在形成下一次 Fourier 变换所需的层级上，目标为：
\begin{theorem}{洁净性}{6.2.1}
对 $1\le k<n$，$F_k$ 跨越非单射轨迹的延拓是洁净的，即在 $\overline C'_k$ 上有
\[j_{k,!}F_k\simeq j_{k,*}F_k.\]
因此，Fourier 归纳在所选开子叠上保持所需的不可约反常层。
\end{theorem}

按挠部分的长度过滤。置
\[\Coh_k^{\mathrm{tor},\le\ell}:=\{M_k\in\Coh_k\mid\length(\operatorname{Tor}(M_k))\le\ell\},\qquad C_k^{\mathrm{tor},\le\ell}:=C_k\cap\Coh_k^{\mathrm{tor},\le\ell}.\]
则 $\Coh_k^{\mathrm{tor}=0}=\Bun_k$，$C_k=\bigcup_{\ell\ge0}C_k^{\mathrm{tor},\le\ell}$。对 $\ell$ 归纳，证明 $\overline C'_k{}^{\mathrm{tor},\le\ell}$ 上的洁净性。
\section{基础情形：无挠层}\label{sec:6.3}
\begin{proof}
取 $\ell=0$。从线丛到向量丛的映射要么为零，要么为单射。因此，$C'_k{}^{\mathrm{tor}=0}$ 在 $\overline C'_k{}^{\mathrm{tor}=0}$ 中的补集是零截面
\[i:C_k^{\mathrm{tor}=0}\hookrightarrow\overline C'_k{}^{\mathrm{tor}=0}.\]
令 $q:\overline C'_k{}^{\mathrm{tor}=0}\to C_k^{\mathrm{tor}=0}$ 为向量丛投影。该层沿纤维上的伸缩作用的轨道局部常值。与 $\GL_2$ 情形一样，洁净性等价于边界消失 $i^*j_{k,*}F_k=0$；由 Verdier 对偶，只需对 $F_k$ 及其对偶证明紧支撑版本。收缩原理给出
\[i^!j_{k,!}F_k\simeq q_!j_{k,!}F_k.\]
\begin{transnote}{易错预警}
用 $i^!$ 而不是 $i^*$：$i$ 是零截面的闭嵌入，$i^!$ 才天然配对紧支撑推前，收缩原理给出的也正是 $!$-版本。经 Verdier 对偶，$i^!j_{k,!}\simeq q_!j_{k,!}$ 对偶到 $i^*j_{k,*}\simeq q_*j_{k,*}$，所以只需证 $i^*j_{k,*}F_k=0$ 就能同时得到 $j_!\simeq j_*$，即 $j_{!*}=j_!=j_*$。前提是 $F_k$ 沿 $\Gm$-伸缩轨道局部常值（monodromic），否则收缩原理不适用。[来源：nLab "vanishing cycle"，https://ncatlab.org/nlab/show/vanishing+cycle（Last revised 2015-01-25；2026-11-17 访问）；并见原书 p.39 本段。]
\end{transnote}
所以只需证明 $q_!j_{k,!}F_k=0$；将同样的论证用于对偶，得到 $*$ 边界消失。

使用 $F_k$ 的 Whittaker 公式，这个对象为
\[q_!j_{k,!}F_k\simeq(r\circ\pi)_!(p_1^*\cL_E\otimes p_2^*\cL_\psi)[\text{移位}](\text{扭曲})|_{C_k^{\mathrm{tor}=0}},\]
其中 $p_1$ 记录挠商，$p_2$ 为 Whittaker 特征标。在无挠零截面边界上，紧化 Whittaker 叠限制为普通 Whittaker 叠
\[\Bun_{N_k}\xrightarrow{\pi}\Bun'_k\xrightarrow{r}\Bun_k.\]
令 $W_k$ 为基础 Whittaker 层，即将 $\exp_\psi$ 从 $\Bun_{N_k}$ 直接像到 $\Bun_k$ 所得的层，采用相同的移位与扭曲。在归一化次数 $d$ 的分支上，
\begin{transnote}{含糊点拨}
$\exp_\psi$ 指把 $\Ga$ 上的 Artin–Schreier 层沿 $\Bun_{N_k}$ 的简单根坐标拉回所得的层，其直像即 Whittaker 层；这里取 $*$-直像而非 $!$-直像，是因为 $W_k$ 要作为非紧支撑的核参与后面的 $!$-卷积。移位与扭转跟 $\Bun_k$ 的反常 $t$-结构对齐，相对维数与 §5.2 中 $[r](r/2)$ 的 $r$ 同一个来源：$\Bun_{N_k}\to\Bun_k$ 的纤维维数。[依据：原书 p.39 本段；并见 p.22 §3.1 对 Whittaker 特征标的定义。]
\end{transnote}
\[q_!j_{k,!}F_k\simeq\Av_E^d(W_k)|_{C_k^{\mathrm{tor}=0}}.\]
\srcpage{38}
归一化次数的约定使 $\Bun_{N_k}$ 位于 $\Bun_k$ 的次数零分支之上。\hyperref[cond:6.1.2]{条件(2)}给出
\[C_k^{\mathrm{tor}=0}\subset\bigcup_{d>nk(2g-2)}\Bun_k^d.\]
由于所考虑的归纳层级满足 $k<n$，\thref{5.7.2}说明：对 $d>nk(2g-2)$，$\Av_E^d=0$。因此 $q_!j_{k,!}F_k=0$，基础情形得证。
\end{proof}
\section{归纳步骤与基变换}\label{sec:6.4}
\begin{proof}
假设已知 $C_k^{\mathrm{tor},\le\ell-1}$ 上的洁净性。为避免将目标层与 Whittaker 旗的终点混淆，将截面的目标记为 $M$。新的边界层为
\[Z:=\{\Omega_X^{k-1}\to M\mid\text{不是单射},\ \length(\operatorname{Tor}(M))=\ell\}.\]
等价地，在 $\overline C'_k{}^{\mathrm{tor},\le\ell}$ 中，已经处理的开部分为
\[U:=C'_k{}^{\mathrm{tor},\le\ell}\cup\overline C'_k{}^{\mathrm{tor},\le\ell-1},\]
$Z$ 为剩余的闭部分。

当 $\ell>0$ 时，不存在从挠长度不超过 $\ell$ 的轨迹到挠长度等于 $\ell$ 的轨迹的典范投影。置 $Y_\ell:=\overline C'_k{}^{\mathrm{tor},\le\ell}$，并通过基变换添加这样的表现。定义
\[\widetilde Y:=\left\{\begin{array}{l}\Omega_X^{k-1}\to M,\quad0\to\mathring M\to M\to\mathcal T\to0,\\\mathring M\in\Bun_k,\ \mathcal T\in\Tor^\ell,\ M\in C_k,\ \mathring M\oplus\mathcal T\in C_k\end{array}\right\}.\]
忘去表现的映射 $\widetilde Y\to Y_\ell$ 在挠长度 $\ell$ 的层上光滑。不带截面时，这就是通过向量丛的修改给出 $\Coh_k$ 中挠长度 $\ell$ 的层的通常表现；相对切空间为 $\Hom(\mathring M,\mathcal T)$，所以相对维数局部常值。

固定点轨迹为
\[\widetilde Z:=\{\Omega_X^{k-1}\to\mathcal T,\quad\mathring M\in\Bun_k,\ \mathring M\oplus\mathcal T\in C_k\}.\]
它通过分裂扩张 $M=\mathring M\oplus\mathcal T$ 映入 $\widetilde Y$。注意 $\widetilde Z\ne Z$：在 $\widetilde Z$ 中，分裂／表现是数据的一部分。有投影
\[\widetilde f:\widetilde Y\longrightarrow\widetilde Z,\]
\[(\Omega_X^{k-1}\to M,\ 0\to\mathring M\to M\to\mathcal T\to0)\longmapsto(\Omega_X^{k-1}\to M\to\mathcal T,\mathring M).\]
对 $\Ext^1(\mathcal T,\mathring M)$ 中的扩张类作标量乘法，同时对所提升截面的 $\mathring M$ 分量作标量乘法，定义 $\widetilde Y$ 上的 $\Gm$ 作用，并延伸为收缩到 $\widetilde Z$ 的 $\A^1$ 作用。只缩放扩张类，会在固定点留下一个不需要的截面的向量丛分量。因此，收缩原理将归纳步骤归约为
\[\widetilde f_!\widetilde j_!(F_k|_{\widetilde U})=0,\qquad\widetilde U:=U\times_{\overline C'_k}\widetilde Y.\]
\end{proof}
\begingroup\setlength{\intextsep}{3pt}\setlength{\abovedisplayskip}{3pt}\setlength{\belowdisplayskip}{3pt}\setlength{\belowdisplayshortskip}{2pt}
\titlespacing*{\section}{0pt}{1.5ex plus .5ex}{1ex plus .2ex}
\section{纤维计算与满射情形}\label{sec:6.5}
\begin{proof}
使用紧化 Whittaker 公式
\[F_k\simeq\pi_!(p_1^*\cL_E\otimes p_2^*\cL_\psi)[\text{移位}](\text{扭曲})|_{C'_k},\]
其中 $p_1$ 记录挠商，$p_2$ 是 Whittaker 特征标。基变换后，取定一点 $(s_{\mathcal T}:\Omega_X^{k-1}\to\mathcal T,\mathring M)\in\widetilde Z$。相关纤维由 Whittaker 旗以及固定扩张
\[0\longrightarrow\mathring M\longrightarrow M\longrightarrow\mathcal T\longrightarrow0\]
之上的映射 $M_k\to M$ 组成，
\srcpage{39}
使复合 $\Omega_X^{k-1}=M_1\to M_k\to M\to\mathcal T$ 为固定映射 $s_{\mathcal T}$。诱导映射 $M_k\to\mathcal T$ 未必为满射。

若 $M_k\to\mathcal T$ 为满射，令 $\mathring M_k:=\ker(M_k\to\mathcal T)$，则有正合图
\origdiagram{p039-exact.pdf}{0.4668}
纤维等价于普通修改叠
\[\{\mathring M_k\to\mathring M\mid\mathring M/\mathring M_k\in\Tor^d\}.\]
因此，所需 $!$ 推前是 $\Av_E^d$ 的一个茎。在归一化次数下，Whittaker 旗的终点 $M_k$ 满足 $\ndeg M_k=0$，故 $\ndeg\mathring M_k=-\ell$。由于 $\mathring M\oplus\mathcal T\in C_k$，\hyperref[cond:6.1.2]{条件(2)}给出 $\ndeg\mathring M>nk(2g-2)$。于是
\[d=\ndeg\mathring M-\ndeg\mathring M_k>nk(2g-2)+\ell.\]
所以 $d>nk(2g-2)$，消失定理给出 $\Av_E^d=0$。这处理了满射情形。
\end{proof}
\section{非满射情形的归约}\label{sec:6.6}
\begin{proof}
假设 $M_k\to\mathcal T$ 不是满射，令 $\mathcal T'\subsetneq\mathcal T$ 为其像。置 $\mathring M_k:=\ker(M_k\to\mathcal T')$。令 $M'\subset M$ 为 $\mathcal T'$ 在 $M\to\mathcal T$ 下的逆像。则数据经由下图分解：
\origdiagram{p039-three.pdf}{0.4713}
从而
\[0\longrightarrow M'/M_k\longrightarrow M/M_k\longrightarrow\mathcal T/\mathcal T'\longrightarrow0.\]

令 $H_{\mathcal T}$ 和 $H_{\mathcal T'}$ 为相应修改纤维。映射 $H_{\mathcal T}\to H_{\mathcal T'}$ 与挠层图相容：
\origdiagram{p039-torses.pdf}{0.3825}
其中 $\Tor^{\mathrm{ses}}$ 分类挠层的短正合列；$m$ 记录中间项，$(s,q)$ 记录子对象与商。在当前情形，它们分别是 $M/M_k$、$M'/M_k$ 与 $\mathcal T/\mathcal T'$。
\Needspace{3\baselineskip}

拉回到这个图后，$H_{\mathcal T}\to H_{\mathcal T'}$ 是仿射空间纤维化。仿射方向由诸如 $\Ext^1(\mathcal T/\mathcal T',\mathring M_k)$ 的扩张数据，以及在商 $\mathcal T'$ 上选择相容提升 $M_k\to M'$ 控制。因此，忽略通常的仿射丛移位
\srcpage{40}
与 Tate 扭转后，通过基变换，剩余项归约为沿剩余投影 $\mathrm{pr}$ 的推前：
\[\mathrm{pr}_!((s,q)_!m^*\cL_E|_{H_{\mathcal T'}})=0.\]
所需命题是在 $\Tor$ 所用商叠归一化中的因子化恒等式
\[(s,q)_!m^*\cL_E\simeq\cL_E\boxtimes\cL_E,\]
它将在\chref{7}证明。它将挠商分为 $M'/M_k$ 部分与剩余的 $\mathcal T/\mathcal T'$ 部分，从而将非满射情形归约为满射情形。在假定该命题的情况下，洁净性的证明完成。
\begin{transnote}{含糊点拨}
这条恒等式说明：沿挠商部分把 $\cL_E$ 积掉，与先把挠商拆成 $M'/M_k$ 和 $\mathcal T/\mathcal T'$ 两段、再分别积掉，结果相同——即 $\cL_E$ 关于 $E$ 的乘法是因子化的。证明推迟到 \chref{7}（$C_k$ 的因子化定义）：先把非满射情形归约为满射，满射时 $(s,q)$ 沿纤维平坦，$\cL_E$ 逐纤维常值，恒等式退化为 $E$ 的乘法性。[依据：原书 p.40 本段；并见 \chref{7} 的因子化恒等式。]
\end{transnote}
\end{proof}
\endgroup

\chapter{Laumon 因子化}\label{ch:7}\index{Laumon 因子化|hyperpage}
% 原PDF第40页中部
\section{洁净性中剩下的引理}\label{sec:7.1}
本章完成“消失推出洁净性”的证明的最后一步。固定 $M_k\in\Bun_k^0$、$J\in\Tor^\ell$、$\mathring M\in\Bun_k$ 以及 $a:M_k\to J$，其中 $\Bun_k^0$ 表示作为 Whittaker 旗终点出现的归一化次数零分支。假设 $\mathring M\oplus J\in C_k$。考虑叠
\[H:=\left\{\begin{array}{l}0\to\mathring M\to M\to J\to0,\quad M\in C_k,\\u:M_k\hookrightarrow M,\quad(M_k\xrightarrow{u}M\twoheadrightarrow J)=a\end{array}\right\}.\]
它有映射 $q_H:H\to\Tor$，$(M,u)\mapsto M/M_k$。洁净性证明所需的引理是紧支撑消失
\[\RGamma_c(H,q_H^*\cL_E)\simeq0.\]

还要对良好的开子叠 $C_k$ 作一个技术性加强。可以要求：对每个扩张 $0\to\mathring M\to M\to J\to0$，
\[\mathring M\oplus J\in C_k\quad\Longrightarrow\quad M\in C_k.\]
确保这一蕴涵的一种方法，是选取线丛 $L^{\mathrm{est}}$ 并施加额外的开条件
\[\Hom(M,L^{\mathrm{est}})=0.\]
事实上，对挠层 $J$ 和无挠目标 $L^{\mathrm{est}}$，$\Hom(J,L^{\mathrm{est}})=0$。因此 $\mathring M\oplus J\in C_k$ 等价于 $\mathring M\in C_k$，对扩张应用 $\Hom(-,L^{\mathrm{est}})$ 得到 $M\in C_k$。于是，在纤维计算中可以忽略条件 $M\in C_k$。
\section{满射情形}\label{sec:7.2}
\begin{proof}
首先假设 $a:M_k\twoheadrightarrow J$。置 $\mathring M_k:=\ker(a)$。则 $H$ 的一个对象等同于如下图：
\origdiagram{p040-exact.pdf}{0.4645}
等价地，它是一个修改 $\mathring M_k\hookrightarrow\mathring M$。因此
\[H\simeq\{\mathring M_k\}\times_{\Bun_k}\operatorname{Mod}_k^d\times_{\Bun_k}\{\mathring M\},\qquad\Bun_k\longleftarrow\operatorname{Mod}_k^d\longrightarrow\Bun_k.\]
在这一等同下，$q_H$ 记录 $\mathring M/\mathring M_k$。忽略 $\Av_E$ 定义中无碍于论证的归一化一维线后，
\[\RGamma_c(H,q_H^*\cL_E)\simeq(\Av_E^d(\delta_{\mathring M_k}))_{\mathring M}\otimes(\text{一维线}),\qquad d=\ndeg\mathring M-\ndeg\mathring M_k.\]
\srcpage{41}
在洁净性证明所用的归一化分支中，$\ndeg M_k=0$，所以 $\ndeg\mathring M_k=-\ell$。条件 $\mathring M\oplus J\in C_k$ 给出 $\ndeg\mathring M>nk(2g-2)$。故
\[d=\ndeg\mathring M-\ndeg\mathring M_k>nk(2g-2)+\ell>nk(2g-2).\]
因此，$\Av_E^d$ 的消失定理给出 $\RGamma_c(H,q_H^*\cL_E)=0$。
\end{proof}
\section{归约到 \texorpdfstring{$M_k\to J$}{Mₖ→J} 的像}\label{sec:7.3}
\begin{proof}
对一般映射 $a:M_k\to J$，令 $J'\subset J$ 为其像，置 $J'':=J/J'$。令 $\mathring M_k=\ker(M_k\twoheadrightarrow J')$。对 $H$ 的一个对象，置 $M':=M\times_JJ'$。于是有
\origdiagram{p041-three.pdf}{0.4690}
令 $H'$ 为前两行构成的叠；它是以 $J'$ 为目标的满射版本。

令 $\Tor^{\mathrm{ses}}$ 为挠层短正合列
\[0\longrightarrow J_s\longrightarrow J_m\longrightarrow J_q\longrightarrow0\]
的叠。记
\[m(J_\bullet)=J_m,\qquad s(J_\bullet)=J_s,\qquad q(J_\bullet)=J_q.\]
商序列
\[0\longrightarrow M'/M_k\longrightarrow M/M_k\longrightarrow M/M'\longrightarrow0\]
定义映射 $H\to\Tor^{\mathrm{ses}}$。这里
\[M'/M_k\simeq\mathring M/\mathring M_k,\qquad M/M'\simeq J/J'.\]
诱导图
\origdiagram{p041-diagram.pdf}{0.2914}
不是笛卡尔图。然而，对于所考虑的层，基变换仍然成立。原因在于
\[H\longrightarrow\Tor^{\mathrm{ses}}\times_{\Tor\times\Tor}H'\]
是仿射丛。在由
\[0\to J_s\to J_m\to J_q\to0,\qquad0\to M_k\to M'\to J_s\to0\]
分类的点处，纤维是仿射空间
\[\Ext^1(J_m,M_k)\times_{\Ext^1(J_s,M_k)}\{[0\to M_k\to M'\to J_s\to0]\}.\]
\srcpage{42}
其维数在每个连通分支上恒定；事实上，这一仿射空间以
\[\ker(\Ext^1(J_m,M_k)\to\Ext^1(J_s,M_k))\simeq\Ext^1(J_q,M_k)\]
为平移空间，维数为 $k\length(J_q)$。所以，忽略恒定的仿射丛移位与 Tate 扭转后，只需证明
\[\RGamma_c((s,q)_!m^*\cL_E|_{H'})=0.\]
由满射情形，这将由下面的因子化恒等式推出。
\end{proof}
\section{Laumon 层通过短正合列的因子化}\label{sec:7.4}
\begin{proposition}{Laumon 因子化\cite{KL85}}{7.4.1}
对任何能够定义 Laumon 复形的局部系统 $E$，特别是洁净性论证所用的几何不可约 $E$，在 $\Tor\times\Tor$ 上有
\[(s,q)_!m^*\cL_E\simeq\cL_E\boxtimes\cL_E.\]
\begin{transnote}{含糊点拨}
这条恒等式就是「因子化」的层化陈述：$\cL_E$ 是沿 Hecke 对应把 $E$ 拉回、再把修改方向的层推前所得的核，而 $(s,q)_!m^*\cL_E$ 把 $\Tor\times\Tor$ 上的乘法拆开，恰好对应「把挠层分成两个子挠层、再各自配 $E$」这个操作。命题 7.4.1 把它当作输入，证明在 \chref{9} 的 Laumon 层构造与 §7.1 的 Laumon 复形里完成；「能够定义 Laumon 复形」只要求所选层理论合适，几何不可约是洁净性论证另加的条件，两者不可混为一谈。[依据：原书 p.42 命题 7.4.1 及其前的构造；并见 [KL85]。]
\end{transnote}
等价地，在 $\Tor^{\mathrm{ses}}$ 中子对象长度为 $d_1$、商长度为 $d_2$ 的分支上，在 $\Tor^{d_1}\times\Tor^{d_2}$ 上有
\[(s,q)_!m^*\cL_E^{d_1+d_2}\simeq\cL_E^{d_1}\boxtimes\cL_E^{d_2},\]
其中移位与 Tate 扭转已经包含在 $\cL_E$ 的归一化中。
\end{proposition}
\begin{note}\numberednote{7.4.2}
前后的洁净性论证假设 $E$ 几何不可约，但因子化恒等式本身是 Springer–Laumon 构造的结构性质，不使用这种不可约性。
\end{note}
\begin{proof}
回顾正则半单轨迹
\[\Tor^{d,\mathrm{rss}}=\left\{\bigoplus_{i=1}^dk(x_i)\ \middle|\ x_i\ne x_j\text{ 当 }i\ne j\right\}\subset\Tor^d,\]
它是 $X^{(d),\mathrm{rss}}=X^{(d)}-\Delta$ 上的胚叠，其在几何点处的稳定子为 $(\Gm)^d$。
\begin{transnote}{含糊点拨}
胚叠是指局部同构于 $BG$ 的叠；这里每个 $k(x_i)$ 上的标量自同构各给出一个 $\Gm$，所以稳定子是 $(\Gm)^d$，而商叠局部是 $B(\Gm^d)$。这一步是为下一行的对称下降铺路：正因为 $\Tor^{d,\mathrm{rss}}$ 相对于 $X^{(d),\mathrm{rss}}$ 只剩 $(\Gm)^d$ 的自同构，$E^{\boxtimes d}$ 沿有序支撑覆盖的推前才带 $S_d$-作用，于是可以取 $S_d$-不变部分把它下降到 $X^{(d),\mathrm{rss}}$ 上。[来源：nLab "gerbe"，https://ncatlab.org/nlab/show/gerbe（Last revised 2026-02-14；2026-11-17 访问）；并见原书 p.42 本段与第 4 章对 $(\Gm)^d$-胚叠的讨论。]
\end{transnote}若 $r_{\mathrm{rss}}:X^{d,\mathrm{rss}}\to X^{(d),\mathrm{rss}}$ 为有序支撑覆盖，$s_{\mathrm{rss}}:\Tor^{d,\mathrm{rss}}\to X^{(d),\mathrm{rss}}$ 为支撑映射，则
\[E^{(d)}:=(r_{\mathrm{rss},*}E^{\boxtimes d})^{S_d},\qquad\cL_E^d\simeq j_{\mathrm{rss},!*}(s_{\mathrm{rss}}^*E^{(d)}[d]).\]
\begin{transnote}{含糊点拨}
$j_{\mathrm{rss},!*}$ 是中间延拓（IC 延拓）：对开嵌入 $j_{\mathrm{rss}}$，它是像 $j_!\to j_*$ 所对应的那个中间层，既非紧支撑延拓 $j_!$，也非 $j_*$。取它才得到 Goresky–MacPherson 延拓，与 §9.1 定义的 Laumon 层一致。$[d]$ 把 $E^{(d)}$ 的度数移到 $\Tor^d$ 的维数上，使 $\cL_E^d$ 落在反常 $t$-结构的度 0 层；若忘记这个移位，下一章的洁净性就会差一个平移。[依据：原书 p.42 本式与 §9.1；并见 nLab "intermediate extension" 的中间延拓定义（该条目无 Last revised，2026-11-17 访问）。]
\end{transnote}
令
\[U_{d_1,d_2}:=(\Tor^{d_1,\mathrm{rss}}\times\Tor^{d_2,\mathrm{rss}})_{\mathrm{disj}}\]
为两个支撑互不相交的开轨迹，令 $s_i:\Tor^{d_i,\mathrm{rss}}\to X^{(d_i),\mathrm{rss}}$ 为两个支撑映射。在这一轨迹上，两个挠层之间的每条短正合列都唯一分裂，中间项映射 $m$ 等同于除子的加法。因此
\[(s,q)_!m^*\cL_E^d|_{U_{d_1,d_2}}\simeq\operatorname{add}^*s_{\mathrm{rss}}^*E^{(d)}[d]|_{U_{d_1,d_2}},\]
而
\[(\cL_E^{d_1}\boxtimes\cL_E^{d_2})|_{U_{d_1,d_2}}\simeq(s_1^*E^{(d_1)}[d_1]\boxtimes s_2^*E^{(d_2)}[d_2])|_{U_{d_1,d_2}}.\]
这两个局部系统通过对称下降典范等同。此外，右边是其到 $U_{d_1,d_2}$ 的限制的 IC 延拓：乘积局部系统已经光滑地跨越两个支撑相遇的交叉对角线延拓。因此，只需证明
\[(s,q)_!m^*\cL_E^d\simeq\operatorname{IC}((s,q)_!m^*\cL_E^d|_{U_{d_1,d_2}}).\]

由于 $\cL_E^d=(\operatorname{Spr}_E^d)^{S_d}$ 是 Springer–Laumon 层的直和项，只需对 $\operatorname{Spr}_E^d$ 证明相应陈述，使其对 $S_{d_1}\times S_{d_2}\subset S_d$ 等变，再取
\srcpage{43}
不变直和项。该断言在 $X$ 上是 étale 局部的。作这样的局部化后，局部系统 $E$ 成为纤维为 $V$ 的常值系统；系数因子 $V^{\otimes d}$ 随整个构造带入。因此，几何输入就是常系数的 Springer 陈述。

对 $X=\A^1$，长度 $d$ 挠层是带有一个自同态的 $d$ 维向量空间，因此
\[\Tor^d\simeq\mathfrak{gl}_d/\GL_d.\]
\begin{transnote}{含糊点拨}
$\A^1=\Spec k[x]$ 上长度 $d$ 的挠层就是一个 $d$ 维 $k[x]$-模：$x$ 的作用给出一个 $d\times d$ 矩阵，模同构就是 $\GL_d$-共轭，于是模叠恰是 $\mathfrak{gl}_d/\GL_d$，全空间都允许，不额外要求 $x$ 幂零。$x$ 幂零的那些点构成其中的幂零锥 $\cN_d$，正是第四章局部模型在 $0$ 附近看到的部分；这里取整块 $\mathfrak{gl}_d$，才能让 $\Tor^d$ 在 $0$ 处带出幂零锥。[来源：nLab "coherent sheaf"，https://ncatlab.org/nlab/show/coherent+sheaf（Last revised 2026-01-25；2026-11-17 访问）；并见原书 p.43 本段。]
\end{transnote}
对 $d=d_1+d_2$，叠 $\Tor^{\mathrm{ses},d_1,d_2}$ 是抛物对应
\[\mathfrak g/G\longleftarrow\mathfrak p/P\longrightarrow\mathfrak m/M,\]
其中 $G=\GL_d$，$P=P_{d_1,d_2}=UM$，$M\simeq\GL_{d_1}\times\GL_{d_2}$，$\mathfrak p$、$\mathfrak m$ 为相应 Lie 代数。所需几何输入是 Grothendieck–Springer 层的标准抛物限制定理\cite{AHJR14,GPR20}：复形
\[(\mathfrak p/P\to\mathfrak m/M)_!(\mathfrak p/P\to\mathfrak g/G)^*\operatorname{Spr}_{\mathfrak g}\]
是从 $\mathfrak m/M$ 中两个特征多项式正则且互不相交的轨迹作出的 IC 延拓，其 Weyl 局部系统由将 $S_d$ 单值作用限制到 $S_{d_1}\times S_{d_2}$ 得到。

作为参考，Bruhat 证明将
\[\widetilde{\mathfrak p}:=\{(x,gB)\mid x\in\mathfrak p\cap\operatorname{Ad}_g\mathfrak b\}\]
按双陪集部分 $PwB/B$ 分层，其中 $w\in W_M\backslash W_G$。对最小陪集代表 $w$，到 $\mathfrak m$ 的 Grothendieck–Springer 解消的映射是仿射纤维化，其纤维为
\[\mathfrak u_w:=\mathfrak u\cap\operatorname{Ad}_w\mathfrak b,\]
而非整个 $\mathfrak u$；维数依赖于 $w$。采用通常的反常归一化，这些分层组合成上述抛物限制的 IC 延拓，不会出现支撑在正则且不交轨迹的补集上的直和项。

带着系数 $V^{\otimes d}$ 取 $S_d$ 不变直和项，将这一 IC 延拓等同于 $\cL_E^{d_1}\boxtimes\cL_E^{d_2}$，因为在正则且不交轨迹上，它恰好就是已描述的对称下降等同。这证明了 Laumon 因子化命题。
\end{proof}
\section{完成引理}\label{sec:7.5}
\begin{proof}
将命题\ref{pro:7.4.1}应用于上文所得剩余表达式：
\[(s,q)_!m^*\cL_E|_{H'}\simeq(\cL_E\boxtimes\cL_E)|_{H'}.\]
第一个因子恰好是 $M_k\twoheadrightarrow J'$ 的满射情形，第二个因子是固定的剩余商 $J/J'$。所以，由满射情形，沿 $H'$ 的 $!$ 推前消失。再由上述仿射丛基变换，
\[\RGamma_c(H,q_H^*\cL_E)=0.\]
这证明了剩下的引理，并完成洁净性证明中缺少的一步。
\end{proof}

\chapter{\texorpdfstring{$\Aut_E$}{Aut（E）} 的构造}\label{ch:8}
\srcpage{44}
\section{已经完成的构造}\label{sec:8.1}
继续取 $G=\GL_n$，并令 $E$ 为 $X$ 上几何不可约的秩 $n$ 局部系统。第\ref{ch:5}–\ref{ch:7}章在良好的开子叠 $C'_n\subset\Coh'_n$ 上构造了初步的 Whittaker–Laumon 层
\[\Aut'_E\in D(C'_n).\]
它通过公式
\[\Aut'_E=\pi_!(p_1^*\cL_E\otimes p_2^*\cL_\psi)[\text{移位}](\text{扭曲})\]
从紧化 Whittaker 叠
\origdiagram{p044-diagram.pdf}{0.4258}
得到。已经证明的定理是：在 $C'_n$ 上，这个对象是不可约反常层。

令
\[C_n^\circ:=C_n\cap\Bun_n,\qquad C_n^{\prime,\circ}:=C'_n\cap\Bun'_n,\]
并令
\[p_C:C_n^{\prime,\circ}\longrightarrow C_n^\circ\]
为将带有一般点非零截面 $\Omega_X^{n-1}\to M$ 的丛映为丛 $M$ 的遗忘映射。待证明的下降陈述是：若 $C_n\cap\Bun_n^d\ne\varnothing$，则
\[\Aut'_E|_{C_n^{\prime,\circ}\cap(\Bun_n^d)'}\]
沿 $p_C$ 下降为层
\[\Aut_E|_{C_n\cap\Bun_n^d}\in\operatorname{Perv}(C_n\cap\Bun_n^d).\]
由于 $p_C$ 在所选良好开子叠上光滑，非零的下降层自动是不可约反常层。最终目标是将这些局部部分唯一延拓为整个 $\Bun_n$ 上的 Hecke 特征层 $\Aut_E$。
\section{沿射影丛纤维下降}\label{sec:8.2}\index{下降!射影丛|hyperpage}
在光滑开图 $V\subset C_n^\circ\cap\Bun_n^d$ 上工作。必要时缩小 $V$ 后，截面 $\Omega_X^{n-1}\to M$ 的空间是 $V$ 上的向量丛；再对截面上的标量作用取商，得到射影丛型映射
\[p_V:V'\longrightarrow V.\]
下降的图景是：先将 $\Aut'_E|_{V'}$ 写为 $V'$ 的某个稠密开部分上局部系统的 IC 延拓，再证明这个 IC 延拓实际上是整个 $V'$ 上的局部系统。

使用两个条件。首先，对每个几何点 $m\in V$，良好开部分与纤维的交是稠密的：
\[V'\cap p_V^{-1}(m)\subset p_V^{-1}(m)\quad\text{是开且稠密的。}\]
可以将它纳入良好开部分的定义。其次，$\Aut'_E|_{V'}$ 必须是局部系统，而不只是 IC 层。知道这一点后，局部系统沿 $p_V$ 唯一下降，因为几何纤维是射影空间，且
\[\pi_1(\mathbb P_{\bar k}^m)=1.\]
\srcpage{45}
等价地，从 $V$ 上局部系统到 $V'$ 上局部系统的拉回函子是等价。这给出了所需的从 $C_n^{\prime,\circ}$ 到 $C_n^\circ$ 的下降。

用于证明不可约反常层为局部系统的判据，是当前纯 IC 情形中的有限域迹检验。设 $Y$ 光滑，$\cK\in\operatorname{Perv}(Y)$ 不可约。记
\[\chi_{\cK}(y):=\sum_i(-1)^i\dim H^i(\cK_{\bar y})\]
为茎的 Euler 示性数。对任意反常层，这个函数本身并不是局部系统判据。这里 $\cK$ 是由该构造得到的不可约 IC 延拓；在这一情形，$\chi_{\cK}$ 恒定且非零排除了边界贡献，从而迫使 $\cK$ 是 $Y$ 上移位的局部系统。因此，只需证明：
\[\chi(\Aut'_E)\text{ 在 }\Bun'_n\to\Bun_n\text{ 的纤维上恒定},\qquad\chi(\Aut'_E)\ne0\text{ 在 }C_n^\circ\text{ 上。}\]
\section{沿纤维的迹恒定性}\label{sec:8.3}
$\chi(\Aut'_E)$ 的恒定性归约到已知存在自守函数的情形。作为黑箱输入，使用 Deligne 的真推前不变性：若 $f:Y_1\to Y_2$ 为真态射，$\cK_1,\cK_2\in D(Y_1)$ 在 étale 局部同构，则 $f_!\cK_1$ 与 $f_!\cK_2$ 的 Euler 示性数函数相同。
\begin{transnote}{含糊点拨}
Euler 示性数函数可按分层计算：把 $Y_1$ 细分到 $\cK_1$ 与 $\cK_2$ 都局部常值（带平移与扭转）的层上，两层在每个分层上只差一个常数因子，故 $\chi$ 相同；真性保证 $f_!\cK_i$ 仍是构造性的，$\chi$ 才有定义。书中把它当黑箱，是因为只需借它把「沿 $\Bun'_n\to\Bun_n$ 纤维恒定」转移到某一个方便的自守函数上，并不需要 Deligne 定理的内部机制。[来源：nLab "six operations"，https://ncatlab.org/nlab/show/six+operations（Last revised 2026-09-30；2026-11-17 访问）；并见原书 p.45 本段。]
\end{transnote}

对任意两个秩 $n$ 局部系统 $E_1,E_2$，即使它们可约，在 $X$ 上也 étale 局部同构。因此，相应的 Laumon 层 étale 局部同构：
\[\cL_{E_1}\simeq_{\mathrm{et,loc}}\cL_{E_2}.\]
若定义 $\Aut'_E$ 的映射是真态射，Deligne 定理立即给出
\[\chi(\Aut'_{E_1})=\chi(\Aut'_{E_2}).\]
该映射并非字面上的真态射，但紧化 Whittaker 构造有 Drinfeld 紧化，所以将对应替换为其紧化版本后，同样的论证适用。

因此，对给定 $E$ 证明恒定性命题，只需对方便的秩 $n$ 局部系统 $E'$ 证明。选择 $E'$，使存在以 $E'$ 为 Hecke 特征值的尖点 Hecke 特征函数。由 Frenkel–Gaitsgory–Kazhdan–Vilonen 的比较定理，它在 $\Bun'_n(\mathbb F_q)$ 上的限制是迹函数
\[M'\longmapsto\Tr(\operatorname{Frob}_{M'},\Aut'_{E'}).\]
\begin{transnote}{含糊点拨}
比较定理的两边是：几何侧的特征层 $\Aut'_{E'}$ 在 $\Bun'_n(\mathbb F_q)$ 上的迹函数，与自守侧的 Whittaker 归一化特征函数。定理说这两者相等（在 $E'$ 为 Hecke 特征值的前提下），于是「数点」与「层」可以互换。这里限制到 $\Bun'_n$ 而不是 $\Bun_n$，是因为 Whittaker 归一化正落在 $\Bun'_n$ 上；$E'$ 之所以要选成「以尖点 Hecke 特征函数为特征值」，是为了让迹函数非零且可算。[依据：原书 p.45 本段；并见 [FGKV98]（JAMS 11 (1998), 451–484）的层–函数比较定理。]
\end{transnote}
由于这个函数从 $\Bun_n(\mathbb F_q)$ 拉回，它在 $\Bun'_n\to\Bun_n$ 的每个纤维上恒定。由前述归约，同样的恒定性对 $\Aut'_E$ 成立。这证明了下降定理。
\section{Hecke–Laumon 对应}\label{sec:8.4}\index{Hecke–Laumon 对应|hyperpage}
通常的 Hecke 性质由凝聚层上一个看起来更强的性质推出。对 $k\ge0$、$d\ge0$，定义 Hecke–Laumon 叠
\[\mathrm{HL}_k^d:=\{0\to M'\to M\to\mathcal T\to0\mid M,M'\in\Coh_k,\ \mathcal T\in\Tor^d\}.\]
它具有映射
\[\Coh_k\xleftarrow{\overleftarrow h}\mathrm{HL}_k^d\xrightarrow{\overrightarrow h}\Tor^d\times\Coh_k,\]
其中 $\overleftarrow h$ 记住 $M$，$\overrightarrow h$ 记住 $(\mathcal T,M')$。相应函子为
\[\mathrm{HL}_k^d(\cK):=\overrightarrow h_!\overleftarrow h^*\cK[dk](dk/2),\qquad\cK\in D(\Coh_k).\]
\begin{transnote}{易错预警}
两个因子各管一件事，缺一个就错：$[dk]$ 是 $\mathrm{HL}_k^d$ 在 $\Coh_k$ 上的相对维数（每次修改贡献 $k$，共 $d$ 次），$(dk/2)$ 是配对的 Tate 归一化，与 §5.2 的 $[r](r/2)$ 同一个来源。两者合起来才使 $\mathrm{HL}_k^d$ 对反常 $t$-结构 $t$-正合；$(dk/2)$ 可以是半整数，故系数域须含 $\sqrt q$。[来源：nLab "Tate twist"，https://ncatlab.org/nlab/show/Tate+twist（Last revised 2025-08-10；2026-11-17 访问）；并见原书 p.45–46 本式与 §5.2 的归一化约定。]
\end{transnote}
\srcpage{46}
\begin{definition*}{}
若反常层 $\cK\in\operatorname{Perv}(\Coh_k)$ 配备同构
\[\alpha_d:\mathrm{HL}_k^d(\cK)\xrightarrow{\sim}\cL_E^d\boxtimes\cK\qquad(d\ge0),\]
满足 $d=0$ 时显然的归一化，且这些同构与迭代修改相容，则称它为以 $E$ 为特征值的 Hecke–Laumon 特征层。
\end{definition*}
更明确地，对 $d=d_1+d_2$，比较通过
\[0\to M'\to M\to\mathcal T'\to0,\qquad0\to M''\to M'\to\mathcal T''\to0\]
从 $M$ 转到 $M''$ 的两种方式。这种两步修改的叠既映到 $\mathrm{HL}_k^d$，也映到 $\Tor^{d_1}\times\mathrm{HL}_k^{d_2}$。所需相容性要求由
\[\mathrm{HL}_k^{d_2}\circ\mathrm{HL}_k^{d_1}\simeq((s,q)_!m^*\boxtimes\mathrm{Id})\circ\mathrm{HL}_k^d\]
得到的图与\chref{7}的因子化同构
\[(s,q)_!m^*\cL_E^d\simeq\cL_E^{d_1}\boxtimes\cL_E^{d_2}\]
交换。这里 $m,s,q$ 是挠层短正合列叠
\[0\longrightarrow\mathcal T''\longrightarrow\mathcal T\longrightarrow\mathcal T'\longrightarrow0\]
中的映射。

对 $k=0$，这恰好是 Laumon 因子化定理：$\cL_E\in\operatorname{Perv}(\Tor)$ 是 Hecke–Laumon 特征层。前几章的 Fourier 归纳继而对 $k$ 归纳，给出
\[F_k\in\operatorname{Perv}(C'_k)\]
为 Hecke–Laumon 特征层。特别地，$\Aut'_E|_{C_n^{\prime,\circ}}$ 满足 Hecke–Laumon 特征性质。下降后，层 $\Aut_E|_{C_n^\circ}$ 继承同样的性质。
\section{从 Hecke–Laumon 到通常的 Hecke}\label{sec:8.5}
剩下的问题是，为何 Hecke–Laumon 性质蕴涵 $\Bun_n$ 上通常的 Hecke 特征性质。这里使用的陈述为：
\begin{proposition}{}{8.5.1}
设 $\mathcal S\in\operatorname{Perv}(C_n^\circ)$ 是以 $E$ 为特征值的 Hecke–Laumon 特征层。再假设其 Verdier 对偶 $\mathbb D\mathcal S$ 是以 $E^*$ 为特征值的 Hecke–Laumon 特征层。则 $\mathcal S$ 满足以 $E$ 为特征值的通常 Hecke 特征性质。
\end{proposition}
\begin{proof}

只需证明第一 Hecke 关系。考虑 $\mathrm{HL}_n^1$。去掉分裂的短正合列，得到开子叠
\[\mathrm{HL}_n^{1,\circ}:=\{0\to M'\to M\to\mathcal T\to0\mid\length(\mathcal T)=1,\ \text{序列不分裂}\}.\]
沿 $X\to\Tor^1$ 作基变换；该映射将 $x$ 映为支撑于 $x$ 的长度一挠层。所得叠为通常的第一 Hecke 对应：
\[\mathrm{HL}_n^{1,\circ}\times_{\Tor^1}X=\{(x,0\to M'\to M)\mid x\in X,\ \supp(M/M')=x,\ \length(M/M')=1\}.\]
等价地，
\[M'\subset M\subset M'(x),\qquad\length(M/M')=1.\]
长度一挠商剩余的标量自同构是一个 $\Gm$ 因子；对它取商后得到通常的 Hecke 叠 $\Hecke_n^1$。

这一商所需的反常估计由下面的初等引理提供。
\end{proof}
\srcpage{47}
\begin{lemma}{}{8.5.2}
设 $\mathcal E\xrightarrow{\pi}B$ 为向量丛，$\mathbb P(\mathcal E)\xrightarrow{\bar\pi}B$ 为其射影化。假设 $\cK\in\operatorname{Perv}(\mathcal E)$ 是 $\Gm$-等变的，且 $\pi_!\cK,\pi_*\cK\in\operatorname{Perv}(B)$。令 $\overline{\cK}$ 为 $\cK|_{\mathcal E\setminus0}$ 到 $\mathbb P(\mathcal E)$ 的下降。则 $\bar\pi_!\overline{\cK}$ 是反常层。
\end{lemma}
\begin{proof}

将\lemref{8.5.2}逐纤维应用于从 $\mathrm{HL}_n^{1,\circ}$ 到通常 Hecke 叠的商。\secref{8.4}中的 Hecke–Laumon 特征性质给出 $\pi_!\cK$ 所需的反常性，对 $\mathbb D\mathcal S$ 的同一性质给出 $\pi_*\cK$ 的反常性。由于
\[\Tor^1\simeq X\times B\Gm\]
且 $\cL_E^1|_X=E$，$\mathcal S$ 的第一通常 Hecke 变换是反常的，且特征值为 $E$。较高的微小 Hecke 关系由第一 Hecke 对应在卷积下的相容性得到，特征值为 $\bigwedge^iE$。因此，$\mathcal S$ 是通常的 Hecke 特征层。
\end{proof}

对 $\mathcal S=\Aut_E|_{C_n^\circ}$，Hecke–Laumon 归纳与 Verdier 对偶保证这些假设成立：$\mathbb D(\Aut_E)$ 的特征值为 $E^*$。所以 $\Aut_E|_{C_n^\circ}$ 满足通常的 Hecke 特征性质。特别地，采用\chref{2}的 Hecke 对应约定，
\[(t_x^{d+n})^*\Aut_E^d\simeq\det(E_x)\otimes\Aut_E^{d+n},\]
其中 $t_x^{d+n}:\Bun_n^{d+n}\to\Bun_n^d$ 为同构，$M\mapsto M(-x)$。等价地，对 $\tau_x^d:\Bun_n^d\to\Bun_n^{d+n}$，$M\mapsto M(x)$，有
\[(\tau_x^d)^*\Aut_E^{d+n}\simeq\det(E_x)^{-1}\otimes\Aut_E^d.\]
\newpage
\section{延拓到整个 \texorpdfstring{$\Bun_n$}{Bunₙ} 及尖点性}\label{sec:8.6}
令 $j_d:C_n\cap\Bun_n^d\hookrightarrow\Bun_n^d$ 为次数 $d$ 分支中的良好开子叠。构造首先通过中间延拓扩展下降层：
\[\Aut_E^d:=j_{d,!*}(\Aut_E|_{C_n\cap\Bun_n^d}).\]
在洁净范围内，它与 $j_{d,!}$ 和 $j_{d,*}$ 给出同一对象。然后，上述最高次 Hecke 关系给出次数分支之间的唯一延拓，因为 $M\mapsto M(-x)$ 使次数减少 $n$。$\Hecke_n^1$ 与最高次 Hecke 对应之间的相容性，将第一 Hecke 关系从良好开子叠传播到延拓层。于是得到以 $E$ 为特征值的 Hecke 特征层
\[\Aut_E\in D(\Bun_n).\]

最后，$\Aut_E$ 是尖点层。原因是洁净延拓论证中所用的同一个消失定理：$\GL_n$ 的每个真抛物常数项函子，都经由某个 $k<n$ 的 $\Bun_k$ 上的平均操作分解，而几何不可约秩 $n$ 局部系统 $E$ 的平均定理迫使该贡献消失。因此，$\Aut_E$ 的所有真抛物常数项均消失。
\begin{transnote}{含糊点拨}
「真抛物常数项函子」是沿抛物子群丛 $\Bun_P\to\Bun_n$ 的 $!$-推前再拉回的复合，对应经典自守表示里的常数项。对 $\GL_n$，一个真抛物 $P$ 的 Levi 是 $\GL_{d_1}\times\cdots\times\GL_{d_r}$，其中必有一块 $k<n$，所以常数项可以分解成先到某个 $\Bun_k$、再作平均的形式；$\rank E=n$ 且几何不可约时，平均定理（定理 5.7.2）正给出这一项为零。这就是「$\Aut_E$ 尖点」的全部内容：不是逐项硬算，而是把每个真抛物常数项化成同一个消失定理。[依据：原书 p.47 本段；并见定理 5.7.2 与 [Gai16b]。]
\end{transnote}

在有限域上，与 FGKV 的层–函数比较陈述如下。若给定 $\Bun_n(\mathbb F_q)$ 上以 $E$ 为特征值的尖点 Hecke 特征函数，则它在 $\Bun'_n(\mathbb F_q)$ 上的限制为
\[M'\longmapsto\Tr(\operatorname{Frob}_{M'},\Aut'_E).\]
因此，上文构造的层 $\Aut_E$ 将预期自守特征函数几何化，与文献\cite{FGKV98}中 Whittaker 函数的比较相匹配。

\chapter{消失猜想}\label{ch:9}
\begingroup\setlength{\belowdisplayshortskip}{2pt}
\titlespacing*{\section}{0pt}{1.5ex plus .5ex}{1ex plus .2ex}
\srcpage{48}
本章解释\chref{6}洁净延拓论证和\secref{8.6}尖点性论证所用的平均消失定理\cite{Gai04,Gai16b}。记号与此前相同：$X$ 是亏格为 $g$ 的光滑射影曲线，$D(-)$ 表示所选层理论。以下 de Rham 讨论取 $k=\mathbb C$、$D(-)=\DMod(-)$。
\section{平均函子}\label{sec:9.1}
设 $E$ 为秩 $r$ 的不可约局部系统，$m<r$。对 $d=d_2-d_1\ge0$，考虑秩 $m$ 丛的上修改：
\origdiagram{p048-diagram.pdf}{0.3962}
其中 $\operatorname{Mod}_m^{d_1,d_2}$ 的一点为 $M\subset M'$，$\length(M'/M)=d$，且 $\pi(M\subset M')=M'/M$。令 $\cL_E^d\in D(\Tor^d)$ 为 Laumon 层：在支撑映射 $\Tor^d\to X^{(d)}$ 下，无重数轨迹 $X^{(d),\circ}\subset X^{(d)}$ 的逆像上，它是 $E^{\boxtimes d}$ 的对称下降 $E^{(d)}$；整体上，它是后者到 $\Tor^d$ 的 Goresky–MacPherson 延拓。忽略通常的归一化移位与 Tate 扭转，
\[\Av_E^d(\cF):=\overrightarrow h_!(\overleftarrow h^*\cF\otimes\pi^*\cL_E^d)\in D(\Bun_m^{d_2}).\]
等价地，对各连通分支求和后，这是 $D(\Bun_m)$ 的自函子。称其为“由 $\cL_E^d$ 作上修改”；文献\cite{Gai16b}从谱一侧回顾同一个平均函子。
\begin{theorem}{消失定理}{9.1.1}
假设 $E$ 不可约、$\rank(E)=r>m$，并为简便假设 $g>1$。则只要 $d>rm(2g-2)$，就有
\[\Av_E^d=0.\]
采用移位与扭曲的标准归一化，同样的表述适用于通常的层论情形。
\end{theorem}
\thref{9.1.1}就是此前以 $r=n$、$m=k<n$ 使用的定理。这里的重点不是构造 $\cL_E^d$，而是从几何 Langlands 的谱一侧理解这个函子为何应当消失。
\section{谱作用的启发式解释}\label{sec:9.2}\index{谱作用|hyperpage}\index{温和对象|hyperpage}
在几何 Langlands 的 de Rham 形式中，纲要上预期有文献\cite{GR24a,AG15}的谱作用及幂零支撑表述：
\[\DMod(\Bun_G)^{\mathrm{temp}}\simeq\QCoh(\LocSys_{\checkG}),\qquad\DMod(\Bun_G)\simeq\IndCoh_{\mathrm{Nilp}}(\LocSys_{\checkG}).\]
将右边限制到紧对象，这常常启发性地表现为谱叠上的凝聚层。因此，$\QCoh(\LocSys_{\checkG})$ 应当作用于 $\DMod(\Bun_G)$。
这一作用并非任意的。对 $E\in\LocSys_{\checkG}$，纤维函子
\[(-)_E:\QCoh(\LocSys_{\checkG})\longrightarrow\operatorname{Vect}\]
应通过基变换切出 $E$-Hecke 范畴：
\[\DMod(\Bun_G)\otimes_{\QCoh(\LocSys_{\checkG})}\operatorname{Vect}\simeq\DMod(\Bun_G)_{E\text{-Hecke}}.\]
\srcpage{49}
这里，右边由 Hecke 变换可通过与局部系统 $E$ 作张量积得到的对象组成。

在固定点 $x\in X$ 处，几何 Satake 给出函子
\[\Rep(\checkG)\otimes D(\Bun_G)\longrightarrow D(\Bun_G),\qquad(V,\cK)\longmapsto H_{V,x}(\cK),\]
其中 $\Rep(\checkG)$ 是球面范畴 $D(G(\cO_x)\backslash G(K_x)/G(\cO_x))$ 内的 Satake 心。若让点移动，则得到
\[\Rep(\checkG)\otimes D(\Bun_G)\longrightarrow D(\Bun_G\times X).\]
$E$-Hecke 特征层是配备同构
\[H_V(\cK)\simeq\cK\boxtimes V_E\]
的对象 $\cK$，其中同构对所有 $V\in\Rep(\checkG)$ 给出，并配备在 $X^I$ 上的高阶相容性。

朴素函子 $\Rep(\checkG)\otimes D(X)\to\End(D(\Bun_G))$ 虽然是函子，却不是所需意义下的幺半函子：当两个点碰撞时，卷积由几何 Satake 的幺半结构控制，而不是由 $X^2$ 上的外张量积控制。例如，必须有
\begin{itemize}
\item 当 $x\ne y$ 时，$H_{V,x}\circ H_{W,y}=H_{V,W;x,y}$；
\item $H_{V,W;x,x}=H_{V\otimes W,x}$。
\end{itemize}
这恰好是因子化条件，所以正确的参数空间必须是 Ran 空间。
\section{Ran 作用与局部到整体}\label{sec:9.3}\index{Ran 空间|hyperpage}\index{因子化范畴|hyperpage}
令
\[\operatorname{Ran}(X)=\{I\subset X\mid0<|I|<\infty\}.\]
它理解为各幂 $X^I$ 沿满射所形成的余极限。因子化范畴 $\Rep(\checkG)_{\mathrm{Ran}}$ 在有限集合 $I\subset X$ 上的纤维为 $\bigotimes_{x\in I}\Rep(\checkG)$，点碰撞时进行融合。Hecke 函子组合成幺半作用

\begin{transnote}{含糊点拨}
取「沿满射」的余极限，是因为 $X^I$ 只记住带标号的点，而满射 $I\twoheadrightarrow J$ 允许把若干点合并；把这些粘起来才得到「任意有限子集都能塞进去、点碰撞时发生融合」的 Ran 空间。纤维写成 $\bigotimes_{x\in I}$ 而不是 $\prod$：$\Rep(\check G)$ 是张量范畴，纤维里的对象是各点处表示的张量积，融合把两点的张量积在碰撞极限下乘成一个。缺了这层，定理 9.3.2 的幺半作用无从谈起。[来源：nLab "Ran space"，https://ncatlab.org/nlab/show/Ran+space（Last revised 2020-07-27；2026-11-17 访问）；并见 nLab "factorization algebra"（Last revised 2026-08-16）。]
\end{transnote}
\[\Rep(\checkG)_{\mathrm{Ran}}\longrightarrow\End(\DMod(\Bun_G)).\]
还有幺半“局部化”函子
\[\operatorname{Loc}_{\checkG}:\Rep(\checkG)_{\mathrm{Ran}}\longrightarrow\QCoh(\LocSys_{\checkG}),\]
它将集合 $(V_x)_{x\in I}$ 映为 $\LocSys_{\checkG}$ 上的取值向量丛 $\bigotimes_{x\in I}V_{E,x}$。

以下结果见文献\cite{Gai10}。
\begin{theorem}{Gaitsgory–Lurie，局部到整体}{9.3.1}
局部化函子 $\operatorname{Loc}_{\checkG}$ 将 $\QCoh(\LocSys_{\checkG})$ 实现为 $\Rep(\checkG)_{\mathrm{Ran}}$ 的 Verdier 商。等价地，$\operatorname{Loc}_{\checkG}$ 的右伴随全忠实。
\end{theorem}

证明在形式上来自描述
\[\LocSys_{\checkG}=\{B\checkG\times X\to X\text{ 的水平截面}\};\]
同一论证适用于任何配备相对于 $X$ 的联络的叠 $Z\to X$。因此，知道 Ran Hecke 作用后，谱作用不再需要额外选择。
\srcpage{50}
\begin{theorem}{广义消失定理\cite{Gai10,Gai04}}{9.3.2}
幺半作用
\[\Rep(\checkG)_{\mathrm{Ran}}\longrightarrow\End(\DMod(\Bun_G))\]
唯一地经由 $\QCoh(\LocSys_{\checkG})$ 分解：
\origdiagram{p050-triangle.pdf}{0.4668}
这是广义消失猜想的范畴形式。其证明使用 Kac–Moody 表示的 Beilinson–Drinfeld 局部化。
\end{theorem}
\section{恢复函子 \texorpdfstring{$\Av_E^d$}{Av（E，d）}}\label{sec:9.4}
取 $G=\GL_m$，从而 $\checkG=\GL_m$，令 $\mathrm{St}_m$ 为标准表示。$\Rep(\GL_m)_{\mathrm{Ran}}$ 中与平均核对应的对象为
\[\cA_{E,d}:=(\Sym^d(\mathrm{St}_m\otimes E))_{X^{(d)}}.\]
它支撑于自然映射 $X^{(d)}\to\operatorname{Ran}(X)$。采用与几何 Satake 对应相同的归一化，Ran Hecke 作用将 $\cA_{E,d}$ 映为 $\Av_E^d$。

原因在局部 Satake 纤维上可见。在点 $x\in X$ 上，Cauchy 分解给出
\[\Sym^d(\mathrm{St}_m\otimes E_x)\simeq\bigoplus_{\substack{\mu\vdash d\\\ell(\mu)\le\min(m,r)}}S_\mu(\mathrm{St}_m)\otimes S_\mu(E_x).\]

\begin{transnote}{含糊点拨}
$S_\mu$ 是 Schur 函子，$\mu\vdash d$ 对应 $S_d$ 的不可约表示；Cauchy 分解把 $\Sym^d(V\otimes W)$ 按对称群的不可约分量拆开，所以右端是分拆求和而不是单一张量。截断 $\ell(\mu)\le m$ 来自 $\mathrm{St}_m$ 只有 $m$ 维：$\wedge^i\mathrm{St}_m=0$ 当 $i>m$，长度更长的分拆没有贡献；$\ell(\mu)\le r$ 同理来自 $\rank E_x=r$。这正是把 $\Sym^d$ 化成有限和、可逐项计算的依据。[依据：原书 p.50 本式与 §9.4 关于 Cauchy 分解的说明。]
\end{transnote}
在几何 Satake 下，$S_\mu(\mathrm{St}_m)$ 对应于由 $\mu$ 标记的 Schubert 分层的 IC 层，而 $S_\mu(E_x)=E_x^\mu$。等价地，右边是不变空间 $(\mathrm{St}_m^{\otimes d}\otimes E_x^{\otimes d})^{S_d}$，并不是对每个直和项再取一次 $S_\mu$ 不变量。这是修改核中出现的 Laumon 层 $\cL_E^d$ 的 Satake 含义。

由\thref{9.3.2}，$\Av_E^d$ 从 $\operatorname{Loc}_{\GL_m}(\cA_{E,d})$ 得到。若 $L$ 为 $X$ 上的局部系统，记 $L^{(d)}$ 为它到 $X^{(d)}$ 的对称下降。忽略对偶约定的差异，该拟凝聚层在秩 $m$ 局部系统 $E'\in\LocSys_{\GL_m}$ 处的纤维为 $\RGamma(X^{(d)},(E'\otimes E)^{(d)})$。因此，只需证明以下初等消失。
\begin{lemma}{}{9.4.1}
设 $E$ 不可约且秩 $r>m$，$E'$ 是任意秩 $m$ 局部系统。若 $d>rm(2g-2)$，则
\[\RGamma(X^{(d)},(E'\otimes E)^{(d)})=0.\]
\end{lemma}
\begin{proof}
置 $L=E'\otimes E$。由于 $E$ 不可约且 $\rank(E)>\rank(E')$，任意映射 $(E')^\vee\to E$ 都为零；故 $H^0(X,L)=0$。对 $E^\vee$ 应用同一论证得到 $H^0(X,L^\vee)=0$，Verdier 对偶给出 $H^2(X,L)=0$。因此 $\RGamma(X,L)$ 集中在次数 $1$，且
\[\dim H^1(X,L)=(2g-2)\rank(L)=rm(2g-2).\]
Künneth 公式给出
\[\RGamma(X^{(d)},L^{(d)})\simeq\Sym^d\RGamma(X,L)\simeq\bigwedge^dH^1(X,L)[-d].\]
当 $d>\dim H^1(X,L)$ 时，最后的外幂为零。
\end{proof}

于是，谱作用将 $\cA_{E,d}$ 映为零，这就在 de Rham 情形给出了 $\Av_E^d$ 的消失定理。
\section{其他层理论与一般约化群}\label{sec:9.5}\index{受限局部系统叠|hyperpage}
\srcpage{51}
对于 de Rham $D$ 模以外的层理论，前述论证中的谱叠要替换为受限局部系统叠\cite{AGK20a}，非正式地记为 $\LocSys_{\checkG}^{\mathrm{restr}}$。可以用点函子描述它：
\[\LocSys_{\checkG}^{\mathrm{restr}}(S)=\{\Rep(\checkG)\xrightarrow{\otimes}\QCoh(S)\otimes\operatorname{QLisse}(X)\},\]
而在 de Rham 情形，将 $\operatorname{QLisse}(X)$ 替换为 $\DMod(X)$ 就有平行的表达式。在 Betti／可构造情形，已知的因子化陈述落在具有幂零奇异支撑的范畴中，纲要上为
\[\Rep(\checkG)_{\mathrm{Ran}}^{\operatorname{QLisse}}\longrightarrow\QCoh(\LocSys_{\checkG}^{\mathrm{restr}})\longrightarrow\End(D_{\mathrm{Nilp}}(\Bun_G)),\]
因此，它本身并不蕴涵整个 $D(\Bun_G)$ 上完整的消失定理。

对一般约化群，预期陈述相同。设 $V\in\Rep(\checkG)$，$E$ 为不可约向量局部系统，$\rank(E)>\dim V$。当
\[d>(2g-2)\dim(V\otimes E)\]
时，Ran 表示范畴中的对象 $\Sym^d(V\otimes E)$ 应在 $D(\Bun_G)$ 上作用为零。在 de Rham 理论中，由与上文相同的上同调计算，这从广义消失定理推出。在 $\ell$-进情形，这一表述中的相应陈述仍作为猜想性输入处理。
\section{Gaitsgory 最初的消失证明}\label{sec:9.6}
本章还概述 Gaitsgory 关于消失定理的最初证明\cite{Gai04}。论证是几何的，但对应于解析陈述：相关 Rankin–Selberg $L$ 函数是次数有界的多项式。

首先，定义下修改和函子 $\Av_{E^*}^{-d}$。忽略标准移位，$\Av_{E^*}^{-d}$ 同时是 $\Av_E^d$ 的左伴随和右伴随，Verdier 对偶给出

\begin{transnote}{含糊点拨}
下修改函子用 $E^*$ 与负次数构造，与 $\Av_E^d$ 互为伴随：忽略标准移位后它同时是左、右伴随。于是 Hom 里的 $\Av_E^d$ 可以左右挪动，$\Av_E^d=0$ 等价于 $\Hom(\Av_E^d(-),-)=0$；Verdier 对偶把紧支撑版本与 $*$-版本互换，就得到后半式 $\mathbb D\circ\Av_{E^*}^{-d}\simeq\Av_E^d\circ\mathbb D$。书中随后正是用这一条把 $F=0$ 从 Hom 计算里推出，而不是硬算核。[依据：原书 p.51 本段与随后的 Hom 计算；并见 nLab "Verdier duality" 的伴随刻画（该条目无 Last revised，2026-11-17 访问）。]
\end{transnote}
\[\mathbb D\circ\Av_{E^*}^{-d}\simeq\Av_E^d\circ\mathbb D.\]
要证明 $\Av_E^d=0$，只需证明
\[\Av_E^d\circ\Av_{E^*}^{-d}=0.\]
事实上，若 $G:=\Av_{E^*}^{-d}$ 是 $F:=\Av_E^d$ 的左伴随且 $FG=0$，则 $\Hom(G\cK,G\cK)\simeq\Hom(\cK,FG\cK)=0$，所以 $G=0$；由于 $G$ 也是 $F$ 的右伴随，便迫使 $F=0$。对消失范围内的 $d$，一个归纳引理说明 $\Av_E^d(\cF)$ 和 $\Av_{E^*}^{-d}(\cF)$ 是尖点的。因此，只要 $\Av_E^d$ 消去尖点层，所需复合就为零。

其次，证明 Euler 示性数归约所需的正合性。更精确地，首先证明 $\Av_E^d$ 在转到下文描述的辅助 Verdier 商后正合；由于其像是尖点的，且该商在尖点对象上保守，这给出消失论证所需的正合性。然后，标准的 Euler 示性数论证将消失归约为
\[\chi(\Av_E^d(\cF))=0.\]
对 Euler 示性数而言，可以将局部系统 $E$ 替换为另一个同秩局部系统；于是归约到平凡局部系统，尽管正合性输入本身使用了 $E$ 的不可约性。
\srcpage{52}

对平凡秩 $r$ 系统 $\mathbf1^{\oplus r}$，Braverman 的计算给出
\[\Av_{\mathbf1^{\oplus r}}^d\simeq\bigoplus_{d_1+\cdots+d_r=d}\Av_{\mathbf1}^{d_1}\circ\cdots\circ\Av_{\mathbf1}^{d_r}.\]
所以，只需证明：对尖点 $\cF$ 及 $d>m(2g-2)$，
\[\Av_{\mathbf1}^d(\cF)=0.\]
最后一步是直接计算。在函数层面，它是 Jacquet–Godement 的陈述：$\GL_m$ 的尖点自守表示的标准 $L$ 函数，是次数至多 $m(2g-2)$ 的多项式；Braverman 的计算将它几何化。
\section{为何平均函子的像是尖点的}\label{sec:9.7}
令 $m=m_1+m_2$，并令
\[\CT:D(\Bun_m)\longrightarrow D(\Bun_{m_1}\times\Bun_{m_2})\]
为 Levi 子群为 $\GL_{m_1}\times\GL_{m_2}$ 的标准抛物子群的常数项函子。$\CT\circ\Av_E^d$ 上有一个过滤，其相伴分次部分为
\[\bigoplus_{d_1+d_2=d}(\Av_E^{d_1}\boxtimes\Av_E^{d_2})\circ\CT.\]
这是修改叠的纯几何过滤。由于对两个 Levi 因子都有 $\rank(E)>m_i$，当 $d>rm(2g-2)$ 时，较小秩的消失定理通过归纳消去每个分次部分：对每个 $d=d_1+d_2$，至少有一个 $d_i$ 超过 $rm_i(2g-2)$。故
\[\CT\circ\Av_E^d=0,\]
所以 $\Av_E^d$ 的像是尖点的。

这也迅速证明了上文用到的一个事实。若 $\cF\in D(\Bun_m)$ 是 $E$-Hecke 特征层，则如谱作用所预言，
\[\Av_E^d(\cF)\simeq\cF\otimes\RGamma(X^{(d)},(E^*\otimes E)^{(d)}).\]
应用 $\CT$ 得到
\[\CT(\Av_E^d(\cF))\simeq\CT(\cF)\otimes\RGamma(X^{(d)},(E^*\otimes E)^{(d)}).\]
另一方面，上述过滤将同一个对象表达为作用在 $\CT(\cF)$ 上的较小秩平均函子的和，因此对 $d\gg0$ 它消失。由于 $E^*\otimes E$ 包含平凡局部系统，所示上同调对适当的大 $d$ 非零。因此 $\CT(\cF)=0$，即 $\cF$ 是尖点的。
\newpage
\section{\texorpdfstring{$t$}{t} 正合性背后的商范畴}\label{sec:9.8}\index{Verdier 商|hyperpage}\index{t正合性@$t$ 正合性|hyperpage}\index{保守函子|hyperpage}
最初证明中的困难是 $t$-正合性。看似诱人的陈述“$\Av_E^1$ 是 $t$-正合的”是错误的。正确做法是构造一个 Verdier 商
\[D(\Bun_m)\twoheadrightarrow\widetilde D(\Bun_m),\]
具有三项性质：
\origdiagram{p052-square.pdf}{0.3370}
其中 $\widetilde{\Av}_E^1$ 是 $t$-正合的，该商继承反常 $t$-结构，而且复合
\[D(\Bun_m)_{\mathrm{cusp}}\longrightarrow D(\Bun_m)\longrightarrow\widetilde D(\Bun_m)\]
\srcpage{53}
是保守的。这是对 $\Av_E^1$ 那个错误正合性陈述的 Springer 型替代。Gaitsgory 最初构造 $\widetilde D$ 的方法在技术上相当复杂。

现代的解释是应将这个商取为温和部分：
\[\widetilde D(\Bun_G)=D(\Bun_G)^{\mathrm{temp}}.\]
几何 Ramanujan 定理说明，尖点对象是温和的\cite{Ber21}；Eisenstein／常数项的形式体系补上其余的非尖点部分。这为旧的商范畴提供了概念上更清楚的替代。

\begin{transnote}{背景直觉}
几何 Ramanujan 定理是经典 Ramanujan 猜想的几何版本：尖点自守表示的「特征值」都落在温度带内，等价地说，尖点对象在谱侧是温和的。经典对应是 $\GL_2$ 上尖点自守表示的 $L$ 函数在临界带上无极点（Ramanujan 界）；这里把它翻译成「尖点层属于 $\Shv_{\mathrm{Nilp}}$ 的温和子范畴」。它补上的正是旧的商范畴（尖点部分）缺掉的那一半：非尖点部分交给 Eisenstein／常数项形式体系。[依据：原书 p.53 本段；并见 [Ber21]。]
\end{transnote}
\endgroup

\chapter{范畴猜想}\label{ch:10}
\begingroup
\setlength{\abovedisplayskip}{3pt}\setlength{\belowdisplayskip}{3pt}
\setlength{\belowdisplayshortskip}{3pt}
\setlength{\intextsep}{8pt}
% 原PDF第53页，第10章开始；第9章结尾不在本文件中重复。
现在，我们从单个 Hecke 特征层转向几何 Langlands 的范畴形式。本章取 $k=\mathbb C$，系数域也取 $\mathbb C$。除非另有说明，均在 de Rham 情形下讨论。记 $\LocSys_{\checkG}$ 为 $X$ 上 de Rham $\checkG$-局部系统的导出叠。忽略惯常的半扭转，预期的等价采用 Arinkin–Gaitsgory／Gaitsgory–Raskin 的表述\cite{AG15,GR24a}：

\begin{transnote}{含糊点拨}
「半扭转」指 Hecke 特征层与 Eisenstein 函子上各差一个 Tate 扭转 $(\text{--})(1/2)$：几何 Satake 把 $V$ 送到 $\mathrm{IC}_V$ 时，$\ell$-进归一化要乘上权 $q^{1/2}$ 的那一半，两侧的权因此差半格。Arinkin–Gaitsgory 与 Gaitsgory–Raskin 的两种表述差的正是这个半整扭转，忽略它之后两个函子才能逐项对上；后面 §10.4、§10.5 的归一化都以此为基准。[来源：nLab "Tate twist"，https://ncatlab.org/nlab/show/Tate+twist（Last revised 2025-08-10；2026-11-17 访问）；并见原书 p.53 本段与 [AG15]、[GR24a]。]
\end{transnote}
\begin{conjecture*}{}
\[\DMod(\Bun_G)\simeq\IndCoh_{\mathrm{Nilp}}(\LocSys_{\checkG}).\]
\end{conjecture*}
右端之所以采用这个范畴，是因为 $\LocSys_{\checkG}$ 作为导出叠通常是奇异的。因此，$\QCoh(\LocSys_{\checkG})$、$\IndCoh(\LocSys_{\checkG})$，以及通过奇异支撑条件得到的中间范畴 $\IndCoh_{\mathrm{Nilp}}(\LocSys_{\checkG})$，确实是不同的范畴。

\section{为何出现 ind-凝聚层}\label{sec:10.1}\index{ind-凝聚层|hyperpage}\index{完美复形|hyperpage}\index{奇异支撑!谱侧|hyperpage}
对于奇异导出概形或导出叠 $Y$，完美对象与凝聚对象未必相同；一般而言，只存在包含关系
\[\operatorname{Perf}(Y)\subseteq\Coh(Y),\]
而且这个包含往往是严格的。
\begin{example}
取 $A=k[\varepsilon]$，其中 $\deg(\varepsilon)=-1$，并令 $Y=\Spec A$。则 $k\in\Coh(Y)$，但 $k\notin\operatorname{Perf}(Y)$；从下面的无限周期解消可以看出这一点：
\[\cdots\xrightarrow{\varepsilon}A\xrightarrow{\varepsilon}A\xrightarrow{\varepsilon}A\longrightarrow k.\]
\end{example}
进行 ind-完备化后，对满足适当条件的 $Y$，有
\[\QCoh(Y)=\operatorname{Ind}(\operatorname{Perf}(Y))\subseteq\IndCoh(Y)=\operatorname{Ind}(\Coh(Y)),\]
在上述奇异例子中，包含仍是严格的。因此，$\IndCoh$ 保留了 $\QCoh$ 可能遗漏的凝聚对象。

不过，在 $\operatorname{Perf}$ 与 $\Coh$ 之间，存在一族可以控制的范畴。依照 Arinkin–Gaitsgory\cite{AG15}，若 $Y$ 是拟光滑的，定义它的奇异概形为
\begin{definition*}{}
\[\operatorname{Sing}(Y):=\Spec_{Y^{\mathrm{cl}}}\Sym_{\cO_{Y^{\mathrm{cl}}}}\bigl(H^1(T_Y)\bigr).\]
等价地，它是 $(-1)$ 次移位余切叠的底层经典叠。若 $Y$ 光滑，则 $\operatorname{Sing}(Y)$ 是 $Y$ 上的零向量丛，也就是零截面。对 $\cF\in\Coh(Y)$，可以赋予一个锥形支撑 $\supp(\cF)\subset\operatorname{Sing}(Y)$。对锥形闭子集 $N\subset\operatorname{Sing}(Y)$，置

\begin{transnote}{易错预警}
$\operatorname{Sing}(Y)$ 不是 $Y$ 的奇点集，而是「$(-1)$ 次移位余切叠的底层经典叠」：$\operatorname{Sing}(Y)=\Spec_{Y^{\mathrm{cl}}}\Sym(H^1(T_Y))$，把 $H^1(T_Y)$ 当作余切方向、并且取了负一次移位。$Y$ 光滑时 $H^1(T_Y)=0$，只剩零截面，这才是「光滑时 $\operatorname{Sing}(Y)$ 是零向量丛」的意思。用 $H^1$ 而不是 $H^0$，正是准光滑导出概形与通常余切丛的差别所在。[依据：原书 p.53 本段；并见 [AG15] 对 $\operatorname{Sing}(Y)$ 的定义。]
\end{transnote}
\[\Coh_N(Y):=\{\cF\in\Coh(Y)\mid\supp(\cF)\subset N\},\qquad
\IndCoh_N(Y):=\operatorname{Ind}(\Coh_N(Y)).\]
\end{definition*}
零支撑范畴恢复完美对象：
\[\Coh_0(Y)\simeq\operatorname{Perf}(Y),\qquad\IndCoh_0(Y)\simeq\QCoh(Y).\]
\srcpage{54}
\Needspace{6\baselineskip}
\begin{example}
一个有用的局部模型是导出交 $Y=\mathrm{pt}\times_V\mathrm{pt}$，其中两个到向量空间 $V$ 的映射都是零映射。此时 $H^1(T_Y)=V$，所以 $\operatorname{Sing}(Y)=V^*$。Koszul 对偶给出
\[\IndCoh(\mathrm{pt}\times_V\mathrm{pt})\simeq\Sym(V[-2])\text{-}\mathrm{Mod}.\]
因此，$H^\bullet(\cF)$ 是一个分次 $\Sym(V)$-模，等价地，是 $V^*$ 上的拟凝聚层；$\supp(\cF)$ 就是这个模通常意义下的锥形支撑。
\end{example}
对于一般概形，形式邻域 $\widehat Y_y$ 的一阶行为由 $T_yY$ 上的导出自交控制；这给出同样的支撑。对于叠，则通过下降将这一构造粘合起来。

\section{谱侧的幂零奇异支撑}\label{sec:10.2}\index{幂零奇异支撑|hyperpage}
令 $(E,\nabla)\in\LocSys_{\checkG}$。伴随丛 $\check{\mathfrak g}_E$ 带有诱导的联络。切复形为
\[T_{(E,\nabla)}\LocSys_{\checkG}\simeq\RGamma_{\mathrm{dR}}(X,\check{\mathfrak g}_E)[1]
\simeq\RGamma\bigl(X,\check{\mathfrak g}_E\xrightarrow{\operatorname{ad}(\nabla)}\check{\mathfrak g}_E\otimes\Omega_X\bigr)[1].\]
$\operatorname{Sing}(\LocSys_{\checkG})$ 在 $(E,\nabla)$ 上的纤维是 $H^1(T_{(E,\nabla)}\LocSys_{\checkG})^\vee$。由于 $H^1(T)=H^2_{\mathrm{dR}}(X,\check{\mathfrak g}_E)$，Poincaré 对偶将这个对偶空间识别为 $H^0_{\mathrm{dR}}(X,\check{\mathfrak g}_E^*)$。因而，经典奇异叠的点为
\[\operatorname{Sing}(\LocSys_{\checkG})=
\{(E,\nabla,A)\mid A\in H^0_{\mathrm{dR}}(X,\check{\mathfrak g}_E^*)\}.\]
等价地，$A$ 是余伴随局部系统的水平截面。利用不变双线性形式 $\check{\mathfrak g}^*\simeq\check{\mathfrak g}$，定义
\begin{definition*}{}
\[\mathrm{Nilp}:=\{(E,\nabla,A)\mid A_x\in\check{\mathfrak g}\text{ 对每个 }x\in X\text{ 都是幂零的}\}
\subset\operatorname{Sing}(\LocSys_{\checkG}).\]
\end{definition*}
因此，$\IndCoh_{\mathrm{Nilp}}(\LocSys_{\checkG})$ 由奇异方向为全局幂零水平场的 ind-凝聚层组成。
\begin{example}

即使最简单的曲线也能说明导出结构为何重要。对于 $X=\mathbb P^1$，非导出的 de Rham 局部系统叠本质上是 $\mathrm{pt}/\checkG$，但导出叠则是导出零纤维
\[\LocSys_{\checkG}(\mathbb P^1)\simeq\mathrm{pt}/\checkG\times_{\check{\mathfrak g}/\checkG}\mathrm{pt}/\checkG.\]
若只保留 $\QCoh(\mathrm{pt}/\checkG)\simeq\Rep(\checkG)$，就会失去在自守侧可见的导出 $\Ext$ 数据，例如球面范畴
\[\DMod(G(\cO)\backslash G(\mathcal K)/G(\cO))\]
中的数据。导出自交上的范畴 $\IndCoh_{\mathrm{Nilp}}$ 才是正确的细化。
\end{example}

\section{相容性检验一：导出 Satake}\label{sec:10.3}\index{导出 Satake 等价|hyperpage}\index{谱 Hecke 叠|hyperpage}
与 Arinkin–Gaitsgory 的奇异支撑形式体系\cite{AG15}相容的局部球面范畴为
\[\operatorname{Sph}_{G,x}:=\DMod(G(\cO_x)\backslash G(\mathcal K_x)/G(\cO_x)).\]
通常的几何 Satake 等价给出心范畴 $\Rep(\checkG)$。在导出范畴层面，正确的谱对象是谱 Hecke 叠
\[\Hecke^{\mathrm{spec}}_{\checkG}:=\mathrm{pt}/\checkG\times_{\check{\mathfrak g}/\checkG}\mathrm{pt}/\checkG
\simeq(\mathrm{pt}\times_{\check{\mathfrak g}}\mathrm{pt})/\checkG.\]
\begin{theorem*}{导出 Satake}
\[\operatorname{Sph}_{G,x}\simeq\IndCoh_{\mathrm{Nilp}(\check{\mathfrak g}^*)/\checkG}(\Hecke^{\mathrm{spec}}_{\checkG}).\]

\begin{transnote}{背景直觉}
导出 Satake 把几何 Satake 的心范畴等价提升到导出层：左端是球面 Hecke 范畴（$\operatorname{Sph}_{G,x}$），右端是谱 Hecke 叠 $\Hecke^{\mathrm{spec}}_{\check G}$ 上「奇异支撑落在幂零锥内」的 ind-凝聚层。下标 $\mathrm{Nilp}(\check{\mathfrak g}^*)/\check G$ 就是允许的奇异支撑范围——只有支撑落在幂零锥的部分才对应球面对象的谱侧。这解释了为什么范畴猜想里必须用 $\IndCoh_{\mathrm{Nilp}}$ 而不能用普通 $\QCoh$：$\QCoh$ 看不见奇异支撑。[来源：nLab "geometric Satake equivalence"，https://ncatlab.org/nlab/show/geometric+Satake+equivalence（Last revised 2024-07-22；2026-11-17 访问）；并见 [AG15]。]
\end{transnote}
\end{theorem*}
这里，$\operatorname{Sing}(\Hecke^{\mathrm{spec}}_{\checkG})\simeq\check{\mathfrak g}^*/\checkG$，而幂零锥 $\mathrm{Nilp}(\check{\mathfrak g}^*)/\checkG$ 恰好是允许的奇异支撑。这是全局范畴 $\IndCoh_{\mathrm{Nilp}}(\LocSys_{\checkG})$ 的局部模型。

\section{相容性检验二：Eisenstein 级数}\label{sec:10.4}\index{Eisenstein 函子!谱侧|hyperpage}
\srcpage{55}
令 $P\subset G$ 为抛物子群，Levi 商为 $M$。在自守侧，有
\origdiagram{p055-automorphic.pdf}{0.3734}
以及 Eisenstein 函子
\[\operatorname{Eis}_{P,!}:=p_!q^*:\DMod(\Bun_M)\longrightarrow\DMod(\Bun_G).\]
它的连续右伴随是常数项函子
\[\CT_{P,*}:=q_*p^!:\DMod(\Bun_G)\longrightarrow\DMod(\Bun_M).\]
因此，自守侧的形式体系在取伴随时是稳定的。

另一方面，谱侧的图表为
\origdiagram{p055-spectral.pdf}{0.4622}
朴素的函子
\[\check p_*\check q^*:\QCoh(\LocSys_{\check M})\longrightarrow\QCoh(\LocSys_{\checkG})\]
没有所需的连续右伴随。障碍在于，$\check p$ 的真性足以保证它保持凝聚层，但它并不保持完美复形。恰当的修正正是采用幂零 ind-凝聚范畴。示意地，谱 Eisenstein 函子为
\[\operatorname{Eis}^{\mathrm{spec}}_{\check P}:=(\check p)_*^{\IndCoh}\circ\check q^!,\]
而且它将 Levi 侧的完美对象送入具有幂零奇异支撑的范畴：
\[\operatorname{Eis}^{\mathrm{spec}}_{\check P}(\operatorname{Perf}(\LocSys_{\check M}))\subset\Coh_{\mathrm{Nilp}}(\LocSys_{\checkG}).\]
此外，$\IndCoh_{\mathrm{Nilp}}(\LocSys_{\checkG})$ 由所有 Levi 子群 $\check M$ 对应的对象 $\operatorname{Eis}^{\mathrm{spec}}_{\check P}(\operatorname{Perf}(\LocSys_{\check M}))$ 生成，其中也包括 $\check M=\checkG$。这说明幂零条件不能任意舍弃：Eisenstein 相容性迫使我们采用这一条件。

\section{温和对象与 Whittaker 检测}\label{sec:10.5}\index{温和对象|hyperpage}\index{Whittaker 系数!检测|hyperpage}
还有一个范围较小的预期等价
\begin{conjecture*}{}
\[\DMod(\Bun_G)^{\mathrm{temp}}\simeq\QCoh(\LocSys_{\checkG}).\]
\end{conjecture*}
从基本思想上说，温和部分是 $\DMod(\Bun_G)$ 中能够被 Whittaker 系数检测到的部分。对一般位置定义的 $N$-约化，有相应的 Whittaker 范畴和一个自然的 $!$-推前
\[\operatorname{Wh}(\Bun_G^{N\text{-}\mathrm{gen}})\longrightarrow\DMod(\Bun_G).\]
所需的局部输入是：仿射 Grassmann 簇的 Whittaker 范畴的手征同调，与 $\Rep(\checkG)$ 的手征同调相匹配。这提供了生成对象，但单凭这一点，并不能识别定义温和子范畴的 Verdier 商。实际使用的现代定义采用球面 Hecke 作用。
\srcpage{56}

对每个 $x\in X$，范畴 $\operatorname{Sph}_{G,x}$ 以幺半方式作用于 $\DMod(\Bun_G)$。在导出 Satake 等价下，
\[\operatorname{Sph}_{G,x}\simeq\IndCoh_{\mathrm{Nilp}}(\mathrm{pt}/\checkG\times_{\check{\mathfrak g}/\checkG}\mathrm{pt}/\checkG),\]
而温和的，也称为左完备化的谱 Hecke 范畴是
\[\QCoh(\mathrm{pt}/\checkG\times_{\check{\mathfrak g}/\checkG}\mathrm{pt}/\checkG).\]
\begin{definition*}{}
若由 $\cF\in\DMod(\Bun_G)$ 生成的 $\operatorname{Sph}_{G,x}$-模经由上述左完备化范畴分解，则称 $\cF$ 是温和的。
\end{definition*}
\begin{theorem*}{Færgeman–Raskin\cite{FR23}}
由此得到的子范畴与所选的点 $x$ 无关。
\end{theorem*}

整个自守范畴由温和紧对象的 Eisenstein 级数生成：
\[\DMod(\Bun_G)=\bigl\langle\operatorname{Eis}_P(\DMod(\Bun_M)^{\mathrm{temp},c})\mid P\subset G\bigr\rangle.\]
因此，温和猜想连同上述 Eisenstein 相容性，等价于完整的范畴猜想。

\newpage
\section{Betti 理论与其他层理论}\label{sec:10.6}\index{奇异支撑!微局部|hyperpage}\index{全局幂零锥|hyperpage}\index{Higgs 叠|hyperpage}
对于其他层理论，特别是 Betti 层或 $\ell$-进层，自守范畴本身也必须受到奇异支撑的限制\cite{BZN16,NY19,AGK20a}。在 Betti 情形，$\LocSys_{\checkG}$ 是 $\pi_1(X)$ 到 $\checkG$ 的表示所组成的导出特征叠；在最精确的层论表述中，还可以将它替换为相应的受限变体。猜想的形式为
\begin{conjecture*}{}
\[\Shv_{\mathrm{Nilp}}(\Bun_G)\simeq\IndCoh_{\mathrm{Nilp}}(\LocSys_{\checkG}).\]
\end{conjecture*}
这里，$\Shv_{\mathrm{Nilp}}(\Bun_G)$ 用微局部奇异支撑定义。对层 $\cF\in\Shv(Y)$，它的奇异支撑是闭锥形子集 $\operatorname{SS}(\cF)\subset T^*Y$，通过微局部分析定义；等价地，可以用非特征方向上邻近圈的消失来描述。一个层是局部系统，当且仅当 $\operatorname{SS}(\cF)$ 包含于零截面。

\begin{transnote}{含糊点拨}
两种定义等价，但方向要分清：微局部定义说 $\mathrm{SS}(\cF)$ 由「微茎非零的余方向」组成；非特征方向 $v$ 上的邻近圈消失，则说沿 $v$ 没有微局部信息，即 $\cF$ 沿 $v$ 局部常值。两者配对给出同一锥。「局部系统 $\Leftrightarrow$ $\mathrm{SS}\subset$ 零截面」是同一说法的极端情形：只有零余方向允许非零微茎，其余方向全消失。在 de Rham 情形 $\mathrm{SS}$ 可与 $D$ 模的特征簇比较，Betti 情形则与微支撑直接对应。[来源：nLab "singular support"（该条目无 Last revised，2026-11-17 访问）；并见原书 p.56 本段与 [BZN16]。]
\end{transnote}

在 $G$-丛 $P$ 处，
\[T_P\Bun_G\simeq\RGamma(X,\mathfrak g_P)[1],\qquad
T_P^*\Bun_G\simeq\RGamma(X,\mathfrak g_P^*\otimes\Omega_X).\]
经典余切叠就是 Higgs 叠：
\[T^*\Bun_G=\operatorname{Higgs}_G,\qquad
\operatorname{Higgs}_G(\mathbb C)=\{(P,\varphi)\mid\varphi\in H^0(X,\mathfrak g_P^*\otimes\Omega_X)\}.\]
\begin{definition*}{}
全局幂零锥是
\[\mathrm{Nilp}_{\mathrm{glob}}:=\{(P,\varphi)\in\operatorname{Higgs}_G\mid
\varphi_x\in\mathcal N_{\mathfrak g^*}\text{ 对每个 }x\in X\text{ 都成立}\}.\]
置
\[\Shv_{\mathrm{Nilp}}(\Bun_G):=\{\cF\in\Shv(\Bun_G)\mid\operatorname{SS}(\cF)\subset\mathrm{Nilp}_{\mathrm{glob}}\}.\]
\end{definition*}
这就是在谱侧的 ind-凝聚范畴上施加幂零奇异支撑条件的 Betti 对应版本。
\endgroup

\chapter{范畴 \texorpdfstring{$\Shv_{\mathrm{Nilp}}(\Bun_G)$}{ShvNilp(BunG)}}\label{ch:11}
\srcpage{57}
\chref{10}表述了 de Rham 范畴猜想\cite{AG15,GR24a}
\[\DMod(\Bun_G)\simeq\IndCoh_{\mathrm{Nilp}}(\LocSys^{\mathrm{dR}}_{\checkG}).\]
本章讨论其他层理论中的对应版本，特别是\cite{BZN16,NY19}中的 Betti 表述。在 Betti 情形，我们使用复解析叠 $\Bun_G(\mathbb C)$ 上的普通层。正确的自守范畴并非所有层，而是具有幂零微局部奇异支撑的满子范畴。

\section{代数叠上的所有层}\label{sec:11.1}\index{层!代数叠上的所有层|hyperpage}\index{光滑下降|hyperpage}
对 $\mathbb C$ 上有限型仿射概形 $S$，置
\[\Shv^{\mathrm{all}}(S):=\Shv^{\mathrm{all}}(S(\mathbb C)),\]
不加任何可构造性条件。若 $Y$ 为代数叠，则通过光滑下降定义 $Y$ 上的层：
\begin{definition*}{}
\[\Shv^{\mathrm{all}}(Y):=\lim_{(S\to Y)\in(\operatorname{Aff}^{\mathrm{sm}}_{/Y})^{\mathrm{op}}}\Shv^{\mathrm{all}}(S),\]
其中过渡函子为拉回。
\end{definition*}
等价地，利用伴随关系，也可以将它写成余极限
\[\Shv^{\mathrm{all}}(Y)\simeq\underset{\substack{S\to Y\\\text{光滑}}}{\operatorname{colim}}\Shv^{\mathrm{all}}(S),\]
其中过渡函子为相应的左伴随。一般的范畴原理如下：若 $L_{ij}:\mathcal C_i\rightleftarrows\mathcal C_j:R_{ij}$ 是可呈示范畴之间的一对伴随函子，则

\begin{transnote}{含糊点拨}
这条原理要求范畴可呈示：有紧生成元、滤余极限够多，否则按左伴随取的余极限与按右伴随取的极限未必对得上。它把「光滑下降给出的极限」与「按左伴随写出的余极限」等同起来，正是下面把 $\Shv^{\mathrm{all}}(Y)$ 写成余极限的依据；没有它，两种写法只是形式记号。[依据：原书 p.57 本段；并见 nLab "locally presentable category"（该条目无 Last revised，2026-11-17 访问）。]
\end{transnote}
\[\underset{L_{ij}}{\operatorname{colim}}\mathcal C_i\simeq\lim_{R_{ij}}\mathcal C_i.\]
因此，$\Shv^{\mathrm{all}}(\Bun_G)$ 是一个由 $\Bun_G$ 的解析图册构成的大可呈示范畴。

\section{Betti 局部系统与 de Rham 局部系统}\label{sec:11.2}\index{局部系统!Betti 与 de Rham 比较|hyperpage}\index{Riemann–Hilbert 对应|hyperpage}\index{刚性化|hyperpage}
在 de Rham 情形，谱参数是 de Rham 局部系统：
\[\begin{split}
\LocSys^{\mathrm{dR}}_{\checkG}(X)&=\operatorname{Conn}_{\checkG}(X)\\
&:=\{(E_{\checkG},\nabla)\mid E_{\checkG}\text{ 为 }\checkG\text{-挠子，}\nabla\text{ 为联络}\}.
\end{split}\]
在 Betti 情形，谱参数是 Betti 局部系统：
\[\LocSys^{\mathrm{Betti}}_{\checkG}(X):=\operatorname{Map}(X^{\mathrm{top}},B\checkG).\]
选定基点 $x\in X(\mathbb C)$ 后，刚性化 Betti 模空间的经典截断为
\[\LocSys^{\mathrm{Betti},\mathrm{rig},\mathrm{cl}}_{\checkG,x}(X)\simeq\Hom(\pi_1(X(\mathbb C),x),\checkG).\]
未刚性化的经典叠通过取共轭作用的商得到。这个描述只是第一近似：范畴 Langlands 中使用的 Betti 模空间是导出的，它的切复形由 $X(\mathbb C)$ 上以伴随局部系统为系数的上链控制。

这两个模叠在代数意义上并不相同。de Rham 叠依赖于复代数曲线 $X$，而 Betti 叠只依赖于 $X(\mathbb C)$ 的拓扑。
\begin{example}
令 $X=E$ 为椭圆曲线，并选取 $e\in E$。则
\[\LocSys^{\mathrm{dR},\mathrm{rig}}_{\Gm,e}(E)\longrightarrow\Pic_E^0\simeq E^\vee\]
是普遍向量扩张，因此是 $E^\vee$ 上的 $\Ga$-挠子，等价地，是一个 $\mathbb A^1$-丛。相比之下，
\[\LocSys^{\mathrm{Betti},\mathrm{rig}}_{\Gm,e}(E)\simeq\Hom(\mathbb Z^2,\Gm)\simeq\Gm^2.\]
\srcpage{58}
Riemann–Hilbert 对应在解析意义下识别它们的复点，但这一对应使用指数映射，并不识别代数结构。
\end{example}

\section{为何朴素的 Betti 猜想过大}\label{sec:11.3}\index{拟光滑层|hyperpage}\index{Fourier–Mellin 变换|hyperpage}
最初可能猜想
\[\Shv^{\mathrm{all}}(\Bun_G)\overset{?}{\simeq}\IndCoh_{\mathrm{Nilp}}(\LocSys^{\mathrm{Betti}}_{\checkG}).\]
\begin{example}
即使 $G=\GL_1$、$X=E$ 为椭圆曲线，这个猜想也是错误的。在自守侧，
\[\Bun_{\Gm}\simeq\coprod_{d\in\mathbb Z}\Pic_E^d\times B\Gm\simeq E\times\mathbb Z\times B\Gm.\]
拓扑 Fourier–Mellin 变换将有限性或可构造性的归一化识别为
\[\Shv(\mathbb Z)_{\mathrm{fin}}\simeq\QCoh(B\Gm),\qquad
\Shv(B\Gm)_{\mathrm{fin}}\simeq\QCoh(\mathrm{pt}\times_{\mathbb A^1}\mathrm{pt}).\]
若按字面取所有层，则这些范畴需要完备化；在这个例子中，完备化问题并不是两侧不匹配的原因。秩一局部系统的导出 Betti 叠具有相应的因子：
\[\LocSys^{\mathrm{Betti}}_{\Gm}(E)\simeq\Gm^2\times B\Gm\times(\mathrm{pt}\times_{\mathbb A^1}\mathrm{pt}).\]
因此，唯一不匹配的是椭圆曲线因子：
\[\Shv^{\mathrm{all}}(E)\ne\QCoh(\Gm^2).\]
正确的替代是光滑层组成的满子范畴
\[\begin{split}
\operatorname{QLisse}(E):=\{\cF\in\Shv^{\mathrm{all}}(E)\mid{}
& H^i(\cF)\text{ 是由有限维局部系统}\\[-2pt]
&\text{生成的 ind-对象}\}.
\end{split}\]
等价地，
\[\operatorname{QLisse}(E)=\Shv^{\mathrm{all}}_0(E):=\{\cF\mid\operatorname{SS}(\cF)\subset T_E^*E\}.\]
这里 $T_E^*E$ 是零截面。因此，朴素猜想的错误恰好在于包含了具有非零微局部支撑的层。
\end{example}

\section{层的微局部奇异支撑}\label{sec:11.4}\index{奇异支撑!微局部|hyperpage}\index{特征簇|hyperpage}\index{拟光滑层|hyperpage}
令 $S$ 为光滑复概形。对可构造层 $\cF$，微局部分析赋予一个闭锥形子集\cite{Bei16}
\[\operatorname{SS}(\cF)\subset T^*S,\]
称为奇异支撑。它度量 $\cF$ 不能局部常值的程度。更准确地说，若切方向 $v\in T_sS$ 被 $\operatorname{SS}(\cF)_s$ 中的所有余向量消去，则 $\cF$ 沿这一方向局部无圈，或者说是光滑的。对闭锥形子集 $N\subset T^*S$，通过 ind-完备化定义所有层的版本：
\begin{definition*}{}
\[\Shv^{\mathrm{all}}_N(S):=\operatorname{Ind}\{\cF\in\Shv_{\mathrm{cons}}(S)\mid\operatorname{SS}(\cF)\subset N\}.\]
\end{definition*}
等价地，也可以将它看作 $\Shv^{\mathrm{all}}(S)$ 中由奇异支撑位于 $N$ 的可构造对象生成的余完备满子范畴。零支撑范畴就是光滑层范畴，也称拟光滑层范畴：
\[\Shv^{\mathrm{all}}_0(S)\simeq\operatorname{QLisse}(S).\]
在 AGKRRV 的归一化中，这指的是有限秩局部系统的 ind-范畴的左完备化。

\begin{transnote}{易错预警}
「左完备化」是相对 $t$-结构取的（Lurie SAG 意义），不是 $\operatorname{Ind}$-完备化：它只把沿 $t$-结构正向收敛的对象补进来，补完恰是零奇异支撑层的完备化，即拟光滑（quasi-lisse）层。若按普通 $\operatorname{Ind}$ 去补，会把不满足收敛条件的对象一并放进来，微局部意义就丢了。[来源：nLab "t-structure"，https://ncatlab.org/nlab/show/t-structure（Last revised 2023-04-20；2026-11-17 访问）；并见原书 p.58 本段。]
\end{transnote}

为了与 de Rham 侧比较，若 $M$ 为正则完整 $D$ 模，而 $\operatorname{Sol}(M)$ 是它的解层，则
\[\operatorname{SS}(\operatorname{Sol}(M))=\operatorname{Char}(M).\]
\srcpage{59}
因此，Betti 层的微局部支撑是完整 $D$ 模的特征簇的 Betti 对应物。

\section{全局幂零锥}\label{sec:11.5}\index{全局幂零锥|hyperpage}\index{Hitchin 纤维化|hyperpage}
对 $G$-丛 $P$，记 $\mathfrak g_P=P\times^G\mathfrak g$。$\Bun_G$ 的经典余切叠是 Higgs 叠
\[T^*\Bun_G\simeq\operatorname{Higgs}_G,\qquad
\operatorname{Higgs}_G(\mathbb C)=\{(P,\varphi)\mid\varphi\in H^0(X,\mathfrak g_P^*\otimes\Omega_X)\}.\]
\begin{definition*}{}
\[\mathrm{Nilp}_{\mathrm{glob}}:=\{(P,\varphi)\in T^*\Bun_G\mid
\varphi_x\in\mathcal N_{\mathfrak g^*}\text{ 对每个 }x\in X\text{ 都成立}\}.\]
\end{definition*}
它是通常的幂零锥 $\mathcal N_{\mathfrak g^*}\subset\mathfrak g^*$ 的全局对应物。

等价地，$\mathrm{Nilp}_{\mathrm{glob}}$ 是 Hitchin 纤维化的零纤维。若 $p_i\in\mathbb C[\mathfrak g^*]^G$ 是次数为 $d_i$ 的齐次生成元，置
\[\cA_G:=\bigoplus_i H^0(X,\Omega_X^{d_i}).\]
Hitchin 映射为
\[h:T^*\Bun_G\longrightarrow\cA_G,\qquad(P,\varphi)\longmapsto(p_i(\varphi))_i.\]
于是
\[\mathrm{Nilp}_{\mathrm{glob}}=h^{-1}(0).\]
用不变记号表示，这就是映射
\[\Gamma(X,\mathfrak g_P^*\otimes\Omega_X)\longrightarrow\Gamma(X,(\mathfrak g^*\!/\!/ G)_{\Omega_X}),\]

\begin{transnote}{含糊点拨}
扭转 $(\text{--})_{\Omega_X}$ 把权重为 $d_i$ 的不变量与 $\Omega_X^{\otimes d_i}$ 张量：$\mathbb C[\mathfrak g^*]^{\check G}$ 的齐次生成元 $p_i$ 次数为 $d_i$，所以不变量截面是 $\bigoplus_i H^0(X,\Omega_X^{d_i})$，而不是朴素的 $H^0(X,(\mathfrak g^*\!/\!/\check G)\otimes\Omega_X)$。必须逐生成元各自扭转，因为各生成元权不同，不能先取商再统一扭。这个直和正是 Hitchin 基 $\cA_{\check G}$。[依据：原书 p.59 本段；并见\chref{14}开头对 Hitchin 底空间的直和描述。]
\end{transnote}
其中 $(\text{--})_{\Omega_X}$ 表示由 $\mathfrak g^*$ 上的伸缩作用产生的 $\Gm$-扭转。取次数为 $d_i$ 的齐次生成元 $p_i$ 后，该目标为 $\bigoplus_iH^0(X,\Omega_X^{d_i})$。

\section{具有幂零支撑的 Betti 范畴猜想}\label{sec:11.6}\index{Betti 几何 Langlands|hyperpage}\index{紧生成|hyperpage}\index{Hecke 光滑性|hyperpage}
修正后的 Betti 猜想是 Betti 几何 Langlands 的幂零支撑版本\cite{BZN16,NY19}：
\begin{conjecture*}{}
\[\Shv^{\mathrm{all}}_{\mathrm{Nilp}}(\Bun_G)\simeq\IndCoh_{\mathrm{Nilp}}(\LocSys^{\mathrm{Betti}}_{\checkG}),\]
其中
\[\Shv^{\mathrm{all}}_{\mathrm{Nilp}}(\Bun_G):=\{\cF\in\Shv^{\mathrm{all}}(\Bun_G)\mid\operatorname{SS}(\cF)\subset\mathrm{Nilp}_{\mathrm{glob}}\}.\]
\end{conjecture*}
在适用于一般层理论，特别是 $\ell$-进层的表述中，谱叠被替换为受限变体 $\LocSys^{\mathrm{restr}}_{\checkG}$（定义见 \secref{12.5}），但自守范畴仍然是 $\Bun_G$ 上层范畴的幂零支撑子范畴。

这个条件通过三项基本检验。第一，它对环面成立：由于 $\mathfrak t^*$ 中的幂零锥是零截面，$\Shv^{\mathrm{all}}_{\mathrm{Nilp}}(\Bun_T)$ 就是光滑部分，而 Mellin 变换将它识别为 Betti 对偶环面局部系统叠上的拟凝聚层。第二，它不会过小：Hecke 特征层应当具有幂零奇异支撑。事实上，Nadler–Yun 证明，对具有幂零奇异支撑的层，Hecke 函子在 $X$ 的点这一参数上局部常值\cite{NY19}；结合 AGKRRV 中使用的逆命题和推广\cite{AGK20a}，Hecke 光滑行为等价于幂零奇异支撑。特别地，特征层属于 $\Shv^{\mathrm{all}}_{\mathrm{Nilp}}(\Bun_G)$。第三，它不会过大：$\Shv^{\mathrm{all}}_{\mathrm{Nilp}}(\Bun_G)$ 是紧生成的，而 $\Shv^{\mathrm{all}}(\Bun_G)$ 一般不是。

更一般地，对光滑概形 $S$ 和锥形子集 $N\subset T^*S$，范畴 $\Shv^{\mathrm{all}}_N(S)$ 受一个适配于 $N$ 的分层控制：奇异支撑位于 $N$ 的对象沿各个分层局部常值。这就是幂零支撑恢复有限性性质的微局部原因。

\chapter{谱分解与 \texorpdfstring{$\LocSys^{\mathrm{restr}}_{\checkG}$}{LocSys 的受限变体}}\label{ch:12}
\srcpage{60}
上一章介绍了范畴 Langlands 的 Betti 形式\cite{BZN16,NY19}：
\[\Shv^{\mathrm{all}}_{\mathrm{Nilp}}(\Bun_G)\simeq\IndCoh_{\mathrm{Nilp}}(\LocSys^{\mathrm{Betti}}_{\checkG}).\]
特别地，紧谱范畴 $\operatorname{Perf}(\LocSys^{\mathrm{Betti}}_{\checkG})$ 作用于 $\Shv^{\mathrm{all}}_{\mathrm{Nilp}}(\Bun_G)$；在可呈示形式中，这常写为如下作用：
\[\QCoh(\LocSys^{\mathrm{Betti}}_{\checkG})\curvearrowright\Shv^{\mathrm{all}}_{\mathrm{Nilp}}(\Bun_G).\]
本章的要点是，在 Betti 情形，这一作用可以像\cite{NY19,AGK20a}那样，从局部 Hecke 作用清楚地构造出来。同一机制还产生受限变体 $\LocSys^{\mathrm{restr}}_{\checkG}$，它是可构造层理论和 $\ell$-进层理论中正确的谱对象。

\section{同伦型上的局部系统}\label{sec:12.1}\index{同伦型上的局部系统|hyperpage}\index{Kan 延拓|hyperpage}
令 $Y\in\operatorname{Spc}$ 为一个同伦型；例如，$Y$ 可以是某个拓扑空间的同伦型。定义
\begin{definition*}{}
\[\LocSys(Y):=\{Y\text{ 上的局部系统}\}\simeq\operatorname{Fun}(Y,\operatorname{Vect}).\]
\end{definition*}
等价地，将 $Y$ 看作 $\infty$-群胚，并采用取值为 $\operatorname{Vect}$ 的常值图表，则
\[\LocSys(Y)\simeq\lim_Y\operatorname{Vect}\simeq\underset{Y}{\operatorname{colim}}\operatorname{Vect}.\]
最后一个识别使用 $\operatorname{Vect}$ 作为 DG 范畴的自对偶性。
\begin{note}
应当注意，$\LocSys(Y)$ 依赖于 $Y$ 的同伦型，而不只依赖于 $\pi_1(Y)$，即使 $Y$ 连通也是如此。
\end{note}

对映射 $f:Y_1\to Y_2$，拉回给出
\[f^*:\LocSys(Y_2)\longrightarrow\LocSys(Y_1).\]
在可呈示情形，它具有伴随；当这些伴随由通常的 Kan 延拓公式定义时，分别记作 $f_*$ 与 $f_!$。函子 $f_!$ 即使在拓扑上也已经相当微妙：例如，对 $\mathrm{pt}\to S^1$，所得结果对应一个具有无限 Jordan 块的局部系统。因此，不能期待所有推前都留在较小的可构造范畴或有限秩范畴内。

\section{Betti 局部到整体定理}\label{sec:12.2}\index{局部到整体定理!Betti 情形|hyperpage}
对于同伦型 $Y$，考虑 Betti 谱叠
\[\operatorname{Maps}(Y,B\checkG).\]
它由以下点函子描述刻画。对测试仿射概形 $S$，
\[\operatorname{Maps}(S,\operatorname{Maps}(Y,B\checkG))\simeq\operatorname{Maps}(S\times Y,B\checkG).\]
等价地，一个 $S$-点是一个对称幺半函子
\[\Rep(\checkG)\longrightarrow\QCoh(S)\otimes\LocSys(Y),\]
当所选的层理论带有相应结构时，还须满足通常的右正合或右 $t$-正合条件。对于 $Y=X(\mathbb C)$，这就是 $\LocSys^{\mathrm{Betti}}_{\checkG}(X)$。
\begin{theorem*}{Betti 局部到整体定理\cite{NY19,AGK20a}}
谱幺半范畴 $\operatorname{Perf}(\operatorname{Maps}(Y,B\checkG))$ 在可对偶 DG 范畴 $\mathcal C$ 上的作用，或者在允许相应完备化时，它的可呈示 $\QCoh$-完备化的作用，由下列因子化数据编码。对每个有限集合 $I$，有幺半函子
\[\Rep(\checkG)^{\otimes I}\longrightarrow\End(\mathcal C)\otimes\LocSys(Y^I),\]
\srcpage{61}
并且与有限集合之间的所有映射相容。若 $I\to J$ 是这样的映射，则相应的对角映射 $\Delta:Y^J\to Y^I$ 给出相容性方块
\origdiagram{p061-square.pdf}{0.4781}
\end{theorem*}
要点在于，单个谱范畴 $\operatorname{Perf}(\operatorname{Maps}(Y,B\checkG))$，或它的 $\QCoh$-完备化，将所有有限集合上的 Hecke 相容性统一编码起来。

\begin{transnote}{含糊点拨}
这族函子完全决定作用，靠两点：$\mathcal{C}$ 可对偶，作用被 $\End(\mathcal{C})$ 中的对象记录；$\QCoh$ 的幺半结构让 $\LocSys(Y^I)$ 的张量积承接因子化。相容性由对角映射 $Y^J\to Y^I$ 给出，所以「一个作用」与「一族与有限集映射相容的幺半函子」是同一份数据——后者的好处是每个 $I$ 上都可以逐点计算。[依据：原书 p.61 本段；并见 [NY19]、[AGK20a]。]
\end{transnote}

\section{与 de Rham Ran 构造的比较}\label{sec:12.3}\index{Ran 空间!Betti 与 de Rham 比较|hyperpage}
在 de Rham 情形，将 $Y=X(\mathbb C)$ 与 $\LocSys(Y^I)$ 替换为 de Rham 层理论 $\DMod(X^I)$。局部 Hecke 范畴组织成 Ran 范畴表达式
\[\Rep(\checkG)_{\mathrm{Ran}}\simeq\underset{I}{\operatorname{colim}}\bigl(\Rep(\checkG)^{\otimes I}\otimes\DMod(X^I)\bigr),\]

\begin{transnote}{含糊点拨}
Ran 空间是有限子集配置空间的余极限：$\Rep(\check G)_{\mathrm{Ran}}$ 就是把各 $X^I$ 上的局部 Hecke 范畴沿对角映射粘起来所得的范畴，过渡映射来自 $X^I\to X^J$。取余极限而不是直积，是因为 Hecke 数据只在有限个点上局部存在，粘起来才能把因子化结构编码成单个范畴；核很大、难以描述，正是消失猜想要处理的对象。[来源：nLab "Ran space"，https://ncatlab.org/nlab/show/Ran+space（Last revised 2020-07-27；2026-11-17 访问）；并见 nLab "factorization algebra"（Last revised 2026-08-16）。]
\end{transnote}
其中过渡映射由对角映射构造。有一个自然的幺半函子
\[\Rep(\checkG)_{\mathrm{Ran}}\longrightarrow\QCoh(\LocSys^{\mathrm{dR}}_{\checkG}).\]
它的核很大，而且难以描述；这就是前面各章消失猜想的来源。

Betti 情形较为清楚，但也不是完全形式化的。对于对角映射 $\Delta:Y^J\to Y^I$，构造 Ran 型余极限所需的左 Kan 延拓
\[\Delta_!:\LocSys(Y^J)\longrightarrow\LocSys(Y^I)\]
一般并非幺半函子；甚至像 $\mathrm{pt}\to S^1$ 这样的映射，也会产生具有无限 Jordan 块的局部系统。因此，不能朴素地对取值于 $\LocSys(Y^I)$ 的 Hecke 数据取余极限，就形成单个幺半范畴。\thref{9.3.1}（局部到整体定理）说明，正确的幺半对象仍然是 $\operatorname{Perf}(\operatorname{Maps}(Y,B\checkG))$，或它的 $\QCoh$-完备化。

\section{Hecke 光滑层与幂零奇异支撑}\label{sec:12.4}\index{Hecke 光滑性|hyperpage}
令 $X$ 为曲线，并取 $\cF\in\Shv^{\mathrm{all}}(\Bun_G)$。整体 Hecke 对应为每个 $V\in\Rep(\checkG)$ 给出函子
\[H(V,\text{--}):\Shv^{\mathrm{all}}(\Bun_G)\longrightarrow\Shv^{\mathrm{all}}(\Bun_G\times X).\]
\begin{definition*}{}
若对每个有限集合 $I$ 及每个 $V_I\in\Rep(\checkG)^{\otimes I}$，迭代 Hecke 变换都位于沿曲线方向光滑的子范畴中，即
\[\begin{split}
H_I(V_I,\cF)&\in\Shv^{\mathrm{all}}(\Bun_G)\otimes\LocSys(X(\mathbb C)^I)\\
&\subset\Shv^{\mathrm{all}}(\Bun_G\times X^I),
\end{split}\]
则称 $\cF$ 是 Hecke 光滑的。等价地，所有 Hecke 变换都沿 $X^I$ 方向局部常值。
\end{definition*}

\secref{12.2}中讨论的 Betti 局部到整体定理说明，幂零支撑条件蕴涵 Hecke 光滑性质\cite{NY19}，而 AGKRRV 使用的形式将这一结论加强为等价\cite{AGK20a}
\[\Shv^{\mathrm{all}}_{\mathrm{Nilp}}(\Bun_G)\simeq\Shv^{\mathrm{all}}(\Bun_G)_{\mathrm{Hecke}\text{-}\mathrm{lisse}}.\]
因此，$\Shv^{\mathrm{all}}_{\mathrm{Nilp}}(\Bun_G)$ 上的 Hecke 作用满足 Betti 局部到整体定理的假设，从而得到谱作用
\[\operatorname{Perf}(\LocSys^{\mathrm{Betti}}_{\checkG})\curvearrowright\Shv^{\mathrm{all}}_{\mathrm{Nilp}}(\Bun_G),\]
\srcpage{62}
或者，在进行可呈示完备化之后，得到相应的 $\QCoh$ 作用。

\section{其他层理论与受限变体}\label{sec:12.5}\index{受限变体|hyperpage}\index{拟光滑层|hyperpage}
同一个定理在多种层论背景下都有相应版本：
\begin{itemize}
\item $\mathbb C$ 上的全部 Betti 层，使用 $\Shv^{\mathrm{all}}$；
\item 可构造 Betti 层，通常先转到 $\operatorname{Ind}(\Shv_{\mathrm{cons}})$；
\item étale 层或 $\ell$-进层，系数可取 $\mathbb Z_\ell$、$\mathbb Q_\ell$ 或 $\Qell$；
\item de Rham 层，包括全部 $D$ 模，以及完整或正则完整 $D$ 模的 ind-完备化。
\end{itemize}
在可构造情形，范畴 $\LocSys(X(\mathbb C))$ 过大。正确的替代是
\[\operatorname{QLisse}(X):=\Shv_0(X),\]
其中下标表示零奇异支撑。在 AGKRRV 的构造\cite{AGK20a}中，更准确地说，这是 $\operatorname{Ind}(\operatorname{Lisse}(X))$ 的左完备化，其中 $\operatorname{Lisse}(X)$ 表示所选层理论中的有限秩局部系统。因此，这种限制并非“只允许有限秩”；它保留拟光滑部分，而移除具有非零微局部支撑的层。记号 $\mathrm{r.\,t.\,ex.}$ 意为右 $t$-正合；它是普通 Betti 叠中张量函子右正合条件的可构造对应版本。
\begin{definition*}{}

按以下点函子规则定义受限变体：
\[\LocSys^{\mathrm{restr}}_{\checkG}(S):=\operatorname{Fun}^{\otimes,\mathrm{r.t.ex.}}\bigl(\Rep(\checkG),\operatorname{QLisse}(S\times X)\bigr).\]

\begin{transnote}{易错预警}
右 $t$-正合这个条件不能省：不加的话，$\operatorname{Fun}^{\otimes}$ 里平移和任意扭转都算函子，点函子会多到失去几何意义；加上后函子由 $\Rep(\check G)$ 的紧对象在 $t$-结构上的行为决定，$\LocSys^{\mathrm{restr}}$ 的点才只能沿零奇异支撑方向变动——拟光滑（quasi-lisse）层 正是它的局部模型。[来源：nLab "t-structure"，https://ncatlab.org/nlab/show/t-structure（Last revised 2023-04-20；2026-11-17 访问）；并见原书 p.62 本式与 §12.5。]
\end{transnote}
\end{definition*}
对于仿射测试概形，目标常被启发性地写为 $\QCoh(S)\otimes\operatorname{QLisse}(X)$；要点是，这个族只允许沿 $X$ 在零奇异支撑方向，也就是光滑方向变动。于是，在任何定义了 $\Shv_{\mathrm{Nilp}}$ 的可构造或 $\ell$-进情形中，局部到整体定理给出作用
\[\QCoh(\LocSys^{\mathrm{restr}}_{\checkG})\curvearrowright\Shv_{\mathrm{Nilp}}(\Bun_G).\]

在 $\mathbb C$ 上，受限变体最好被看作 Betti 叠的形式细化，通过 Riemann–Hilbert 对应，可以在解析意义下与 de Rham 图景比较。有自然映射
\[\LocSys^{\mathrm{restr}}_{\checkG}\longrightarrow\LocSys^{\mathrm{Betti}}_{\checkG}.\]
粗略地说，若 $\pi:\LocSys^{\mathrm{Betti}}_{\checkG}\to\LocSys^{\mathrm{coarse}}_{\checkG}$ 表示到半单局部系统的 GIT 粗模空间的映射，则在一个半单化 $\bar\sigma$ 的局部，受限变体以形式完备化 $(\pi^{-1}(\bar\sigma))^\wedge$ 为模型。
\begin{note}
这是一个形式局部描述，并非所有这类完备化的普通不交并。在这个描述中，$\pi(\sigma_1)=\pi(\sigma_2)$ 当且仅当 $\sigma_1$ 与 $\sigma_2$ 具有相同的半单化。de Rham 对应物只有经过解析 Riemann–Hilbert 比较才能得到，因此，不应将它理解为字面上的代数映射 $\LocSys^{\mathrm{restr}}\to\LocSys^{\mathrm{dR}}$。
\end{note}

\section{受限范畴 Langlands}\label{sec:12.6}
由此得到的 Betti 可构造猜想的受限形式是
\begin{conjecture*}{}
\[\Shv_{\mathrm{Nilp}}(\Bun_G)\overset{?}{\simeq}\IndCoh_{\mathrm{Nilp}}(\LocSys^{\mathrm{restr}}_{\checkG}).\]
\end{conjecture*}
\begin{note}
这里有一个需要注意的限制：在某些可构造表述中，这个朴素等价可能过大，或者需要限制到一个直和因子，例如限制到由系数理论所选定的
\srcpage{63}
连通分支。
\end{note}
可靠且用于谱分解的是作用
\[\QCoh(\LocSys^{\mathrm{restr}}_{\checkG})\curvearrowright\Shv_{\mathrm{Nilp}}(\Bun_G).\]
这是按 Langlands 参数分解自守对象的层论替代。

\newpage
\section{范畴迹与有限域上的函数}\label{sec:12.7}\index{范畴迹|hyperpage}\index{谱分解|hyperpage}\index{游走算子|hyperpage}\index{shtuka 模空间|hyperpage}
现在取 $k=\mathbb F_q$，并使用 $\ell$-进层。Frobenius 作用于自守范畴：
\[\Fr:\Shv_{\mathrm{Nilp}}(\Bun_G)\longrightarrow\Shv_{\mathrm{Nilp}}(\Bun_G).\]
预期的层与函数对应陈述是范畴迹恒等式
\[\Tr(\Fr;\Shv_{\mathrm{Nilp}}(\Bun_G))\simeq C^\infty_{\mathrm{stk}}(\Bun_G(\mathbb F_q)),\]
它在通常的完备化选择、有限性条件以及自守函数上来自叠自同构的权重的意义下成立。

回顾范畴迹的定义。对于有限维向量空间 $V$ 和自同态 $f:V\to V$，通常的迹是以下复合得到的标量：

\begin{transnote}{背景直觉}
可对偶范畴指每个对象有对偶、$\mathrm{coev}$ 与 $\mathrm{ev}$ 满足 zig-zag 恒等式；迹的这个复合两端都落在 $\mathrm{Vect}$ 上，所以结果是一个复形而不是 $\mathcal{C}$ 中的对象。取 $F=\Fr$ 时它应当就是 $\Bun_G(\mathbb F_q)$ 上的函数空间；取 $F=\mathrm{Id}$ 时 $\Tr(\mathrm{Id};\mathcal{C})$ 给出 Hochschild 同调，这就是「范畴迹」与通常迹的对照。[依据：原书 p.63 本段；并见 §12.6 的范畴迹形式体系。]
\end{transnote}
\[\bigl[k\xrightarrow{\operatorname{coev}}V\otimes V^*
\xrightarrow{f\otimes\operatorname{Id}}V\otimes V^*
\xrightarrow{\operatorname{ev}}k\bigr].\]
\begin{definition*}{}
对于可对偶 DG 范畴 $\mathcal C$ 和自函子 $F:\mathcal C\to\mathcal C$，相同的公式给出 $\operatorname{Vect}$ 中的对象：
\[\Tr(F;\mathcal C):=\bigl[\operatorname{Vect}\xrightarrow{\operatorname{coev}}\mathcal C\otimes\mathcal C^\vee
\xrightarrow{F\otimes\operatorname{Id}}\mathcal C\otimes\mathcal C^\vee
\xrightarrow{\operatorname{ev}}\operatorname{Vect}\bigr].\]
\end{definition*}
范畴 $\Shv_{\mathrm{Nilp}}(\Bun_G)$ 预期是可对偶的，而且在受限理论中，这一点已在所需的意义下成立\cite{AGK20b}，所以这个迹有意义。

更一般地，对 $V_I\in\Rep(\checkG)^{\otimes I}$，构造以 $X^I$ 为参数的自函子
\[\begin{split}
\Shv_{\mathrm{Nilp}}(\Bun_G)&\xrightarrow{\Fr}\Shv_{\mathrm{Nilp}}(\Bun_G)\\
&\xrightarrow{H_I(V_I,\text{--})}\Shv_{\mathrm{Nilp}}(\Bun_G)\otimes\operatorname{QLisse}(X^I).
\end{split}\]
它的范畴迹是对象
\[\Tr(H_I(V_I,\text{--})\circ\Fr;\Shv_{\mathrm{Nilp}}(\Bun_G))\in\operatorname{QLisse}(X^I).\]
猜想中的识别是：这个对象就是相应 shtuka 模空间的上同调。因此，范畴迹形式体系把从几何 Langlands 到函数论 Langlands 的过渡、以及 V. Lafforgue 的游走算子都几何化了\cite{Gai16a,Laf18}。

最后，$\QCoh(\LocSys^{\mathrm{restr}}_{\checkG})$ 在 $\Shv_{\mathrm{Nilp}}(\Bun_G)$ 上的谱作用，将迹 $\Tr(\Fr;\Shv_{\mathrm{Nilp}}(\Bun_G))$ 提升为 Frobenius 扭曲不动点叠上的拟凝聚层，其中
\[\LocSys^{\mathrm{restr},\Fr}_{\checkG}:=
\LocSys^{\mathrm{restr}}_{\checkG}\underset{\substack{\LocSys^{\mathrm{restr}}_{\checkG}\times\LocSys^{\mathrm{restr}}_{\checkG}\\\Delta,\,(\operatorname{Id},\Fr)}}{\times}\LocSys^{\mathrm{restr}}_{\checkG}.\]
它的点是受限 Weil $\checkG$-局部系统。取全局截面，就恢复上述自守迹空间。这就是函数论 Langlands 参数化背后的谱分解陈述。

\chapter{Whittaker 系数的保守性}\label{ch:13}
\srcpage{64}
本章旨在解释 Whittaker 系数为何能够检测自守范畴的温和部分\cite{FR25,Gai18}。在范畴 Langlands 的 de Rham 形式中，预期有
\[\DMod(\Bun_G)\simeq\IndCoh_{\mathrm{Nilp}}(\LocSys_{\checkG}),\]

\begin{transnote}{含糊点拨}
$\IndCoh_{\mathrm{Nilp}}$ 是 $\LocSys_{\check G}$ 上奇异支撑落在全局幂零锥内的 ind-凝聚层：幂零锥由 $\check{\mathfrak g}^*$ 的幂零锥沿 $\LocSys$ 拉回。de Rham 侧用 $\IndCoh$ 而不是 $\QCoh$，是因为 Hecke 特征层的谱侧天然带奇异支撑；只有温和部分支撑为零，才退化为 $\QCoh$。所以「温和 $\leftrightarrow$ $\QCoh$」是一条独立陈述，不是定义。[来源：nLab "ind-coherent sheaves"（该条目无 Last revised，2026-11-17 访问）；并见原书 p.64 本段与 [Gai18]、[FR25]。]
\end{transnote}
以及
\[\DMod(\Bun_G)^{\mathrm{temp}}\simeq\QCoh(\LocSys_{\checkG}).\]
在 Betti 理论或可构造理论中，相应的谱叠是 $\LocSys^{\mathrm{restr}}_{\checkG}$，而自守侧使用 $\Shv_{\mathrm{Nilp}}(\Bun_G)$，即层范畴中具有幂零奇异支撑的部分。本章记录的定理是重整化 Whittaker 系数函子的保守性：
\begin{theorem*}{}
\[\begin{aligned}
\DMod(\Bun_G)^{\mathrm{temp}}&\xrightarrow{\operatorname{Coeff}^{\mathrm{ren}}}\QCoh(\LocSys_{\checkG}),\\
\Shv_{\mathrm{Nilp}}(\Bun_G)^{\mathrm{temp}}&\xrightarrow{\operatorname{Coeff}^{\mathrm{ren}}_{\mathrm{restr}}}\QCoh(\LocSys^{\mathrm{restr}}_{\checkG}).
\end{aligned}\]
这里，“保守”意味着：若系数为零，则对象本身为零。
\end{theorem*}
在 de Rham 实现中，对应的可构造范畴是特征簇包含于全局幂零锥的正则完整 $D$ 模的 ind-完备化。它并非整个 $\DMod(\Bun_G)$，而是正则完整且具有幂零支撑的部分。

\section{普通系数函子与重整化系数函子}\label{sec:13.1}\index{Whittaker 系数!重整化|hyperpage}\index{余作用|hyperpage}
固定 Borel 子群 $B=TN$，以及非退化加性特征标 $\psi:N\to\Ga$。令 $\Bun_N^\Omega$ 为通常的 $\Omega$-扭转 $N$-丛叠，它带有投影 $p:\Bun_N^\Omega\to\Bun_G$ 和 Whittaker 特征标 $\psi:\Bun_N^\Omega\to\mathbb A^1$。忽略标准的上同调归一化，第一个 Whittaker 系数是紧支撑 de Rham 上同调
\begin{definition*}{}
\[\operatorname{coeff}(\cF):=C_{\mathrm{dR},c}(\Bun_N^\Omega,p^!(\cF)\otimes\psi^!(\exp))\in\operatorname{Vect}.\]
\end{definition*}
等价地，它是将 $\cF$ 拉回到 Whittaker 叠后，以指数层扭转，再 $!$-推前到一点。移位的选择与 Verdier 对偶和下文的微局部指标公式相容。

谱作用给出
\[\QCoh(\LocSys_{\checkG})\otimes\DMod(\Bun_G)\longrightarrow\DMod(\Bun_G).\]
通过对偶，一个标量值泛函 $c:\DMod(\Bun_G)\to\operatorname{Vect}$ 产生一个谱值泛函
\[\begin{split}
\DMod(\Bun_G)&\xrightarrow{\operatorname{coact}}\QCoh(\LocSys_{\checkG})\otimes\DMod(\Bun_G)\\
&\xrightarrow{\operatorname{Id}\otimes c}\QCoh(\LocSys_{\checkG}).
\end{split}\]
取 $c=\operatorname{coeff}$，便得到 $\operatorname{Coeff}^{\mathrm{ren}}$。它的余单位，也就是取全局截面后得到的标量形式，恢复原来的标量系数：
\[\Gamma(\LocSys_{\checkG},\operatorname{Coeff}^{\mathrm{ren}}(\cF))\simeq\operatorname{coeff}(\cF).\]
\begin{note}
单凭这一恢复性质，不能刻画取值于 $\QCoh(\LocSys_{\checkG})$ 的对象，因为 $\LocSys_{\checkG}$ 一般不是仿射的；其定义是上面的余作用公式。同样的说明也适用于受限 Betti／可构造情形。在谱范畴满足通常的刚性假设时，给出到 $\QCoh(\LocSys_{\checkG})$ 的 $\QCoh(\LocSys_{\checkG})$-线性函子，等价于给出取值于 $\operatorname{Vect}$ 的底层标量函子。这就是标量 Whittaker 泛函能够定义谱值泛函的形式原因。
\end{note}

证明使用两个基本事实。第一，
\[\operatorname{coeff}:\DMod(\Bun_G)\longrightarrow\operatorname{Vect}\]
\srcpage{65}
经由到温和商的投影分解。第二，$\QCoh(\LocSys_{\checkG})$ 在 $\DMod(\Bun_G)^{\mathrm{temp}}$ 上的作用在所有 Satake／Hecke 算子下稳定，所以 $\operatorname{Coeff}^{\mathrm{ren}}$ 也经由温和部分分解。

\section{局部 Whittaker 范畴与温和性}\label{sec:13.2}\index{Whittaker 范畴!局部|hyperpage}\index{双 Whittaker 范畴|hyperpage}\index{温和性|hyperpage}
令 $\mathcal K=\mathbb C((t))$、$\cO=\mathbb C[[t]]$。若 DG 范畴 $\mathcal C$ 带有 $G(\mathcal K)$ 作用，则它的非退化 Whittaker 范畴为
\[\operatorname{Wh}(\mathcal C):=\mathcal C^{N(\mathcal K),\psi}.\]
在\cite{Gai18,FR23}的 Whittaker 形式体系中，局部几何 Langlands 的基本思想预言
\[\DMod(G(\mathcal K))\text{-}\mathrm{Mod}^{\mathrm{temp}}\simeq\operatorname{ShvCat}(\LocSys_{\checkG}(D^\circ)),\]
其中 $D^\circ=\Spec\mathcal K$。在这一图景下，双 Whittaker 范畴
\[\operatorname{BiWhit}(G):=\DMod((N(\mathcal K),\psi)\backslash G(\mathcal K)/(N(\mathcal K),\psi))\]
应当识别为
\[\operatorname{BiWhit}(G)\simeq\QCoh(\LocSys_{\checkG}(D^\circ)).\]
对一个 $G(\mathcal K)$-范畴 $\mathcal C$，预期的重构映射是
\[\operatorname{Wh}(\DMod(G(\mathcal K)))\otimes_{\operatorname{BiWhit}(G)}\operatorname{Wh}(\mathcal C)\longrightarrow\mathcal C.\]

\begin{transnote}{含糊点拨}
$\operatorname{BiWhit}(G)$ 是双 Whittaker 范畴，充当这里的「基环」：$\operatorname{Wh}$ 把 $G(\mathcal K)$-范畴送到它的模，$\otimes_{\operatorname{BiWhit}(G)}$ 把两边的 $\operatorname{Wh}$ 值张量起来。重构映射全忠实当且仅当 $\mathcal{C}$ 温和，这是 [Gai18] 的定理；它与 13.1 的谱定义（温和 $\leftrightarrow$ $\QCoh$）是同一件事的两种说法。[依据：原书 p.65 本段；并见 [Gai18]、[FR23]。]
\end{transnote}
重构映射全忠实，其像是由 $\operatorname{Wh}(\mathcal C)$ 生成的 $\DMod(G(\mathcal K))$-子模。因此，$\mathcal C$ 温和，当且仅当这个像是整个 $\mathcal C$。

非分歧情形更为明确。令
\[\operatorname{Sph}_{G,x}:=\DMod(G(\cO)\backslash G(\mathcal K)/G(\cO))\]
为 $x\in X$ 处的球面 Hecke 范畴。带有强 $G(\cO)$ 作用的 $G(\mathcal K)$-范畴 $\mathcal C$，受其球面不变量 $\mathcal C^{G(\cO)}$ 控制，后者是一个 $\operatorname{Sph}_{G,x}$-模。导出 Satake 给出伴随／重构图景
\[\DMod(G(\mathcal K)/G(\cO))\otimes_{\operatorname{Sph}_{G,x}}\mathcal C^{G(\cO)}\rightleftarrows\mathcal C,\]
而几何 Casselman–Shalika 等价，将从球面到 Whittaker 的比较识别为通常的从 $\operatorname{Sph}_{G,x}$ 到 $\Rep(\checkG)$ 的过渡。因此，在非分歧情形中，$\mathcal C$ 温和，当且仅当从球面不变量到 Whittaker 不变量的函子在相关的生成子范畴上保守。

对于 $\Bun_G$，每个 $x\in X$ 都通过 Hecke 修改给出 $\operatorname{Sph}_{G,x}$ 的作用。Arinkin–Gaitsgory 对 $\DMod(\Bun_G)^{\mathrm{temp}}$ 的定义，是从这一作用转到温和谱 Hecke 范畴得到的。Færgeman–Raskin 的定理证明，这一定义与 $x$ 无关\cite{FR23}。还有左温和与右温和的变体；在局部，预期它们相同，而在全局，这反映为 Whittaker 系数与所选点无关。

\section{函数论中的对应图景}\label{sec:13.3}
这一图景与经典自守理论平行。自守函数定义在
\[G(F)\backslash G(\A)/G(\cO)\]
上，而 Whittaker 系数是沿
\[(N(\A),\psi)\backslash G(\A)/G(\cO)\]
对特征标 $\psi$ 作积分。
\srcpage{66}
在几何上，后一个商被替换为一般位置定义的 $N$-约化的 Whittaker 叠，积分被替换为函子 $\operatorname{coeff}$。本章所述的保守性定理说明，限制到温和范畴后，这个几何 Whittaker 泛函能够检测所有对象的非零性。

\section{归约到幂零奇异支撑}\label{sec:13.4}\index{反温和对象|hyperpage}\index{幂零锥!非一般位置正则轨迹|hyperpage}
证明从范畴 Langlands 的受限／幂零形式所提供的比较方块开始。按此处采用的方向，示意地写为
\origdiagram{p066-comparison.pdf}{0.5829}
\begin{note}
竖直箭头应当理解为所选形式体系中可用的比较函子或投影函子，而不是同一公共范畴中的字面包含。
\end{note}
因此，只需在 $\Shv_{\mathrm{Nilp}}(\Bun_G)$ 内理解保守性的机制；在这一范畴中，可以使用微局部奇异支撑与特征圈。

回顾
\[T^*\Bun_G=\operatorname{Higgs}_G=\{(P_G,\varphi)\mid\varphi\in H^0(X,\mathfrak g_{P_G}^*\otimes\Omega_X)\}.\]
Hitchin 映射是
\[h:T^*\Bun_G\longrightarrow\cA_G:=\Gamma(X,(\mathfrak g^*/\!/G)_{\Omega_X}),\]
全局幂零锥为 $\mathrm{Nilp}:=h^{-1}(0)$。令 $\mathrm{Nilp}_{\mathrm{irreg}}\subset\mathrm{Nilp}$ 为幂零 Higgs 场一般位置正则轨迹的闭补集。因此，它的开补集就是一般位置正则幂零 Higgs 丛的轨迹。
\begin{note}
这里的“非正则”是指不在一般位置正则，并不要求在 $X$ 的每一点都非正则。
\end{note}
\begin{theorem*}{幂零层情形中的微局部定理\cite{FR25,NT25}}
作为满子范畴，有
\[\Shv_{\mathrm{Nilp}_{\mathrm{irreg}}}(\Bun_G)=\Shv_{\mathrm{Nilp}}(\Bun_G)^{\mathrm{anti}\text{-}\mathrm{temp}}.\]

\begin{transnote}{含糊点拨}
反温和指谱 Hecke 作用中负向部分（$t$-结构负次数）作用为零的层。$\mathrm{Nilp}_{\mathrm{irreg}}$ 去掉一般正则幂零轨迹，于是「支撑落在 $\mathrm{Nilp}_{\mathrm{irreg}}$」等价于「不撞该轨迹」，也就是反温和；这条微局部定理把几何条件翻译成表示论条件，后面才能把 $\operatorname{SS}(F)\subset\mathrm{Nilp}_{\mathrm{irreg}}$ 直接读成 $F$ 反温和，从而它与温和子范畴只交于零对象。[依据：原书 p.66 本段；并见 [FR25]、[NT25]。]
\end{transnote}
\end{theorem*}
它说明：若奇异支撑从不与一般位置正则幂零轨迹相交，则该对象在温和商中不可见；反之亦然。

\section{Kostant 分支与第一个系数}\label{sec:13.5}\index{Kostant 分支|hyperpage}\index{微局部指标定理|hyperpage}\index{特征圈|hyperpage}\index{微茎|hyperpage}
全局 Kostant 截面是 Chevalley 映射的 Kostant 截面 $\mathfrak g^*/\!/G\to\mathfrak g^{*,\mathrm{reg}}$ 的全局对应物，并在取 $G$ 的商后理解。进行通常的全局扭转选择后，它在 Hitchin 底空间零点上的值给出一个特殊的正则幂零轨迹，记作 $\mathfrak f^{\mathrm{glob}}\subset\mathrm{Nilp}$。令 $\mathrm{Nilp}_{\mathrm{Kos}}$ 为包含这个轨迹的唯一不可约分支。在微局部意义下，$\mathrm{Nilp}_{\mathrm{Kos}}$ 是指数 Whittaker 核的对应物，如\cite{NT25}中移位微茎的解释所示：将 Lagrange 图像 $d\psi\subset T^*\Bun_N^\Omega$ 与从 $\Bun_N^\Omega$ 到 $\Bun_G$ 的 Lagrange 对应复合，就得到全局 Kostant 切片。

特征圈计算首先对紧的可构造对象表述；可呈示范畴中的保守性陈述则由连续性与紧生成推出。

对可构造对象 $\cF\in\Shv_{\mathrm{Nilp}}(\Bun_G)$，将其特征圈写为
\[\operatorname{CC}(\cF)=\sum_{Z\subset\mathrm{Nilp}}c_Z(\cF)[Z],\qquad c_Z(\cF)\in\mathbb Z,\]
\srcpage{67}
其中求和遍及 $\mathrm{Nilp}$ 的不可约分支。
\begin{theorem*}{Færgeman–Raskin 微局部指标定理\cite{FR25}}
\[\chi(\operatorname{coeff}(\cF))=(-1)^{\dim\Bun_G}c_{\mathrm{Nilp}_{\mathrm{Kos}}}(\cF).\]
\end{theorem*}
因此，若 $\mathrm{Nilp}_{\mathrm{Kos}}$ 作为 $\operatorname{SS}(\cF)$ 的不可约分支出现（等价地，以非零特征圈重数出现），则 $\operatorname{coeff}(\cF)\ne0$。证明使用正则幂零轨迹的光滑性，以及全局 Kostant 截面沿特殊轨迹 $\mathfrak f^{\mathrm{glob}}$ 与 Hitchin 零纤维相交这一事实。

\newpage
\section{其他 Whittaker 系数与 Hecke 平移}\label{sec:13.6}\index{Hecke 移动定理|hyperpage}
令 $D=\sum_i\check\lambda_i x_i$ 为 $X$ 上取值于 $\check\Lambda^+$ 的有效除子。它对应一个对象 $V^D\in\Rep(\checkG)_{\mathrm{Ran}}$。在点互异的轨迹上，这是放置在各点 $x_i$ 的表示 $V^{\check\lambda_i}$ 的张量积；在 Ran 空间上，则是相应的因子化延拓。还有相应的 $D$-移位 Whittaker 叠；为避免固定某一种符号约定，此处将它记为 $\Bun_N^{\Omega,D}$，其系数记为 $\operatorname{coeff}_D$。
\begin{theorem*}{几何 Casselman–Shalika 公式}
\[\operatorname{coeff}_D(\cF)\simeq\operatorname{coeff}(V^D\star\cF).\]
\end{theorem*}
因此，若 $\mathrm{Nilp}_{\mathrm{Kos}}$ 作为 $\operatorname{SS}(V^D\star\cF)$ 的分支出现，则 $\operatorname{coeff}_D(\cF)\ne0$。由于 $\operatorname{Coeff}^{\mathrm{ren}}$ 与 Hecke／谱作用相容，这一非消失蕴涵 $\operatorname{Coeff}^{\mathrm{ren}}(\cF)\ne0$。

剩余的几何输入是\cite{FR25}中的 Hecke 移动定理：
\begin{theorem*}{}
若 $\operatorname{SS}(\cF)\not\subset\mathrm{Nilp}_{\mathrm{irreg}}$，则存在 $D$，使 $\mathrm{Nilp}_{\mathrm{Kos}}$ 作为 $\operatorname{SS}(V^D\star\cF)$ 的分支出现。
\end{theorem*}
该定理通过仿射 Springer 纤维的局部计算，以及一个全局铺展论证证明。非正式地说，Hecke 修改可以将任意一般位置正则的幂零奇异方向移动到 Kostant 分支。
\enlargethispage{\baselineskip}
\begin{proof}

现在设 $\cF\in\Shv_{\mathrm{Nilp}}(\Bun_G)^{\mathrm{temp}}$，并假设相应的重整化系数为零，即在受限理论中假设
\[\operatorname{Coeff}^{\mathrm{ren}}_{\mathrm{restr}}(\cF)=0.\]
Hecke 相容性于是蕴涵每个 $D$-系数都消失。由 Hecke 移动定理，$\operatorname{SS}(\cF)\subset\mathrm{Nilp}_{\mathrm{irreg}}$。因此，$\cF$ 是反温和的。由于温和子范畴与反温和子范畴的交为零，得到 $\cF=0$。这就证明了保守性。通过上面的比较方块，同一论证给出 de Rham 情形的陈述。
\end{proof}

\chapter{Hitchin 系统的量子化}\label{ch:14}
% 原PDF第67页，第14章开始。
本章旨在解释 Hitchin Hamilton 量的 Beilinson–Drinfeld 量子化，以及由此得到的 Hecke 特征层\cite{BD91,BD05}。仍记自守群为 $G$，其 Langlands 对偶群为 $\checkG$，并假设 $\checkG$ 为伴随型。记 $\check{\mathfrak g}=\operatorname{Lie}(\checkG)$，固定 Borel 子群 $\check B\subset\checkG$，置
\[\mathfrak c_{\checkG}:=\check{\mathfrak g}/\!/\checkG=\Spec k[\check{\mathfrak g}]^{\checkG},\qquad
\mathfrak c_{\checkG,\Omega}:=(\mathfrak c_{\checkG})_{\Omega_X}.\]
若 $d_1,\ldots,d_r$ 为 $\check{\mathfrak g}$ 的指数，则 $\mathfrak c_{\checkG,\Omega}\simeq\bigoplus_i\Omega_X^{d_i+1}$，Hitchin 底空间为
\[\operatorname{Hitch}_{\checkG}(X):=\Gamma(X,\mathfrak c_{\checkG,\Omega}).\]
本章的要点是，同一个仿射空间出现两次：在经典情形，它是 Hitchin 底空间；在量子情形，它是临界水平中心的相伴分次。

\section{oper 作为 Hitchin 底空间上的挠子}\label{sec:14.1}\index{oper|hyperpage}\index{Hitchin 底空间|hyperpage}\index{射流|hyperpage}
\srcpage{68}
令 $\pi:Y\to X$ 为 $D_X$-概形。
\begin{definition*}{}
依照 Beilinson–Drinfeld\cite{BD05}，$Y$ 上相对于 $X$ 的 $\checkG$-oper 是一个二元组 $(\mathcal F_{\check B},\nabla)$，其中 $\mathcal F_{\check B}$ 是 $Y$ 上的 $\check B$-挠子，$\nabla$ 是诱导的 $\checkG$-挠子 $\mathcal F_{\checkG}:=\checkG\times^{\check B}\mathcal F_{\check B}$ 上相对于 $\pi$ 的联络，并且 $\check B$-约化的第二基本形式满足 oper 条件：
\[c(\nabla)\in\Gamma(Y,(\check{\mathfrak g}/\check{\mathfrak b})_{\mathcal F_{\check B}}\otimes\pi^*\Omega_X)\]
位于由简单负根方向组成的子丛中的开 $\check B$-轨道；等价地，它的简单负根分量处处非零。
\end{definition*}
函子 $Y\mapsto\Op_{\checkG}(Y)$ 由一个 $D_X$-概形表示，记为 $\Op^D_{\checkG}\to X$。它的全局截面是通常的 oper 空间 $\Op_{\checkG}(X)$。
\begin{proposition}{}{14.1.1}
预叠 $\Op_{\checkG}(X)$ 是向量空间 $\operatorname{Hitch}_{\checkG}(X)$ 的挠子。因此，若

\begin{transnote}{含糊点拨}
底空间 $\operatorname{Hitch}_{\checkG}(X)=\bigoplus_i H^0(X,\Omega_X^{d_i+1})$ 是仿射空间，$d_i$ 是 $\check{\mathfrak g}$ 齐次生成元的次数；oper 空间是它的挠子：任取两个 oper 之差给出一个 Hitchin 场，但没有典范的 oper 可以固定，所以只能说是「挠子」。过滤来自 Rees 构造，因此 $\operatorname{gr} A\simeq A^{\mathrm{cl}}$——这就是「量子化在经典极限回到 Hitchin 系统」的确切含义。[来源：nLab "Hitchin fibration"（该条目无 Last revised，2026-11-17 访问）；并见原书 p.68 本段。]
\end{transnote}
\[A:=\Gamma(\Op_{\checkG}(X),\cO),\qquad A^{\mathrm{cl}}:=\Gamma(\operatorname{Hitch}_{\checkG}(X),\cO),\]
则 $A$ 的自然过滤的相伴分次为 $\operatorname{gr}A\simeq A^{\mathrm{cl}}$。
\end{proposition}

命题\ref{pro:14.1.1}的证明归约到引理\ref{lem:14.1.2}、\ref{lem:14.1.3}和\ref{lem:14.1.4}。
\begin{lemma}{}{14.1.2}
$\Op_{\checkG}(X)$ 非空。
\end{lemma}
\Needspace{8\baselineskip}
\begin{lemma}{}{14.1.3}
令 $W$ 为遗忘函子 $\operatorname{Sch}_{/X_{\mathrm{dR}}}\to\operatorname{Sch}_{/X}$ 的右伴随。对 $X$ 上的向量丛 $E$，可以将 $W(E)$ 理解为射流的 $D_X$-概形，使得
\[W(E)(D_x)=\operatorname{Jets}_x(E).\]
则 $\Op^D_{\checkG}$ 是 $D_X$-概形 $W(\mathfrak c_{\checkG,\Omega})$ 的挠子。
\end{lemma}
\begin{lemma}{}{14.1.4}
在 $\check{\mathfrak g}$ 中选定主 $\mathfrak{sl}_2$ 三元组 $(e,h,f)$。存在依赖于这一选择的 $X$ 上典范 $\check B$-挠子 $\mathcal F^{\mathrm{op}}_{\check B}$，使得对于任意 $Y\to X$ 上的每个 $\checkG$-oper，其底层 $\check B$-挠子都典范地等于 $\mathcal F^{\mathrm{op}}_{\check B}$ 的拉回。相应的 $\Bun_{\checkG}$ 中的点，在连通分支的意义下与辅助选择无关。
\end{lemma}

特别地，遗忘联络给出到 $\Bun_{\checkG}$ 的常值映射：
\origdiagram{p068-square.pdf}{0.3484}

\section{局部正规形}\label{sec:14.2}\index{Drinfeld–Sokolov 规范变换|hyperpage}\index{Kostant 切片|hyperpage}
在 $X$ 上局部讨论，选择坐标 $t$，并先假设典范的 $\check B$-挠子平凡。在这样的平凡化下，oper 联络具有形式
\[\nabla=d+(f+\rho(t))\,dt,\qquad\rho(t)\in\check{\mathfrak b}((t)),\]
正则情形下则有 $\rho(t)\in\check{\mathfrak b}[[t]]$。这种形式并不唯一。oper 的 Drinfeld–Sokolov 规范变换陈述\cite{BD05,Fre02}说，每个 $\check N((t))$-规范等价类都有唯一的代表元
\[\nabla=d+(f+v(t))\,dt,\qquad v(t)\in\check{\mathfrak g}^{e}((t)),\]
正则 oper 满足 $v(t)\in\check{\mathfrak g}^{e}[[t]]$。这里 $\check{\mathfrak g}^{e}$ 是 $e$ 的中心化子。这就是 Kostant 切片：真正的参数是 $\check{\mathfrak g}^{e}$ 中的 Kostant 坐标，而非整个 $\check{\mathfrak b}$。

\begin{transnote}{含糊点拨}
$\check N((t))$ 以 $n(t)\mapsto\nabla+n(t)$ 的方式作用在 oper 上；Kostant 的 $\mathfrak{sl}_2$ 理论给出横截切片：$\check{\mathfrak g}^{e}$ 与 $\check{\mathfrak b}$ 交于唯一代表。$f$ 定 oper 的最低次项，$e$ 定切片方向，因此真正的参数是 $\check{\mathfrak g}^{e}$ 里的 Kostant 坐标，而不是整个 $\check{\mathfrak b}$——这就是「Kostant 切片」在此的含义。[来源：nLab "Drinfeld-Sokolov reduction"，https://ncatlab.org/nlab/show/Drinfeld-Sokolov+reduction（Last revised 2013-06-10；2026-11-17 访问）；并见原书 p.68 本段与 [BD05]、[Fre02]。]
\end{transnote}
\srcpage{69}
\begin{proof}

先取 $\checkG=\operatorname{PGL}_2$。oper 条件平凡化 $(\check{\mathfrak g}/\check{\mathfrak b})_{\mathcal F_{\check B}}\otimes\pi^*\Omega_X$。因此，$\Op_{\operatorname{PGL}_2}(Y)$ 是 $\pi^*\Omega_X^{\otimes2}$ 的一个可能为空的挠子。当 $g>1$ 时，由于 $H^1(X,\Omega_X^{\otimes2})=0$，它非空；当 $g=0,1$ 时，可以直接构造。因此，三个引理对 $\operatorname{PGL}_2$ 成立。在这一情形中，通常的陈述是：射影联络组成二次微分空间的挠子。

对于一般伴随型 $\checkG$，沿主同态 $\operatorname{PGL}_2\to\checkG$ 诱导一个 $\operatorname{PGL}_2$-oper。这给出一个 $\checkG$-oper，从而证明非空性。Kostant 同构给出
\[(\check{\mathfrak g}^{e})_{\mathcal F^{\mathrm{op}}_{\check B}}\otimes\Omega_X\simeq\mathfrak c_{\checkG,\Omega},\]
所以同一论证将一般 oper 的自由参数识别为 Hitchin 向量丛。这证明 $\Op^D_{\checkG}$ 是 $W(\mathfrak c_{\checkG,\Omega})$-挠子，进而 $\Op_{\checkG}(X)$ 是 $\operatorname{Hitch}_{\checkG}(X)$-挠子。
\end{proof}

\section{Feigin–Frenkel 中心}\label{sec:14.3}\index{Feigin–Frenkel 中心|hyperpage}\index{临界水平|hyperpage}\index{真空模|hyperpage}\index{手征代数|hyperpage}
令 $x\in X$，$D_x=\Spec k[[t]]$，$D_x^\times=\Spec k((t))$。记 $\widehat{\mathfrak g}_{\mathrm{crit},x}$ 为临界水平的仿射 Kac–Moody 代数，
\[\mathbb V_{\mathrm{crit},x}:=\operatorname{Ind}_{\mathfrak g[[t]]\oplus k\mathbf1}^{\widehat{\mathfrak g}_{\mathrm{crit},x}}k\]
为真空模。有两个相关的中心：
\[Z_{\mathrm{crit},x}:=Z(\widehat U_{\mathrm{crit},x}),\qquad
\mathfrak z_{\mathrm{crit}}(D_x):=\End_{\widehat{\mathfrak g}_{\mathrm{crit},x}}(\mathbb V_{\mathrm{crit},x}).\]
真空模随 $x$ 变动，组成临界手征代数 $\mathbb V_{\mathrm{crit}}$。它的中心 $\mathfrak z_{\mathrm{crit}}:=Z(\mathbb V_{\mathrm{crit}})$ 是交换手征代数；等价地，它是因子化意义下的交换 $D_X$-代数，所以 $\Spec\mathfrak z_{\mathrm{crit}}$ 是一个 $D_X$-概形。
\begin{theorem}{Feigin–Frenkel，参见\cite{Fre02}}{14.3.1}
存在典范识别
\[\operatorname{Spf}Z_{\mathrm{crit},x}\simeq\Op_{\checkG}(D_x^\times),\qquad
\Spec\mathfrak z_{\mathrm{crit}}(D_x)\simeq\Op_{\checkG}(D_x),\]
在全局，有 $D_X$-概形的识别
\[\Spec\mathfrak z_{\mathrm{crit}}\simeq\Op^D_{\checkG}.\]
在 PBW 过滤下，经典极限是
\[\Spec\operatorname{gr}\mathfrak z_{\mathrm{crit}}\simeq W(\mathfrak c_{\checkG,\Omega}).\]
\end{theorem}
\begin{proof}

证明思路如下。首先，$\mathbb V_{\mathrm{crit}}$ 的手征代数结构使 $\mathfrak z_{\mathrm{crit}}$ 具有因子化性质；在局部，这产生一个从通常临界中心的形式谱到穿孔圆盘中心的映射。\thref{14.3.1}说明这个映射是同构。因此，剩下的任务是将 $\Spec\mathfrak z_{\mathrm{crit}}(D)$ 与 $D$ 上的正则 oper 识别，并使这一识别与 $\Aut(D)$ 相容。

正则映射 $\Spec\mathfrak z_{\mathrm{crit}}(D)\to\Op_{\checkG}(D)$ 通过 Tannaka 方法构造。需要在 $\Spec\mathfrak z_{\mathrm{crit}}(D)\times D$ 上构造一个 $\checkG$-挠子、一个联络，以及满足 oper 条件的 $\check B$-约化，并要求这些数据都是 $\Aut(D)$-等变的。几何 Satake 与临界全局截面给出对称幺半函子
\[F_D:\Rep(\checkG)\longrightarrow\mathfrak z_{\mathrm{crit}}(D)\text{-}\mathrm{Mod},\]
由性质
\[\Gamma(\Gr_G,\operatorname{Sat}(V))_{\mathrm{crit}}\simeq F_D(V)\otimes_{\mathfrak z_{\mathrm{crit}}(D)}\mathbb V_{\mathrm{crit}}\]
刻画。由 Tannaka 对偶，$F_D$ 确定 $\Spec\mathfrak z_{\mathrm{crit}}(D)$ 上的 $\checkG$-挠子。$\check B$-约化从最低旋转权得到。对不可约 $\checkG$-表示 $V^{\check\lambda}$，令 $\operatorname{IC}_{\check\lambda}=\operatorname{Sat}(V^{\check\lambda})$。若 $L_0=-t\partial_t\in\operatorname{Lie}\Aut^0(D)$，则 $\Gamma(\Gr_G,\operatorname{IC}_{\check\lambda})_{\mathrm{crit}}$ 的最低 $L_0$-特征值为 $-\langle\check\lambda,\rho\rangle$；相应特征空间是一维的，并且 $\check B$-不变。这些直线满足 Plücker 关系，从而给出所需的 $\check B$-约化。最后一步是证明得到的到 $\Op_{\checkG}(D)$ 的映射为同构。

\begin{transnote}{含糊点拨}
$L_0=-t\partial_t$ 来自形式圆盘的旋转（伸缩）作用；最低特征值 $-\langle\check\lambda,\rho\rangle$ 里的 $\rho$ 是半和正根，这正是临界扭转与 $\rho$-移位的来源。最低权空间一维且 $\check B$-不变，才能当纤维化的一根直线；这些直线之间满足 Plücker 关系，来自 $\mathrm{IC}_{\check\lambda}=\operatorname{Sat}(V^{\check\lambda})$ 的函子性，而不是额外构造。[依据：原书 p.69 本段；并见 [Fre02]。]
\end{transnote}
\end{proof}
\section{全局微分算子与量子 Hitchin 映射}\label{sec:14.4}\index{量子 Hitchin 映射|hyperpage}\index{临界扭转微分算子|hyperpage}
\srcpage{70}
令 $\cD_{\mathrm{crit}}$ 为 $\Bun_G$ 上临界扭转微分算子层。局部临界中心通过局部化作用于全局微分算子。在局部，使用如下标准构造：若一对 $(\mathfrak h,K)$ 作用于光滑叠 $Y$，而 $y\in Y$ 是 $K$-不动点，则齐性模型上的不变微分算子映入 $Y$ 上的微分算子：
\[(U(\mathfrak h)/U(\mathfrak h)\mathfrak k)^K\longrightarrow\Gamma(Y,\cD_Y).\]
取 $Y=\Bun_G$、$\mathfrak h=\mathfrak g((t))$、$K=G[[t]]$，得到
\[\mathfrak z_{\mathrm{crit}}(D_x)=\End(\mathbb V_{\mathrm{crit},x})\longrightarrow\Gamma(\Bun_G,\cD_{\mathrm{crit}}).\]
这些局部映射与变量 $x$ 中的联络相容，因此下降到水平截面：
\[\mathfrak z_{\mathrm{crit}}(X)\longrightarrow\Gamma(\Bun_G,\cD_{\mathrm{crit}}).\]
取相伴分次给出经典 Hitchin 映射
\[\operatorname{gr}\mathfrak z_{\mathrm{crit}}(X)\longrightarrow\Gamma(T^*\Bun_G,\cO).\]

利用 Feigin–Frenkel 识别与上述挠子陈述，置
\[\begin{aligned}
A&:=\Gamma(\Op_{\checkG}(X),\cO)=\Gamma(X_{\mathrm{dR}},\mathfrak z_{\mathrm{crit}}),\\
A^{\mathrm{cl}}&:=\Gamma(\operatorname{Hitch}_{\checkG}(X),\cO)\simeq\operatorname{gr}A.
\end{aligned}\]
\begin{theorem*}{}
上述映射为同构\cite{BD91}：
\[A\xrightarrow{\sim}\Gamma(\Bun_G,\cD_{\mathrm{crit}}),\qquad
A^{\mathrm{cl}}\xrightarrow{\sim}\Gamma(T^*\Bun_G,\cO).\]
\end{theorem*}
对 $\GL_n$ 方向上的正特征方法，参见\cite{BCTZ16}。因此，量子 Hitchin Hamilton 量的交换代数，就是 $\checkG$-oper 空间上的函数代数。第二个同构是经典陈述：$T^*\Bun_G$ 上的全局函数来自 Hitchin 底空间。

\section{oper 对应的自守模}\label{sec:14.5}\index{Hecke 特征层!由 oper 构造|hyperpage}\index{Beilinson–Drinfeld 局部化|hyperpage}
令 $\sigma\in\Op_{\checkG}(X)$。通过 $\Op_{\checkG}(X)\subset\LocSys_{\checkG}(X)$，它有一个底层 $\checkG$-局部系统，仍记为 $\sigma$。令 $k_\sigma$ 为 $A=\Gamma(\Op_{\checkG}(X),\cO)$ 的相应特征标。定义临界扭转 $D$ 模
\begin{definition*}{}
\[\Aut_\sigma:=\cD_{\mathrm{crit}}\otimes_{A,\sigma}k_\sigma\in\DMod(\Bun_G).\]
\end{definition*}
\begin{theorem}{Beilinson–Drinfeld\cite{BD91}}{14.5.1}
$\Aut_\sigma$ 是特征值为 $\sigma$ 的 Hecke 特征层。等价地，对每个 $x\in X$ 及每个 $V\in\Rep(\checkG)$，
\[H_{V,x}(\Aut_\sigma)\simeq V_{\sigma,x}\otimes\Aut_\sigma,\]
其中 $H_{V,x}$ 是对应于 $V$ 的 Hecke 函子，而 $V_{\sigma,x}$ 是由 $V$ 和 $\sigma$ 关联的局部系统在 $x$ 处的纤维。
\end{theorem}
\begin{proof}

证明使用 Beilinson–Drinfeld 局部化
\[\operatorname{Loc}_x:\widehat{\mathfrak g}_{\mathrm{crit},x}\text{-}\mathrm{Mod}^{G[[t]]}\longrightarrow\DMod_{\mathrm{crit}}(\Bun_G),\]
它满足
\[\operatorname{Loc}_x(\mathbb V_{\mathrm{crit},x})\simeq\cD_{\mathrm{crit}}.\]
\srcpage{71}
对 $V\in\Rep(\checkG)$，Hecke 作用与局部化相容：
\[\begin{aligned}
\operatorname{Sat}_x(V)\star\cD_{\mathrm{crit}}
&\simeq\operatorname{Loc}_x(\operatorname{Sat}_x(V)\star\mathbb V_{\mathrm{crit},x})\\
&\simeq\operatorname{Loc}_x\bigl(F_{D_x}(V)\otimes_{\mathfrak z_{\mathrm{crit}}(D_x)}\mathbb V_{\mathrm{crit},x}\bigr).
\end{aligned}\]
令 $\mathcal F$ 为 $\Op_{\checkG}(X)$ 上的普遍 $\checkG$-挠子，等价地，是普遍 oper 局部系统。函子 $F_{D_x}$ 是这一普遍挠子在形式圆盘上的限制，所以上式变为
\[\operatorname{Sat}_x(V)\star\cD_{\mathrm{crit}}\simeq V_{\mathcal F,x}\otimes_A\cD_{\mathrm{crit}}.\]
在 $A$ 上与 $k_\sigma$ 作张量积，得到
\[\operatorname{Sat}_x(V)\star\Aut_\sigma\simeq V_{\sigma,x}\otimes\Aut_\sigma.\]
这正是 Hecke 特征性质。
\end{proof}

\section{真空生成与下一步的局部图景}\label{sec:14.6}\index{基本局部等价|hyperpage}\index{非分歧 oper|hyperpage}
\secref{14.5}的构造只使用临界 Kac–Moody 范畴中由真空模生成的部分。示意地，有
\origdiagram{p071-square.pdf}{0.4918}
其中，$\mathbb V$-生成意为由真空对象生成。下一步是用整个临界范畴替换这个真空生成的子范畴。预期的局部形式是基本局部等价。置
\[\Op^{\mathrm{unr}}_{\checkG,x}:=\Op_{\checkG}(D_x^\times)\times_{\LocSys_{\checkG}(D_x^\times)}\LocSys_{\checkG}(D_x).\]
于是，忽略通常的 $*$ 型与 $!$ 型 $\IndCoh$ 约定的区别，
\[\operatorname{KL}_{G,\mathrm{crit},x}\simeq\IndCoh(\Op^{\mathrm{unr}}_{\checkG,x}).\]
局部化应当提升为如下形式的相对构造：
\[\begin{split}
\operatorname{KL}_{G,\mathrm{crit},x}\otimes_{\QCoh(\Op^{\mathrm{unr}}_{\checkG,x})}
\QCoh(\Op_{\checkG}(X\setminus\{x\})^{\mathrm{unr}})\\
\longrightarrow\DMod_{\mathrm{crit}}(\Bun_G),
\end{split}\]
其中定义
\[\Op_{\checkG}(X\setminus\{x\})^{\mathrm{unr}}:=\Op_{\checkG}(X\setminus\{x\})\times_{\LocSys_{\checkG}(X\setminus\{x\})}\LocSys_{\checkG}(X).\]
这就是从 oper 的 Beilinson–Drinfeld 特征层，通向\cite{FG05,FG07}中范围更广的临界水平局部理论的桥梁。


\chapter{基本局部等价}\label{ch:15}
% 原PDF第71页，第15章开始。
本章解释整体局部化论证背后的局部输入\cite{FG05,FG07,ABC24a}。其基本思想是：临界水平的球面 Kac–Moody 表示，等同于非分歧，也就是单值性平凡的 oper 空间上的层。这将\chref{14}的真空生成图景提升到整个 Kazhdan–Lusztig 范畴。

\section{从 Beilinson–Drinfeld 到局部等价}\label{sec:15.1}\index{基本局部等价|hyperpage}
\srcpage{72}
回顾证明 $\Aut_\sigma$ 的 Hecke 特征性质时使用的两个输入。第一，对点 $x\in X$，Beilinson–Drinfeld 局部化函子给出
\[\operatorname{Loc}_x:\widehat{\mathfrak g}_{\mathrm{crit},x}\text{-}\mathrm{Mod}^{L^+G_x}\longrightarrow\DMod_{\mathrm{crit}}(\Bun_G),\]
而中心通过将局部 oper 上的函数拉回到全局 oper 来作用：
\origdiagram{p072-square.pdf}{0.4850}
这里，上标 $\mathbb V$-生成表示由真空模生成的子范畴。第二，对 $V\in\Rep(\checkG)$，几何 Satake 与临界全局截面给出
\[\operatorname{Sat}_x(V)\star\mathbb V_{\mathrm{crit},x}\simeq F_x(V)\otimes_{\mathfrak z_{\mathrm{crit}}(D_x)}\mathbb V_{\mathrm{crit},x},\]
其中 $F_x(V)$ 是 $\Op_{\checkG}(D_x)\simeq\Spec\mathfrak z_{\mathrm{crit}}(D_x)$ 上与普遍 $\checkG$-局部系统关联的向量丛。基本局部等价用整个球面临界范畴的等价，替换这一真空生成的陈述。

\section{非分歧 oper}\label{sec:15.2}\index{非分歧 oper|hyperpage}\index{Kazhdan–Lusztig 范畴|hyperpage}\index{临界 Weyl 对象|hyperpage}
令 $D=\Spec k[[t]]$、$D^\times=\Spec k((t))$、$L^+G=G[[t]]$、$LG=G((t))$。记
\[\operatorname{KL}_{G,\mathrm{crit}}:=\widehat{\mathfrak g}_{\mathrm{crit}}\text{-}\mathrm{Mod}^{L^+G}\]
为临界 Kazhdan–Lusztig 范畴。FLE 中的谱空间是比 $D$ 上正则 oper 空间更大的纤维积

\begin{transnote}{易错预警}
非分歧只要求底层局部系统延拓到 $D$，并不要求 oper 约化在 $D$ 上也正则延拓；所以谱空间取纤维积 $\Op_{\checkG}(D^\times)\times_{\LocSys_{\checkG}(D^\times)}\LocSys_{\checkG}(D)$，即「在 $D^\times$ 上是 oper、在 $D$ 上只是局部系统」。它比 $\Op_{\checkG}(D)$ 大：后者的点还要在 $t=0$ 处满足一般位置条件。缺这一对比，就看不懂 15.2 为何要用 $\check\lambda$ 分层去度量「缺陷」。[依据：原书 p.72 本段；并见 [FG07]、[ABC24a]。]
\end{transnote}
\begin{definition*}{}
\[\Op^{\mathrm{unr}}_{\checkG}:=\Op_{\checkG}(D^\times)\times_{\LocSys_{\checkG}(D^\times)}\LocSys_{\checkG}(D).\]
\end{definition*}
因此，它的一个点是穿孔圆盘上的 oper，其底层 $\checkG$-局部系统可延拓到 $D$。映射 $\Op_{\checkG}(D)\to\Op^{\mathrm{unr}}_{\checkG}$ 不是同构：局部系统可以延拓，并不迫使 oper 约化也正则地延拓。

具体地，选择平凡化后，$D^\times$ 上的 oper 具有形式 $\nabla=d+\omega$，其中 $\omega\in\check{\mathfrak g}((t))\,dt$。非分歧意味着，经过某个 $\checkG((t))$-规范变换后，它成为 $d+\eta$，其中 $\eta\in\check{\mathfrak g}[[t]]\,dt$。但同一个规范变换未必使 $\eta\bmod\check{\mathfrak b}[[t]]$ 在 $t=0$ 处满足一般位置条件。缺陷由优势权 $\check\lambda\in X^*(\check T)^+$ 度量，等价地，由 $G$ 的优势余权度量。在相应的正则分层上，可以写为
\[\eta\bmod\check{\mathfrak b}[[t]]=\sum_{i\in I}t^{\langle\check\lambda,\check\alpha_i\rangle}\psi_i(t)f_i\,dt,\]
其中 $\psi_i\in k[[t]]$ 且 $\psi_i(0)\ne0$，$\alpha_i$ 是 $\check{\mathfrak g}$ 的简单根，$f_i\in\check{\mathfrak g}_{-\alpha_i}$ 是简单负根向量。将此轨迹记为 $\Op^{\mathrm{reg},\check\lambda}_{\checkG}$。在集合意义下，
\[\Op^{\mathrm{unr}}_{\checkG}(k)=\coprod_{\check\lambda\in X^*(\check T)^+}\Op^{\mathrm{reg},\check\lambda}_{\checkG}(k).\]
更准确地，$\Op^{\mathrm{unr}}_{\checkG}$ 具有分支或形式片 $\Op^{\mathrm{unr},\check\lambda}_{\checkG}$，且
\[\Op^{\mathrm{reg},\check\lambda}_{\checkG}\subset(\Op^{\mathrm{unr},\check\lambda}_{\checkG})_{\mathrm{red}}.\]
\srcpage{73}
\begin{theorem}{临界 FLE，点态形式\cite{FG07,ABC24a}}{15.2.1}
存在 DG 范畴的等价
\[\operatorname{FLE}_{G,\mathrm{crit}}:\widehat{\mathfrak g}_{\mathrm{crit}}\text{-}\mathrm{Mod}^{L^+G}\xrightarrow{\sim}\IndCoh(\Op^{\mathrm{unr}}_{\checkG}),\]
它对标准的 Frenkel–Gaitsgory 阿贝尔 $t$-结构是 $t$-正合的，这里省略惯常的 $*$ 型或 $!$ 型 $\IndCoh$ 选择。此等价与球面范畴的几何 Satake 作用相容。
\end{theorem}

这一相容性在 Weyl 对象上可以看出。令 $\mathbb W^{\check\lambda}_{\mathrm{crit}}$ 为相应最高权的临界 Weyl 对象，并令 $i_{\check\lambda}:\Op^{\mathrm{reg},\check\lambda}_{\checkG}\hookrightarrow\Op^{\mathrm{unr}}_{\checkG}$。Drinfeld–Sokolov 约化给出
\[\Psi(\mathbb W^{\check\lambda}_{\mathrm{crit}})\simeq\mathfrak z^{\mathrm{reg},\check\lambda}\simeq\End(\mathbb W^{\check\lambda}_{\mathrm{crit}}),\qquad
\Spec\mathfrak z^{\mathrm{reg},\check\lambda}\simeq\Op^{\mathrm{reg},\check\lambda}_{\checkG}.\]
在 FLE 下，这说明，在相应 $\IndCoh$ 约定下，$\mathbb W^{\check\lambda}_{\mathrm{crit}}$ 对应于 $(i_{\check\lambda})_*^{\IndCoh}\cO_{\Op^{\mathrm{reg},\check\lambda}_{\checkG}}$。因此，球面临界范畴的支撑恰好是非分歧 oper 轨迹；自守侧的非分歧对应于谱侧单值性平凡。

\section{Drinfeld–Sokolov 约化与 Whittaker 形式}\label{sec:15.3}\index{Drinfeld–Sokolov 约化|hyperpage}\index{半无限上同调|hyperpage}\index{BRST 约化|hyperpage}
FLE 的底层函子是 Drinfeld–Sokolov 约化\cite{FG05,Fre02}。选定标准的非退化加性特征标 $\chi:\mathfrak n((t))\to k$。对临界模 $M$，置
\[\Psi(M):=C^{\frac\infty2}(\mathfrak n((t)),M\otimes k_\chi).\]

\begin{transnote}{含糊点拨}
$\infty/2$ 指半无限维数：先取 $\mathfrak n((t))$ 在临界中心扩张下的真空模，再作半无限 Chevalley 上链，$k_\chi$ 是加性特征标给出的一维模。临界水平时 $\Psi$ 恰好等同于 Drinfeld–Sokolov 约化，因此 $\Psi(M)$ 可以当作 $M$ 的半无限不变量；缺掉这层，$\Psi$ 就只是一个记号。[来源：nLab "Drinfeld-Sokolov reduction"，https://ncatlab.org/nlab/show/Drinfeld-Sokolov+reduction（Last revised 2013-06-10；2026-11-17 访问）；并见原书 p.73 本式与 [FG05]。]
\end{transnote}
这就是 $M$ 的半无限上同调，也称 BRST 约化。在临界水平，它有一个增强版本，取值于穿孔圆盘上 oper 空间上的层。相应的 Whittaker 陈述为
\[\widehat{\mathfrak g}_{\mathrm{crit}}\text{-}\mathrm{Mod}^{\mathrm{Whit}}\simeq\IndCoh(\Op_{\checkG}(D^\times)).\]
从 Whittaker 形式转到球面／Kazhdan–Lusztig 子范畴，在谱侧相当于施加底层局部系统可延拓到 $D$ 的条件，因此将 $\Op(D^\times)$ 替换为 $\Op^{\mathrm{unr}}$。

Satake 作用有平行的谱描述。临界球面 Hecke 范畴被识别为局部系统 Steinberg 叠上的 ind-凝聚层：
\[\begin{split}
\DMod(L^+G\backslash LG/L^+G)_{\mathrm{crit}}\simeq{}\\
\IndCoh\bigl(\LocSys_{\checkG}(D)\times_{\LocSys_{\checkG}(D^\times)}\LocSys_{\checkG}(D)\bigr).
\end{split}\]
\enlargethispage{\baselineskip}
在 FLE 下，这个范畴的卷积变成 $\IndCoh(\Op^{\mathrm{unr}}_{\checkG})$ 上自然的对应作用。在 Frenkel–Gaitsgory 最初的证明中，这出现在阿贝尔范畴和阿贝尔几何 Satake 的层面；整体 Langlands 需要 DG／因子化表述。

\section{与量子局部 Langlands 的关系}\label{sec:15.4}\index{量子局部 Langlands|hyperpage}
临界 FLE 是量子局部几何 Langlands 的一个边界值\cite{FG05}。离开边界后，对偶水平满足关系
\[\kappa=\kappa_{G,\mathrm{crit}}+c\,\kappa_{\mathrm{Kil},G},\qquad
\check\kappa=\kappa_{\checkG,\mathrm{crit}}+c^{-1}\kappa_{\mathrm{Kil},\checkG}.\]

\begin{transnote}{含糊点拨}
对偶侧取 $c^{-1}$ 是量子 Langlands 的水平配方：$\kappa$ 与 $\check\kappa$ 配对成临界值。$c\to0$ 时 $\check G$-Whittaker 侧退化，大水平极限给出 Satake 对应，所以「量子局部 Langlands 的极限等价就是 FLE」；这里 $c$ 不是自由参数，而是两侧水平之和偏离临界值的量。[依据：原书 p.73 本式与随后一段；并见 [FG05]。]
\end{transnote}
当 $c\ne0$ 时，预期的局部等价具有形式
\[\widehat{\mathfrak g}_{\kappa}\text{-}\mathrm{Mod}^{L^+G}\simeq\Whit(D_{\check\kappa}(\Gr_{\checkG})).\]
它在边界上的对应形式是标准的 Satake 与 Casselman–Shalika 等价。在大水平极限中，球面与 Whittaker Grassmann 簇范畴退化为
\[D(\Gr_G)^{L^+G}_{\mathrm{sph}}\simeq\Rep(\checkG),\qquad\Whit(D(\Gr_{\checkG}))\simeq\Rep(G).\]
\srcpage{74}
当 $c\to0$ 时，$G$ 侧变为临界水平，而 $\checkG$-Whittaker 侧变为 oper 的经典谱几何；极限等价就是 FLE
\[\widehat{\mathfrak g}_{\mathrm{crit}}\text{-}\mathrm{Mod}^{L^+G}\simeq\IndCoh(\Op^{\mathrm{unr}}_{\checkG}).\]

这是范围更广的量子局部对应表中的一行：
\[\begin{aligned}
LG\text{-}\mathrm{Mod}_{\kappa}&\longleftrightarrow L\checkG\text{-}\mathrm{Mod}_{\check\kappa},\\
\widehat{\mathfrak g}_{\kappa}\text{-}\mathrm{Mod}&\longleftrightarrow\Whit(D_{\check\kappa}(L\checkG)),\\
\Whit(D_{\kappa}(LG))&\longleftrightarrow\widehat{\check{\mathfrak g}}_{\check\kappa}\text{-}\mathrm{Mod},\\
D_{\kappa}(\Gr_G)&\longleftrightarrow D_{\check\kappa}(\Gr_{\checkG}),\\
D_{\kappa}(\Fl_G)&\longleftrightarrow D_{\check\kappa}(\Fl_{\checkG}).
\end{aligned}\]
在 $\kappa=\kappa_{\mathrm{crit}}$ 时，预期整个局部 $LG$-模理论是 $\LocSys_{\checkG}(D^\times)$ 上的范畴层：
\[LG\text{-}\mathrm{Mod}_{\mathrm{crit}}\rightsquigarrow\operatorname{ShvCat}(\LocSys_{\checkG}(D^\times)).\]
因此，Kac–Moody 范畴是 oper 轨迹上的范畴：它的临界中心是 $\Op_{\checkG}(D^\times)$ 上的函数。经过 Whittaker／Drinfeld–Sokolov 约化后，这个 oper 上的范畴被识别为
\[\IndCoh(\Op_{\checkG}(D^\times)),\]
这与 Feigin–Frenkel 中心的识别 $\operatorname{Spf}Z_{\mathrm{crit}}\simeq\Op_{\checkG}(D^\times)$ 相容。

\section{因子化 FLE 及其与整体局部化的相容性}\label{sec:15.5}\index{基本局部等价!因子化形式|hyperpage}\index{谱 Poincaré 函子|hyperpage}
\thref{15.2.1}中的点态等价有 Ran 空间上的因子化版本\cite{ABC24a}：
\[\operatorname{FLE}_{G,\mathrm{crit},\mathrm{Ran}}:\operatorname{KL}_{G,\mathrm{crit},\mathrm{Ran}}\xrightarrow{\sim}\IndCoh(\Op^{\mathrm{unr}}_{\checkG})_{\mathrm{Ran}}.\]
它与到 $\Bun_G$ 的局部化函子及全局谱函子相容。基本图表为
\origdiagram{p074-square.pdf}{0.5419}
这里 $\mathbb L$ 表示整体 Langlands 函子。函子 $\operatorname{Poinc}^{\mathrm{spec}}$ 是沿如下对应的推前—拉回变换：这一对应比较全局局部系统与它在形式圆盘上的限制。对点 $x$，记
\[\Op_{\checkG}(X\setminus\{x\})^{\mathrm{unr}}:=\Op_{\checkG}(X\setminus\{x\})\times_{\LocSys_{\checkG}(X\setminus\{x\})}\LocSys_{\checkG}(X),\]
相应的谱对应为
\[\Op_{\checkG}(X\setminus\{x\})^{\mathrm{unr}}\longrightarrow\Op_{\checkG}(D_x^\times)\times_{\LocSys_{\checkG}(D_x^\times)}\LocSys_{\checkG}(D_x).\]
\srcpage{75}
后面证明 $\mathbb L$ 为等价时，将使用这个图表，以及系数函子和无穷水平 FLE：
\origdiagram{p075-square.pdf}{0.5464}

\section{多点情形与生成性}\label{sec:15.6}\index{局部化!多点情形|hyperpage}
对有限点集 $\mathbf x=\{x_1,\ldots,x_n\}$，因子化给出
\[\operatorname{KL}_{G,\mathrm{crit},\mathbf x}\simeq\bigotimes_{i=1}^n\widehat{\mathfrak g}_{\mathrm{crit},x_i}\text{-}\mathrm{Mod}^{L^+G_{x_i}}.\]
通常的局部化函子
\[\operatorname{Loc}_{\mathbf x}:\operatorname{KL}_{G,\mathrm{crit},\mathbf x}\longrightarrow\DMod_{\mathrm{crit}}(\Bun_G)\]
可以通过在穿孔曲线上加入全局谱参数来增强。更准确地，记 $\Op^{\mathrm{unr}}_{\checkG,\mathbf x}:=\prod_i\Op^{\mathrm{unr}}_{\checkG,x_i}$，增强函子示意地具有形式
\[\begin{split}
\operatorname{KL}_{G,\mathrm{crit},\mathbf x}\otimes_{\QCoh(\Op^{\mathrm{unr}}_{\checkG,\mathbf x})}
\QCoh(\Op_{\checkG}(X\setminus\mathbf x)^{\mathrm{unr}})\\
\longrightarrow\DMod_{\mathrm{crit}}(\Bun_G),
\end{split}\]
其中
\[\Op_{\checkG}(X\setminus\mathbf x)^{\mathrm{unr}}:=\Op_{\checkG}(X\setminus\mathbf x)\times_{\LocSys_{\checkG}(X\setminus\mathbf x)}\LocSys_{\checkG}(X).\]
等价地，$\operatorname{Rep}(\checkG)_{\mathbf x}$ 在 $\operatorname{Loc}_{\mathbf x}$ 的像上的作用经由 $\LocSys_{\checkG}(X)$ 分解。

当 $\mathbf x$ 足够大时，这些局部范畴生成尖点部分。一个有用的示意表述是：对每个拟紧开集 $U\subset\Bun_G$，函子
\[\operatorname{KL}_{G,\mathrm{crit},\mathrm{Ran}}\longrightarrow\DMod_{\mathrm{crit}}(\Bun_G)\longrightarrow\DMod_{\mathrm{crit}}(U)\]
的像，在限制到 $U$ 后生成尖点贡献。整体论证的其余部分是 Eisenstein 部分：通过归纳，以及 FLE 与 Eisenstein 级数和常数项函子的相容性，处理非尖点对象。

\chapter{整体猜想的证明}\label{ch:16}
% 原PDF第75页，第16章开始。
在\chref{15}的局部输入之后，本章按照现代抛物诱导输入\cite{CCF24}，证明整体范畴 Langlands 定理。我们省略通常的半密度扭转，仍用 $\DMod(\Bun_G)$ 表示上文的自守范畴。
\begin{theorem*}{}
Langlands 函子给出等价
\[\mathbb L_G:\DMod(\Bun_G)\xrightarrow{\sim}\IndCoh_{\mathrm{Nilp}}(\LocSys_{\checkG}),\]
并且与导出 Satake、Eisenstein 级数、常数项函子和对偶相容。
\end{theorem*}
证明先处理抛物部分，即 Eisenstein 部分，再证明剩余的尖点等价。

\section{粗 Langlands 函子}\label{sec:16.1}\index{Langlands 函子!粗形式|hyperpage}\index{Langlands 函子!重整化|hyperpage}
\srcpage{76}
第一个构造通过\cite{GR24a}中的 Langlands 函子构造，产生其粗形式
\[\mathbb L_G^0:\DMod(\Bun_G)\longrightarrow\QCoh(\LocSys_{\checkG}).\]
它由广义消失定理刻画：$\Rep(\checkG)_{\mathrm{Ran}}$ 在 $\DMod(\Bun_G)$ 上的局部 Hecke 作用，经由 $\QCoh(\LocSys_{\checkG})$ 的谱作用分解。采用第\ref{ch:13}–\ref{ch:15}章的记号，Whittaker 系数函子与谱作用给出图表
\origdiagram{p076-coarse.pdf}{0.4485}
这里 $\operatorname{CS}$ 表示几何 Casselman–Shalika 等价。完整的 Langlands 函子是 $\mathbb L_G^0$ 通过遗忘函子

\begin{transnote}{含糊点拨}
$\mathbb L_G$ 是粗函子 $\mathbb L_G^0$ 穿过遗忘函子 $\IndCoh_{\mathrm{Nilp}}\to\QCoh$ 的提升：目标必须换成重整化版本 $\QCoh^{!}$（ind-完备化），因为在紧对象上只能先落进下有界凝聚子范畴，普通 $\QCoh$ 不够，于是函数只在紧对象上与 $\mathbb L_G^0$ 相符。[依据：原书 p.76 本段；并见 [GR24a]。]
\end{transnote}
\[\Psi:\IndCoh_{\mathrm{Nilp}}(\LocSys_{\checkG})\longrightarrow\QCoh(\LocSys_{\checkG})\]
的提升，即 $\mathbb L_G^0=\Psi\circ\mathbb L_G$。在紧对象上，首先落入一个下有界的凝聚子范畴；可呈示函子 $\mathbb L_G$ 通过 ind-完备化得到。方便起见，用 $\QCoh^!(\LocSys_{\checkG})$ 表示相应的重整化谱目标。因此，简记为
\[\DMod(\Bun_G)\xrightarrow{\mathbb L_G}\QCoh^!(\LocSys_{\checkG}).\]
与导出 Satake 的相容性内置于这一构造中。新的工作是与 Eisenstein 和常数项函子的相容性\cite{CCF24}。

\section{Eisenstein 与常数项的相容性}\label{sec:16.2}\index{Eisenstein 函子!Langlands 相容性|hyperpage}\index{常数项函子!Langlands 相容性|hyperpage}
令 $P\subset G$ 为 Levi 商为 $M$ 的抛物子群，并令 $\check P\subset\checkG$ 为对偶抛物子群。忽略标准的移位与扭转，自守 Eisenstein 函子和常数项函子为
\[\operatorname{Eis}_P:\DMod(\Bun_M)\longrightarrow\DMod(\Bun_G),\qquad
\CT_P:\DMod(\Bun_G)\longrightarrow\DMod(\Bun_M).\]
在谱侧，它们由以下对应定义：
\[\LocSys_{\check M}\xleftarrow{q}\LocSys_{\check P}\xrightarrow{p}\LocSys_{\checkG}.\]
目标陈述是以下两个图表交换：
\origdiagramtag{p076-eis.pdf}{0.4645}{16.1}{eq:16.1}
以及
\origdiagramtag{p076-ct.pdf}{0.4713}{16.2}{eq:16.2}
\srcpage{77}
按\cite{CCF24}的形式，\eqref{eq:16.1}的证明归约到两个方块：一个涉及 Whittaker 系数的几何方块，以及一个涉及表示的谱方块。在几何 Casselman–Shalika 等价下，$\Whit(H)_x\simeq\Rep(\check H)_x$，所以这一比较变为 Whittaker 系数的抛物限制与表示上的相应 Chevalley 函子之间的相容性。

\section{谱方块：推前与局部化}\label{sec:16.3}\index{Chevalley 函子|hyperpage}\index{Ran 空间!标记点插入|hyperpage}
基本问题如下。对同态 $H_1\to H_2$，有映射 $\LocSys_{H_1}\to\LocSys_{H_2}$；对点 $x\in X$，还有局部的局部化函子
\[\operatorname{Loc}_x:\Rep(H_i)_x\longrightarrow\QCoh^!(\LocSys_{H_i}).\]
这里不能朴素地用一个方块，将沿 $\LocSys_{H_1}\to\LocSys_{H_2}$ 的拉回与表示限制比较。正确的方块使用推前和 Ran 插入。
\begin{lemma}{}{16.3.1}
假设 $H_1\to H_2$ 的核为幺幂群。令 $\operatorname{Ran}_{x,c}:=\{I\in\operatorname{Ran}(X)\mid x\in I\}$。则下面的图表交换：
\origdiagramtag{p077-general.pdf}{0.4713}{16.3}{eq:16.3}
\end{lemma}

这里，$\operatorname{inv}^+$ 先将对象插入标记点，再对幺幂核取导出不变量，等价地，取 Chevalley 上链。具体地，若 $V_x\in\Rep(H_1)_x$，则在点 $(x,y_1,\ldots,y_m)\in\operatorname{Ran}_{x,c}$，插入为

\begin{transnote}{含糊点拨}
先插入再取不变量，是为了与 Ran 空间的因子化结构相容：插入点 $x$ 上多出一个因子，再取幺幂核的导出不变量（即 Chevalley 上链 $C^{\bullet}(\check{\mathfrak u},-)$），剩下的正好是 Whittaker 归一化。每个插入点各带一个这样的因子，所以后文的公式是逐点乘起来的，而不是全局取一次。[依据：原书 p.77 本段；并见 [CCF24]。]
\end{transnote}
\[\operatorname{ins}(V_x)_{(x,y_1,\ldots,y_m)}=V_x\otimes\mathbf1_{y_1}\otimes\cdots\otimes\mathbf1_{y_m}.\]
\begin{proof}
证明是以下初等原理的 Ran 空间形式。若 $\pi:Y\to X$ 是一个 $D$-概形，且 $Y$ 拟仿射、余仿射，或满足相同的下降条件，则
\[\QCoh(Y)\simeq\pi_*\cO_Y\text{-}\mathrm{Mod}(\QCoh(X)),\]
其中 $\pi_*\cO_Y$ 一般是一个复形。例如，$\mathbb A^2\setminus\{0\}\to\mathrm{pt}$ 和 $B\Ga\to\mathrm{pt}$ 说明，必须保留这个复形，而不能只保留它的零次函数。将这一原理应用于形式圆盘，并使用插入与 Ran 卷积的相容性，就得到\eqref{eq:16.3}。
\end{proof}

对抛物商 $\check P\to\check M$，它的核是幺幂根基 $\check U$，引理给出
\origdiagram{p077-parabolic.pdf}{0.4690}
\srcpage{78}
再与 $\LocSys_{\check P}\to\LocSys_{\checkG}$ 复合，便得到常数项方向的谱方块
\origdiagram{p078-ct.pdf}{0.4690}
这里 $\operatorname{Chev}^+$ 是 Chevalley 函子。若 $V_x\in\Rep(\checkG)_x$，则示意地，
\[\operatorname{Chev}^+(V_x)_{(x,y_1,\ldots,y_m)}=
C^\bullet(\check{\mathfrak u},V_x)_x\otimes C^\bullet(\check{\mathfrak u})_{y_1}\otimes\cdots\otimes C^\bullet(\check{\mathfrak u})_{y_m},\]
其中 $\check{\mathfrak u}=\operatorname{Lie}(\check U)$ 是 $\check P$ 的幺幂根基的 Lie 代数。取对偶之后，等价地，在表示侧转到左伴随之后，出现相反的幺幂根基，给出 Eisenstein 方向的方块
\origdiagram{p078-eis.pdf}{0.4781}

\section{几何方块与半无限范畴}\label{sec:16.4}\index{半无限范畴|hyperpage}\index{第二伴随性|hyperpage}
几何侧是下面的方块：
\origdiagram{p078-geometric.pdf}{0.4485}
这里，$J^-:\Whit(M)_{\mathrm{Ran}_{x,c}}\to\Whit(G)_x$ 是 Eisenstein 方块中诱导方向的函子。同一个半无限核先给出局部限制函子 $r_P^-:\Whit(G)\to\Whit(M)$；正如第二伴随性所预言的，$J^-$ 是它的伴随形式。
\Needspace{4\baselineskip}
在局部，识别
\[\Whit(G)=\DMod(\Gr_G)^{LN,\chi_G},\qquad
\Whit(M)=\DMod(\Gr_M)^{LN_M,\chi_M}.\]
核从负半无限范畴构造：
\[\begin{aligned}
\operatorname{SI}_P^-&:=\DMod(\Gr_G)^{LN^-\cdot L^+M}\\
&\simeq(\DMod(\Gr_G)\otimes\DMod(\Gr_M))^{LP^-}\\
&\simeq\operatorname{Fun}^{LP^-}(\DMod(\Gr_G),\DMod(\Gr_M)).
\end{aligned}\]
因此，任意 $LP^-$-线性函子 $F:\DMod(\Gr_G)\to\DMod(\Gr_M)$ 都给出一个限制型函子
\[r_F:\DMod(\Gr_G)^{LN,\chi_G}\longrightarrow\DMod(\Gr_M)^{LN_M,\chi_M}.\]
上面使用的函子 $J^-$，是与 $\operatorname{SI}_P^-$ 的单位对象对应的伴随，即与轨道
\[LN^-\cdot1\cdot L^+G/L^+G\subset\Gr_G\]
上的标准对象对应的伴随。
\srcpage{79}
然后，将 Whittaker 特征标 $\chi_G$ 退化为一个在 $P$ 的幺幂根基上平凡、并在 $LN_M$ 上诱导 $\chi_M$ 的特征标 $\chi_P$。这给出局部模式
\[\DMod(\Gr_G)^{LN,\chi_G}\xrightarrow{\text{邻近圈}}\DMod(\Gr_G)^{LN,\chi_P}\longrightarrow\DMod(\Gr_M)^{LN_M,\chi_M}.\]
半无限对应使用 $P^-$，而特征标退化使用 $P$，这是第二伴随性的一个表现。

令 $\operatorname{Chev}^-:=(\operatorname{Chev}^+)^L$。在 Casselman–Shalika 等价下，所需的顶面方块如下：
\origdiagram{p079-chevalley.pdf}{0.4440}
证明使用 $\operatorname{SI}_P^-$ 的谱描述：
\begin{equation}
\begin{split}
\operatorname{SI}_P^-\simeq\QCoh^!\bigl(
\LocSys_{\checkG}(D)\times_{\LocSys_{\checkG}(D^\circ)}\LocSys_{\check P}(D^\circ)\\
\times_{\LocSys_{\check M}(D^\circ)}\LocSys_{\check M}(D)\bigr).
\end{split}
\tag{16.4}\label{eq:16.4}
\end{equation}

\begin{transnote}{含糊点拨}
四个因子依次读：$D$ 上的整体局部系统、它到 $D^\circ$ 的限制、$P$ 的抛物约化、$M$ 侧的局部系统。$P=G$ 时 $\LocSys_{\check P}=\LocSys_{\check G}$，右端退化为导出 Satake 的谱侧——这正是检验半无限范畴与谱侧对应是否正确的边界情形。[依据：原书 p.79 本式；并见 [CCF24]。]
\end{transnote}
当 $P=G$ 时，这恰好是导出 Satake 的谱侧。更一般地，它是导出 Satake 的抛物半无限对应版本。

局部 Langlands 的直观思想是：一个 $LG$-模范畴应当对应于 $\LocSys_{\checkG}(D^\circ)$ 上的范畴层。将环路群从 $G$ 限制到 $M$，对应于沿 $\LocSys_{\check M}(D^\circ)\to\LocSys_{\checkG}(D^\circ)$ 拉回。因此，预期有
\[\begin{aligned}
\DMod(\Gr_G)&\rightsquigarrow\QCoh(\LocSys_{\checkG}(D)),\\
\DMod(\Gr_M)&\rightsquigarrow\QCoh(\LocSys_{\check M}(D)),\\
\DMod(\Gr_G)^{LM}&\rightsquigarrow\QCoh\bigl(\LocSys_{\checkG}(D)\times_{\LocSys_{\checkG}(D^\circ)}\LocSys_{\check M}(D^\circ)\bigr).
\end{aligned}\]
在两侧计算 $\operatorname{Fun}(\text{--},\text{--})$，就得到\eqref{eq:16.4}，从而证明 Eisenstein 相容性\eqref{eq:16.1}。

\section{为何常数项更困难}\label{sec:16.5}\index{Kac–Moody–Whittaker 对偶|hyperpage}
证明\eqref{eq:16.1}的方法不能直接证明常数项相容性\eqref{eq:16.2}。在 Eisenstein 图表中，所有箭头都指向推前所需的方向。对于 $\CT_P$，必须反转这些箭头，而简单的几何论证不再适用。

替代方法是 Kac–Moody–Whittaker 对偶。在非临界水平，比较
\origdiagram{p079-noncritical.pdf}{0.4599}
特殊化到临界水平，等价地，特殊化到 FLE 退化中量子参数的边界值后，局部 Whittaker 范畴被球面临界 Kac–Moody 范畴替换。在 $\IndCoh$ 约定下，FLE 将它们识别为非分歧 oper：
\[\operatorname{KL}_{G,\mathrm{crit}}\simeq\IndCoh(\Op^{\mathrm{unr}}_{\checkG}),\qquad
\operatorname{KL}_{M,\mathrm{crit}}\simeq\IndCoh(\Op^{\mathrm{unr}}_{\check M}).\]
相应的局部函子为
\[r_P^{\mathrm{KM}}:\operatorname{KL}_{G,\mathrm{crit}}\longrightarrow\operatorname{KL}_{M,\mathrm{crit}},\qquad
M\longmapsto C^{\infty/2}(\mathfrak u((t)),M),\]
\srcpage{80}
其中 $\mathfrak u$ 是 $P$ 的幺幂根基的 Lie 代数。它的谱对应物是函子
\[r_{\check P}^{\mathrm{op}}:\IndCoh(\Op^{\mathrm{unr}}_{\checkG})\longrightarrow\IndCoh(\Op^{\mathrm{unr}}_{\check M}),\]
由 $\operatorname{SI}_{\check P^-}^{\mathrm{spec}}$ 中的单位对象给出。这些函子证明了常数项立方体的顶面与侧面\cite{CCF24}。剩余的前面与后面并非仅由半无限论证推出，而是通过手征方法证明。这就完成了常数项相容性\eqref{eq:16.2}的证明。

\section{伴随与 Eisenstein 生成的部分}\label{sec:16.6}\index{Eisenstein 生成子范畴|hyperpage}\index{可约局部系统轨迹|hyperpage}\index{不可约局部系统轨迹|hyperpage}
Eisenstein 论证的最后部分解释，相容性\eqref{eq:16.1}和\eqref{eq:16.2}如何进入整体归纳。令 $\DMod(\Bun_G)_{\mathrm{Eis}}$ 为由真抛物子群 $P$ 对应的 $\operatorname{Eis}_P(\DMod(\Bun_M))$ 生成的子范畴，令 $\LocSys^{\mathrm{red}}_{\checkG}\subset\LocSys_{\checkG}$ 为可约轨迹，即所有真抛物子群 $\check P$ 对应的 $\LocSys_{\check P}\to\LocSys_{\checkG}$ 的像的并。等价地，这些局部系统容许真抛物约化。其补集是不可约轨迹 $\LocSys^{\mathrm{irr}}_{\checkG}$。

需要证明粗函子 $\mathbb L_G^0$ 有左伴随。要点是，只需在生成对象上构造左伴随。在谱侧，生成对象来自 $\operatorname{Eis}^{\mathrm{spec}}_{\check P}(\QCoh(\LocSys_{\check M}))$，而 Eisenstein 相容性通过从 $M$ 输运伴随，给出所需的自守函子。等价地，使用以下事实：转到 Ran 作用描述后，粗函子已经有左伴随；$\Rep(\check M)_{\mathrm{Ran}}$ 作用于自守侧，而商 $\Rep(\check M)_{\mathrm{Ran}}\to\QCoh(\LocSys_{\check M})$ 识别谱作用。

因此，假设每个真 Levi 子群 $M$ 的 Langlands 等价已经成立，则函子 $\mathbb L_G$ 在 Eisenstein 直和因子上诱导等价\cite{CCF24}：
\begin{equation}
\mathbb L_G:\DMod(\Bun_G)_{\mathrm{Eis}}\xrightarrow{\sim}\QCoh^!(\LocSys_{\checkG})_{\mathrm{red}}.
\tag{16.5}\label{eq:16.5}
\end{equation}
于是，整体猜想被归约到尖点部分。剩余的陈述是
\[\mathbb L_{G,\mathrm{cusp}}:\DMod(\Bun_G)_{\mathrm{cusp}}\xrightarrow{\sim}\QCoh(\LocSys^{\mathrm{irr}}_{\checkG}).\]
这个归约为下面的尖点论证提供输入。

现在，在上述 Eisenstein 归约之后，按照证明系列中的重数一定理部分\cite{GR24b}，完成证明。仍用记号
\[\mathbb L_G:\DMod(\Bun_G)\longrightarrow\QCoh^!(\LocSys_{\checkG})\]
表示重整化 Langlands 函子，并用 $\mathbb L_G^L$ 和 $\mathbb L_G^R$ 表示它在相关直和因子上限制后的左伴随和右伴随。前面各节构造了 $\mathbb L_G$、构造了 $\mathbb L_G^L$，并证明 Eisenstein 等价
\[\mathbb L_G:\DMod(\Bun_G)_{\mathrm{Eis}}\xrightarrow{\sim}\QCoh^!(\LocSys_{\checkG})_{\mathrm{red}}.\]
剩下的任务是证明尖点等价
\[\mathbb L_{G,\mathrm{cusp}}:\DMod(\Bun_G)_{\mathrm{cusp}}\xrightarrow{\sim}\QCoh(\LocSys^{\mathrm{irr}}_{\checkG}).\]
这里，$\LocSys^{\mathrm{irr}}_{\checkG}$ 是不可约局部系统的开子叠，即不容许约化到真抛物子群 $\check P\subsetneq\checkG$ 的局部系统。

\section{尖点函子的已知输入}\label{sec:16.7}\index{保守函子!尖点 Langlands 函子|hyperpage}\index{单子|hyperpage}
\srcpage{81}
第一个输入是保守性。
\begin{theorem}{}{16.7.1}
函子 $\mathbb L_{G,\mathrm{cusp}}$ 是保守的。
\end{theorem}
\begin{proof}
证明思路大致使用 Whittaker 系数与几何 Ramanujan 输入\cite{FR25,Ber21,GR24b}。自守系数函子给出如下形式的方块：
\origdiagram{p081-cuspidal.pdf}{0.4918}
下方的水平函子是 Casselman–Shalika 识别。由几何 Ramanujan 输入，尖点对象是温和的，而 Whittaker 系数函子在温和子范畴上保守。因此，$\mathbb L_{G,\mathrm{cusp}}$ 是保守的。
\end{proof}

第二个输入是\secref{16.6}中已经构造的伴随的存在性。因此，要证明 $\mathbb L_{G,\mathrm{cusp}}$ 是等价，只需识别单子
\[\mathbb L_{G,\mathrm{cusp}}\circ\mathbb L_{G,\mathrm{cusp}}^L\in\End(\QCoh(\LocSys^{\mathrm{irr}}_{\checkG})).\]
谱范畴是 $\QCoh(\LocSys^{\mathrm{irr}}_{\checkG})$-线性的。因此，这个单子是与结合代数对象
\[\begin{split}
\mathcal A_{G,\mathrm{irr}}&:=(\mathbb L_{G,\mathrm{cusp}}\circ\mathbb L_{G,\mathrm{cusp}}^L)(\cO_{\LocSys^{\mathrm{irr}}_{\checkG}})\\
&\in\operatorname{AssocAlg}(\QCoh(\LocSys^{\mathrm{irr}}_{\checkG}))
\end{split}\]
作张量积。等价地，若在整个谱侧取 $\mathcal A_G=\mathbb L_G\mathbb L_G^L(\cO)$，则 $\mathcal A_{G,\mathrm{irr}}\simeq\mathcal A_G|_{\LocSys^{\mathrm{irr}}_{\checkG}}$。这一限制并非仅由 $\QCoh^!$ 的记号形式地推出，而是来自重整化谱范畴与不可约开集上的普通 $\QCoh$-模结构的相容性。因此，唯一剩下的任务是
\[\cO_{\LocSys^{\mathrm{irr}}_{\checkG}}\xrightarrow{\sim}\mathcal A_{G,\mathrm{irr}}.\]
这个陈述的可约部分，已在上文通过 Eisenstein 归纳证明。

\section{双伴随性与自对偶代数}\label{sec:16.8}\index{双伴随性|hyperpage}\index{自对偶代数|hyperpage}
\begin{theorem}{双伴随性\cite{ABC24b}}{16.8.1}
存在伴随的典范等价
\[\mathbb L_{G,\mathrm{cusp}}^L\simeq\mathbb L_{G,\mathrm{cusp}}^R.\]
\end{theorem}
\begin{proof}
证明概述如下。通过上面使用的局部化与系数图表比较两侧。自守尖点范畴和 $\QCoh(\LocSys^{\mathrm{irr}}_{\checkG})$ 都带有自然对偶。$\mathbb L_G$ 与 FLE、谱 Poincaré 函子和 Whittaker 系数的相容性，将 $\mathbb L_{G,\mathrm{cusp}}$ 的对偶函子识别为同一个函子。因此，左伴随和右伴随相同。
\end{proof}

其结果是：代数 $\mathcal A_{G,\mathrm{irr}}$ 作为 $\LocSys^{\mathrm{irr}}_{\checkG}$ 上的复形是完美的，并且典范自对偶。事实上，右伴随给出与左伴随相同的核，但采用相反的伴随约定。因此，与 $\mathcal A_{G,\mathrm{irr}}$ 对偶的对象是同一个对象，只是现在将它看作余代数。将这一余代数记为 $\mathcal B_{G,\mathrm{irr}}$。用公式表示，
\[\mathcal B_{G,\mathrm{irr}}:=\mathcal A_{G,\mathrm{irr}}^\vee,\qquad
\mathcal A_{G,\mathrm{irr}}^\vee\simeq\mathcal A_{G,\mathrm{irr}}.\]
\srcpage{82}
正是在这一步，双伴随性将单子问题转化为关于谱侧伪真 Ran 推前的几何陈述。

\section{一般位置 oper 与余代数}\label{sec:16.9}\index{oper!扩展一般位置结构|hyperpage}\index{余代数|hyperpage}\index{有限 étale 代数|hyperpage}
令
\[\pi:\Op^{\mathrm{gen}}_{\checkG,\mathrm{irr}}\longrightarrow\LocSys^{\mathrm{irr}}_{\checkG}\]
为不可约局部系统上的一般位置，更准确地说是扩展一般位置 $\checkG$-oper 结构的叠。非正式地说，在 $\sigma\in\LocSys^{\mathrm{irr}}_{\checkG}$ 上的纤维，是 $\sigma$ 到 $\check B$ 的一般位置约化所组成的空间，这些约化在 $X$ 的有限多个点之外满足 oper 横截条件。FLE／局部化形式体系将上述余代数识别为 de Rham 直接像\cite{ABC24a,GR24b}
\[\mathcal B_{G,\mathrm{irr}}\simeq\pi_{\mathrm{dR},!}\omega_{\Op^{\mathrm{gen}}_{\checkG,\mathrm{irr}}},\]
这是 $\QCoh(\LocSys^{\mathrm{irr}}_{\checkG})$ 中的余交换余代数对象。在证明中，此对象由真空 Poincaré 核构造：局部 FLE 说明，Kac–Moody 真空对象被送到非分歧 oper 对象，而 Ran 因子化将局部的一般位置约化粘合为上述全局一般位置 oper 对应。

下面使用几个初等推论。第一，$\mathcal B_{G,\mathrm{irr}}$ 余交换，位于上同调次数 $\le0$，并沿 $\LocSys^{\mathrm{irr}}_{\checkG}$ 带有平坦联络。第二，由自对偶性，$\mathcal A_{G,\mathrm{irr}}$ 交换，位于次数 $\ge0$。由于对偶将这一完美对象送回自身，反向的次数界也成立；所以 $\mathcal A_{G,\mathrm{irr}}$ 位于心范畴中，并且是带有平坦联络的向量丛。第三，$\pi_{\mathrm{dR},!}\omega$ 具有有限单值性。因此：
\begin{theorem}{}{16.9.1}
代数 $\mathcal A_{G,\mathrm{irr}}$ 是有限 étale 交换 $\cO_{\LocSys^{\mathrm{irr}}_{\checkG}}$-代数。等价地，它定义 $\LocSys^{\mathrm{irr}}_{\checkG}$ 的有限 étale 覆盖。
\end{theorem}
\begin{note}\numberednote{16.9.2}
对经典群，Beraldo–Kazhdan–Schlank 的定理加强了这个有限单值性输入：相关的一般位置 oper 纤维是同调可缩的\cite{BKS18}。因此，在 $A$、$B$、$C$、$D$ 型中，该有限 étale 覆盖已被迫为平凡覆盖。证明以此为模型情形，再将一般情形归约到同样的有限 étale 论证。
\end{note}

\section{归约到一个数值秩}\label{sec:16.10}\index{重数一定理|hyperpage}\index{Cohen–Macaulay 性质|hyperpage}
现在限制到证明中使用的标准良好情形：$G$ 为伴随型，且 $g_X>1$；小的例外情形 $g_X=2$ 且 $\checkG=\SL_2$ 通过单独的归约处理。在这一情形中，$\LocSys^{\mathrm{irr}}_{\checkG}$ 在这里所需的意义下没有非平凡的有限 étale 覆盖。因此，它上面的任意有限 étale 代数都是常值的：
\[\mathcal A_{G,\mathrm{irr}}\simeq\cO_{\LocSys^{\mathrm{irr}}_{\checkG}}^{\oplus n},\]
其中 $n\ge1$ 为整数。要证明 $n=1$，只需证明全局截面上的映射
\[\Gamma(\LocSys^{\mathrm{irr}}_{\checkG},\cO)\longrightarrow\Gamma(\LocSys^{\mathrm{irr}}_{\checkG},\mathcal A_{G,\mathrm{irr}})\]
是同构。在良好情形中，不可约轨迹的补集余维至少为二，而且局部系统叠是 Cohen–Macaulay 的；因此，将 $\mathcal A_{G,\mathrm{irr}}$ 延拓为整个 $\LocSys_{\checkG}$ 上的全局代数 $\mathcal A_G$ 后，可以检验同一个问题。只需证明
\[\Gamma(\LocSys_{\checkG},\cO)\xrightarrow{\sim}\Gamma(\LocSys_{\checkG},\mathcal A_G).\]
\srcpage{83}
全局函数方面的输入给出 $\Gamma(\LocSys_{\checkG},\cO)=k$。由全局 Poincaré 真空核的定义，右端是基本 Whittaker 对象 $W_i$ 的自同态代数：
\[\Gamma(\LocSys_{\checkG},\mathcal A_G)\simeq\End(W_i).\]
因此，最后需要计算的是 $\End(W_i)=k$。

计算中所用的 Whittaker 图卡，是通常的 Drinfeld 紧化中的开分层。对充分负的次数数据 $i$，有一个局部闭图：
\origdiagram{p083-whittaker.pdf}{0.4850}
其中 $r(G)=\operatorname{rank}_{\mathrm{ss}}(G)$，$\chi_{\mathrm{glob}}$ 是简单根 Whittaker 特征标的乘积。相关的极小层形如
\[\mathring{\mathcal A}_{\min}\simeq\Gm^{r(G)}\times BU_i,\]
其中 $U_i$ 为幺幂群。对于洁净 Whittaker 对象，$BU_i$ 因子只贡献单位对象；因此，相关的自同态计算归约为
\[\End_{\DMod(\mathring{\mathcal A}_{\min})}\bigl(W_i|_{\mathring{\mathcal A}_{\min}}\bigr)
\simeq\End_{\DMod(\Gm^{r(G)})}\bigl(\mathcal E^{\boxtimes r(G)}\bigr).\]
对象 $W_i$ 在这一图卡上限制为洁净指数对象 $\chi_{\mathrm{glob}}^!\mathcal E$。因此，这一计算就是 $\Gm^{r(G)}$ 上 $\mathcal E^{\boxtimes r(G)}$ 的初等计算，并且
\[\End(W_i)=k.\]
所以 $n=1$，从而 $\mathcal A_{G,\mathrm{irr}}\simeq\cO$。于是，单子 $\mathbb L_{G,\mathrm{cusp}}\mathbb L_{G,\mathrm{cusp}}^L$ 是恒等单子。由于 $\mathbb L_{G,\mathrm{cusp}}$ 是保守的，并且具有所需的伴随，它是一个等价：
\[\mathbb L_{G,\mathrm{cusp}}:\DMod(\Bun_G)_{\mathrm{cusp}}\xrightarrow{\sim}\QCoh(\LocSys^{\mathrm{irr}}_{\checkG}).\]
再与可约部分的 Eisenstein 等价\eqref{eq:16.5}结合，便证明了整体几何 Langlands 等价\cite{CCF24,ABC24b,GR24b}。

% CHAPTER_INPUTS
% 由 work/gen_index.py 生成，勿手改
\chapter*{译注一览}
\addcontentsline{toc}{chapter}{译注一览}
\markboth{译注一览}{译注一览}
\begin{multicols}{2}\small
\noindent\textbf{1}\;\textit{含糊点拨}\;几乎同构：idele 类群的单位分支\hfill p.\,\pageref{tn:1}\\[1pt]
\noindent\textbf{2}\;\textit{含糊点拨}\;两次修改的交换性条件即余圈条件\hfill p.\,\pageref{tn:2}\\[1pt]
\noindent\textbf{3}\;\textit{易错预警}\;上标「心」记号指 t-结构的心\hfill p.\,\pageref{tn:3}\\[1pt]
\noindent\textbf{4}\;\textit{背景直觉}\;幂零奇异支撑与 AG 修正\hfill p.\,\pageref{tn:4}\\[1pt]
\noindent\textbf{5}\;\textit{含糊点拨}\;「无限级」指形式圆盘上的平凡化\hfill p.\,\pageref{tn:5}\\[1pt]
\noindent\textbf{6}\;\textit{背景直觉}\;卷积即融合积，闭包上的 IC 层对应不可约表示\hfill p.\,\pageref{tn:6}\\[1pt]
\noindent\textbf{7}\;\textit{含糊点拨}\;相对位置与 Birkhoff 分解\hfill p.\,\pageref{tn:7}\\[1pt]
\noindent\textbf{8}\;\textit{易错预警}\;移位与 Tate 扭转两个因子\hfill p.\,\pageref{tn:8}\\[1pt]
\noindent\textbf{9}\;\textit{含糊点拨}\;半和正根在基本权上的值\hfill p.\,\pageref{tn:9}\\[1pt]
\noindent\textbf{10}\;\textit{含糊点拨}\;有理平凡化与缺陷除子\hfill p.\,\pageref{tn:10}\\[1pt]
\noindent\textbf{11}\;\textit{易错预警}\;特征标群是 Pontryagin 对偶\hfill p.\,\pageref{tn:11}\\[1pt]
\noindent\textbf{12}\;\textit{含糊点拨}\;mirabolic 子群与两条轨道\hfill p.\,\pageref{tn:12}\\[1pt]
\noindent\textbf{13}\;\textit{含糊点拨}\;沿 Hecke 对应的卷积\hfill p.\,\pageref{tn:13}\\[1pt]
\noindent\textbf{14}\;\textit{含糊点拨}\;Artin–Schreier 层与洁净判别\hfill p.\,\pageref{tn:14}\\[1pt]
\noindent\textbf{15}\;\textit{背景直觉}\;标量自同构胚叠与刚性化\hfill p.\,\pageref{tn:15}\\[1pt]
\noindent\textbf{16}\;\textit{含糊点拨}\;叠的维数为负的算法\hfill p.\,\pageref{tn:16}\\[1pt]
\noindent\textbf{17}\;\textit{含糊点拨}\;ρ 扭曲与紧化整体版本\hfill p.\,\pageref{tn:17}\\[1pt]
\noindent\textbf{18}\;\textit{含糊点拨}\;归一化与反常 t-正合\hfill p.\,\pageref{tn:18}\\[1pt]
\noindent\textbf{19}\;\textit{含糊点拨}\;Deligne–Abel–Jacobi 消失\hfill p.\,\pageref{tn:19}\\[1pt]
\noindent\textbf{20}\;\textit{含糊点拨}\;两个消失界的关系\hfill p.\,\pageref{tn:20}\\[1pt]
\noindent\textbf{21}\;\textit{含糊点拨}\;「相对」Fourier–Deligne 变换\hfill p.\,\pageref{tn:21}\\[1pt]
\noindent\textbf{22}\;\textit{易错预警}\;用上叹号拉回的收缩原理\hfill p.\,\pageref{tn:22}\\[1pt]
\noindent\textbf{23}\;\textit{含糊点拨}\;基础 Whittaker 层\hfill p.\,\pageref{tn:23}\\[1pt]
\noindent\textbf{24}\;\textit{含糊点拨}\;因子化恒等式与挠商\hfill p.\,\pageref{tn:24}\\[1pt]
\noindent\textbf{25}\;\textit{含糊点拨}\;因子化恒等式即层化陈述\hfill p.\,\pageref{tn:25}\\[1pt]
\noindent\textbf{26}\;\textit{含糊点拨}\;标量胚叠与对称群的下降\hfill p.\,\pageref{tn:26}\\[1pt]
\noindent\textbf{27}\;\textit{含糊点拨}\;中间延拓的判别\hfill p.\,\pageref{tn:27}\\[1pt]
\noindent\textbf{28}\;\textit{含糊点拨}\;幂零锥商叠与仿射直线上的挠层\hfill p.\,\pageref{tn:28}\\[1pt]
\noindent\textbf{29}\;\textit{含糊点拨}\;按分层计算 Euler 示性数\hfill p.\,\pageref{tn:29}\\[1pt]
\noindent\textbf{30}\;\textit{含糊点拨}\;FGKV 比较定理的两边\hfill p.\,\pageref{tn:30}\\[1pt]
\noindent\textbf{31}\;\textit{易错预警}\;相对维数与半扭转两个因子\hfill p.\,\pageref{tn:31}\\[1pt]
\noindent\textbf{32}\;\textit{含糊点拨}\;真抛物常数项与尖点性\hfill p.\,\pageref{tn:32}\\[1pt]
\noindent\textbf{33}\;\textit{含糊点拨}\;沿满射的余极限\hfill p.\,\pageref{tn:33}\\[1pt]
\noindent\textbf{34}\;\textit{含糊点拨}\;Cauchy 分解与 Schur 函子\hfill p.\,\pageref{tn:34}\\[1pt]
\noindent\textbf{35}\;\textit{含糊点拨}\;平均函子的双侧伴随\hfill p.\,\pageref{tn:35}\\[1pt]
\noindent\textbf{36}\;\textit{背景直觉}\;几何 Ramanujan 定理\hfill p.\,\pageref{tn:36}\\[1pt]
\noindent\textbf{37}\;\textit{含糊点拨}\;半扭转的定义\hfill p.\,\pageref{tn:37}\\[1pt]
\noindent\textbf{38}\;\textit{易错预警}\;奇异支撑叠不是奇点集\hfill p.\,\pageref{tn:38}\\[1pt]
\noindent\textbf{39}\;\textit{背景直觉}\;导出 Satake 与幂零奇异支撑\hfill p.\,\pageref{tn:39}\\[1pt]
\noindent\textbf{40}\;\textit{含糊点拨}\;微局部支撑的两种定义\hfill p.\,\pageref{tn:40}\\[1pt]
\noindent\textbf{41}\;\textit{含糊点拨}\;可呈示性：余极限等于极限\hfill p.\,\pageref{tn:41}\\[1pt]
\noindent\textbf{42}\;\textit{易错预警}\;左完备化不是 Ind-完备化\hfill p.\,\pageref{tn:42}\\[1pt]
\noindent\textbf{43}\;\textit{含糊点拨}\;扭转与 Hitchin 基\hfill p.\,\pageref{tn:43}\\[1pt]
\noindent\textbf{44}\;\textit{含糊点拨}\;Betti 局部到整体的一族函子\hfill p.\,\pageref{tn:44}\\[1pt]
\noindent\textbf{45}\;\textit{含糊点拨}\;Ran 范畴是余极限\hfill p.\,\pageref{tn:45}\\[1pt]
\noindent\textbf{46}\;\textit{易错预警}\;右 t-正合不能省\hfill p.\,\pageref{tn:46}\\[1pt]
\noindent\textbf{47}\;\textit{背景直觉}\;范畴迹与可对偶性\hfill p.\,\pageref{tn:47}\\[1pt]
\noindent\textbf{48}\;\textit{含糊点拨}\;幂零奇异支撑与温和部分\hfill p.\,\pageref{tn:48}\\[1pt]
\noindent\textbf{49}\;\textit{含糊点拨}\;重构映射与双 Whittaker 范畴\hfill p.\,\pageref{tn:49}\\[1pt]
\noindent\textbf{50}\;\textit{含糊点拨}\;非正则幂零锥与反温和\hfill p.\,\pageref{tn:50}\\[1pt]
\noindent\textbf{51}\;\textit{含糊点拨}\;oper 叠是 Hitchin 挠子\hfill p.\,\pageref{tn:51}\\[1pt]
\noindent\textbf{52}\;\textit{含糊点拨}\;Kostant 切片\hfill p.\,\pageref{tn:52}\\[1pt]
\noindent\textbf{53}\;\textit{含糊点拨}\;最低 L0-权\hfill p.\,\pageref{tn:53}\\[1pt]
\noindent\textbf{54}\;\textit{易错预警}\;FLE 谱空间为何更大\hfill p.\,\pageref{tn:54}\\[1pt]
\noindent\textbf{55}\;\textit{含糊点拨}\;半无限上同调 ∞/2\hfill p.\,\pageref{tn:55}\\[1pt]
\noindent\textbf{56}\;\textit{含糊点拨}\;对偶水平配方\hfill p.\,\pageref{tn:56}\\[1pt]
\noindent\textbf{57}\;\textit{含糊点拨}\;谱侧函子对粗函子的提升\hfill p.\,\pageref{tn:57}\\[1pt]
\noindent\textbf{58}\;\textit{含糊点拨}\;插入与 Chevalley 上链\hfill p.\,\pageref{tn:58}\\[1pt]
\noindent\textbf{59}\;\textit{含糊点拨}\;式 (16.4) 的四个因子\hfill p.\,\pageref{tn:59}\\[1pt]
\end{multicols}

\backmatter
\nocite{*}
\printbibliography[heading=bibintoc,title={References}]

\printindex
\end{document}
